%% ============================================================
%% Metodi Diagnostici Avanzati e di Ricerca:
%% Patologia Digitale e Analisi delle Immagini
%% Dispensa del Corso — Documento Principale
%% ============================================================
\documentclass[12pt, a4paper, twoside]{book}

%% --- Codifica e Lingua ---
\usepackage[utf8]{inputenc}
\usepackage[T1]{fontenc}
\usepackage[italian]{babel}

%% --- Colori (deve precedere tcolorbox e hyperref) ---
\usepackage[dvipsnames, table, x11names]{xcolor}
\definecolor{pathblue}{RGB}{30, 90, 160}
\definecolor{pathgray}{RGB}{245, 245, 245}
\definecolor{pathorange}{RGB}{220, 100, 30}

%% --- Tipografia ---
\usepackage{lmodern}
\usepackage{microtype}

%% --- Impaginazione ---
\usepackage[
  a4paper,
  top=2.5cm,
  bottom=2.5cm,
  inner=3cm,
  outer=2.5cm,
  headheight=14pt
]{geometry}

\usepackage{textcomp}

%% --- Matematica ---
\usepackage{amsmath}
\usepackage{amssymb}

%% --- Grafica e Figure ---
\usepackage{graphicx}
\usepackage{tikz}
\usetikzlibrary{arrows, shapes, shapes.geometric,
                positioning, fit, calc,
                decorations.pathreplacing, arrows.meta, 
                backgrounds, mindmap, trees}

%% --- Tabelle ---
\usepackage{booktabs}
\usepackage{array}
\usepackage{longtable}

%% --- Elenchi ---
\usepackage{enumitem}

%% --- Riquadri e Callout (dopo xcolor) ---
\usepackage{tcolorbox}
\tcbuselibrary{skins, breakable}

\newtcolorbox{obiettiviapprendimento}{
  colback=pathgray,
  colframe=pathblue,
  fonttitle=\bfseries,
  title=Obiettivi di Apprendimento,
  breakable
}

\newtcolorbox{terminechave}[1]{
  colback=blue!5,
  colframe=pathblue!60,
  fonttitle=\bfseries\small,
  title=Termine Chiave: #1,
  breakable
}

\newtcolorbox{notaclinica}{
  colback=orange!8,
  colframe=pathorange,
  fonttitle=\bfseries,
  title=Nota Clinica,
  breakable
}

%% --- Hyperlink (dopo xcolor e tcolorbox) ---
\usepackage[
  colorlinks=true,
  linkcolor=pathblue,
  citecolor=pathblue,
  urlcolor=pathorange,
  bookmarks=true,
  bookmarksnumbered=true,
  pdftitle={Metodi Diagnostici Avanzati e di Ricerca: Patologia Digitale e Analisi delle Immagini},
  pdfauthor={Dispensa del Corso},
  pdfsubject={Patologia Digitale},
  pdfkeywords={patologia digitale, whole slide imaging, intelligenza artificiale, apprendimento profondo}
]{hyperref}

%% --- Bibliografia ---
\usepackage[round, authoryear, sort&compress]{natbib}
\bibliographystyle{plainnat}

%% --- Didascalie ---
\usepackage[font=small, labelfont=bf, margin=10pt]{caption}

%% --- Intestazioni e Piè di Pagina ---
\usepackage{fancyhdr}
\pagestyle{fancy}
\fancyhf{}
\fancyhead[LE]{\small\leftmark}
\fancyhead[RO]{\small\rightmark}
\fancyfoot[C]{\thepage}
\renewcommand{\headrulewidth}{0.4pt}

%% --- Stile dei Capitoli ---
\usepackage{titlesec}
\titleformat{\chapter}[display]
  {\normalfont\huge\bfseries\color{pathblue}}
  {\chaptertitlename\ \thechapter}{20pt}{\Huge}
\titlespacing*{\chapter}{0pt}{30pt}{20pt}

\titleformat{\section}
  {\normalfont\Large\bfseries\color{pathblue!80!black}}
  {\thesection}{1em}{}
\titleformat{\subsection}
  {\normalfont\large\bfseries}
  {\thesubsection}{1em}{}

%% --- Varie ---
\usepackage{csquotes}
\usepackage{setspace}
\onehalfspacing

%% ============================================================
\begin{document}

%% ---- FRONTESPIZIO ----
\begin{titlepage}
  \centering
  \vspace*{2cm}

  {\color{pathblue}\rule{\linewidth}{1.5pt}}
  \vspace{0.5cm}

  {\Huge\bfseries\color{pathblue}
    Patologia Digitale\\[0.3em]
    e Analisi delle Immagini}

  \vspace{0.4cm}
  {\LARGE\bfseries\color{pathblue!70!black}
    Diagnostica avanzata e metodologie della ricerca}

  \vspace{0.3cm}
  {\color{pathblue}\rule{\linewidth}{1.5pt}}

  \vspace{1.5cm}

  {\Large\itshape Dispensa del Corso}

  \vspace{0.5cm}
  {\large a cura di Salvatore Lorenzo Renne*}

  \vfill

  \begin{tcolorbox}[colback=pathgray, colframe=pathblue, width=0.85\linewidth]
    \centering
    \textbf{Panoramica del Corso}\\[0.5em]
    \small
    Questa dispensa accompagna il corso di cinque lezioni sulla patologia digitale e
    l'analisi delle immagini. Tratta la transizione dalla patologia tradizionale a quella
    digitale, i flussi di lavoro per il controllo qualità, i fondamenti dell'intelligenza
    artificiale, le applicazioni dell'apprendimento profondo e multimodale, e le sfide
    etiche della diagnosi assistita dall'IA. Ogni capitolo include obiettivi di
    apprendimento, contenuti teorici, attività pratiche e domande di autovalutazione.
  \end{tcolorbox}

  \vspace{1.5cm}

  {\large Anno Accademico 2025--2026}

  \vspace{0.5cm}
  {\small\color{gray}
    *\textit{Professore Associato di Anatomia Patologica, Humanitas University, Milano. \\
    Email: \href{mailto:salvatore.renne@hunimed.eu}{salvatore.renne@hunimed.eu}}}

\end{titlepage}

%% ---- MATERIALE PRELIMINARE ----
\frontmatter

%% Prefazione
\chapter*{Prefazione}
\addcontentsline{toc}{chapter}{Prefazione}

La patologia digitale sta trasformando la pratica dell'istopatologia. La digitalizzazione dei
vetrini istologici in immagini di vetrini interi (\emph{whole slide images}, WSI), combinata con
i progressi dell'intelligenza artificiale e dell'analisi delle immagini, sta creando nuove
opportunità per una valutazione dei campioni tissutali più efficiente, riproducibile e
quantitativa. Per gli specializzandi in patologia e i tecnici di laboratorio, la comprensione di
queste tecnologie non è più facoltativa: è una componente essenziale della pratica professionale
moderna.

Questa dispensa accompagna un corso di cinque lezioni progettato per introdurre i principi e la
pratica della patologia digitale e dell'analisi delle immagini a studenti e specializzandi in
ambito biomedico e medico. Il corso è strutturato in modo da costruire le conoscenze in modo
progressivo: dai fondamenti dell'imaging digitale e della scansione di vetrini interi, attraverso
il controllo qualità e la gestione del flusso di lavoro, fino ai principi dell'intelligenza
artificiale, dell'apprendimento profondo e dell'integrazione multimodale dei dati. Il capitolo
conclusivo affronta le sfide critiche legate al bias, all'interpretabilità, all'etica e al ruolo
futuro del professionista della patologia in un contesto sempre più assistito dall'IA.

Ogni capitolo è concepito per essere letto in abbinamento alla lezione e alla sessione pratica
corrispondenti. Le domande di revisione al termine di ogni capitolo sono destinate
all'autovalutazione; la chiave delle risposte con le relative spiegazioni è fornita nella sezione
finale di questa dispensa.

Il campo della patologia digitale è in rapida evoluzione. Pur avendo compiuto ogni sforzo per
garantire l'accuratezza e l'aggiornamento dei contenuti, gli studenti sono incoraggiati a
consultare la letteratura primaria e le linee guida correnti per le informazioni più aggiornate.
I riferimenti bibliografici principali sono citati nel testo e raccolti nella bibliografia.

\vspace{1em}
\noindent\textit{Ci auguriamo che questa dispensa costituisca una risorsa utile e duratura nel
corso della vostra formazione e della vostra carriera professionale.}

%% Indice
\tableofcontents
\listoffigures

%% ---- TESTO PRINCIPALE ----
\mainmatter

%% Capitoli
%% ============================================================
%%  Capitolo 1 – Introduzione alla Patologia Digitale
%%  Dispensa del Corso di Patologia Post-Laurea
%%  Questo file è incluso con \input{} da main.tex; NON aggiungere il preambolo qui.
%% ============================================================

\chapter{Introduzione alla Patologia Digitale}
\label{ch:1}

\section{Introduzione}

Per oltre un secolo e mezzo, la diagnosi istopatologica si è fondata su una tecnologia di apparente semplicità: una sottile sezione di tessuto, fissata chimicamente, inclusa in paraffina, colorata con ematossilina ed eosina (H\&E) o con un pannello di colorazioni speciali, e montata sotto un vetrino coprioggetto. Il vetrino così ottenuto viene esaminato al microscopio ottico in campo chiaro, e l'occhio esperto del patologo interpreta l'architettura cellulare, la morfologia nucleare e l'organizzazione tissutale per giungere a una diagnosi. Questo flusso di lavoro ha servito la medicina in modo straordinario, eppure presenta limitazioni intrinseche: i vetrini sono oggetti fisici fragili che si deteriorano nel tempo, i microscopi sono vincolati a una singola postazione, i secondi pareri richiedono il trasporto fisico di materiale insostituibile, e la natura soggettiva e analogica della valutazione visiva rende difficile conseguire una riproducibilità quantitativa.

L'affermarsi della \emph{patologia digitale} nel corso degli ultimi due decenni rappresenta la trasformazione più significativa dell'istopatologia diagnostica dall'introduzione dell'immunoistochimica. Nella sua essenza, la patologia digitale sostituisce il microscopio ottico, quale strumento primario di osservazione, con un'immagine digitale ad alta risoluzione dell'intero vetrino --- una \emph{whole slide image} (WSI) --- che può essere archiviata, trasmessa, visualizzata e analizzata computazionalmente. Questo cambiamento non è una mera comodità tecnologica: modifica in modo fondamentale ciò che è possibile nella pratica patologica, nella didattica, nella ricerca e nella garanzia della qualità. La refertazione a distanza, i tumor board multisede, la ricerca su larga scala di biomarcatori e la diagnosi assistita dall'intelligenza artificiale dipendono tutti dall'immagine digitale come tipo di dato fondamentale \cite{bera2019}.

Il ritmo di adozione si è accelerato notevolmente da quando gli enti regolatori di diverse giurisdizioni hanno concesso l'approvazione per la diagnosi primaria mediante WSI. Nel Regno Unito, il programma di trasformazione della patologia di NHS England ha identificato la patologia digitale come priorità strategica, e molti centri medici accademici gestiscono oggi servizi diagnostici completamente o parzialmente digitali. A livello internazionale, la convergenza tra scanner ad alto rendimento a costi accessibili, infrastrutture di archiviazione cloud e software di analisi delle immagini validati ha abbassato le barriere all'adozione per istituzioni di ogni dimensione \cite{song2023}.

Il presente capitolo fornisce le basi concettuali e pratiche utili per  affrontare gli argomenti più avanzati trattati nei capitoli successivi. Prende avvio dalla definizione di patologia digitale e dalla sua collocazione nel più ampio panorama della patologia computazionale. Esamina quindi in dettaglio la tecnologia del whole slide imaging, trattando l'hardware degli scanner, i formati delle immagini, i concetti di risoluzione e i metadati. Il capitolo affronta poi le modalità di integrazione dei sistemi di patologia digitale con i sistemi informativi di laboratorio e l'infrastruttura informatica ospedaliera, per poi concentrarsi sulle responsabilità specifiche del tecnico di laboratorio biomedico nel flusso di lavoro digitale. Due attività pratiche strutturate consolidano i contenuti teorici, e il capitolo si chiude con un riepilogo e una serie di domande di verifica.

\section{Che cos'è la Patologia Digitale?}

La patologia digitale può essere definita come l'acquisizione, la gestione, la condivisione e l'interpretazione di informazioni patologiche --- incluse immagini, annotazioni e dati clinici associati --- in formato digitale. Nel senso più ristretto, il termine si riferisce alla digitalizzazione dei vetrini istologici per produrre whole slide images; nel senso più ampio, comprende l'intero ecosistema informatico che circonda tali immagini, dai sistemi informativi di laboratorio e dalle piattaforme di gestione delle immagini agli algoritmi di intelligenza artificiale e alle reti di telepatologia. Il Royal College of Pathologists e gli organismi equivalenti in Nord America e in Europa hanno adottato definizioni che sottolineano sia la dimensione tecnologica sia quella clinica: la patologia digitale non riguarda semplicemente il possesso di immagini digitali, ma l'utilizzo di tali immagini a supporto della diagnosi, della didattica, della ricerca e del miglioramento della qualità \cite{bankhead2022}.

È importante distinguere la patologia digitale dal campo correlato ma distinto della \emph{patologia computazionale}. La patologia digitale è la disciplina più ampia, che comprende tutti gli aspetti del lavoro con le immagini digitali di vetrini, indipendentemente dal fatto che sia coinvolta l'elaborazione computazionale. La patologia computazionale è una sotto-disciplina che applica metodi computazionali quantitativi --- analisi statistica delle immagini, apprendimento automatico, apprendimento profondo --- per estrarre informazioni da tali immagini che sarebbero difficili o impossibili da ottenere mediante la sola ispezione visiva \cite{campanella2019}. Un patologo che esamina una WSI su schermo pratica la patologia digitale; un modello di apprendimento profondo addestrato a rilevare figure mitotiche nella stessa immagine è un esempio di patologia computazionale. I due campi sono profondamente intrecciati e la distinzione è talvolta sfumata nella letteratura, ma è concettualmente utile mantenerla.

L'ambito della patologia digitale si estende a tutte le sottodiscipline dell'istopatologia e della citopatologia. La patologia chirurgica, la neuropatologia, la dermatopatologia, l'ematopatologia e la citologia hanno tutte sviluppato flussi di lavoro digitali, ciascuno con requisiti tecnici specifici. Le tecniche basate sulla fluorescenza, tra cui l'immunofluorescenza e l'ibridazione fluorescente in situ (FISH), richiedono scanner in grado di acquisire più canali di fluorescenza, aggiungendo complessità sia all'acquisizione sia all'archiviazione. I campioni citologici, spesso preparati come preparati in fase liquida su vetrini di grande formato, presentano sfide di scansione diverse rispetto alle sezioni istologiche standard. La diversità dei tipi di campione implica che una singola soluzione tecnica raramente soddisfi tutte le esigenze, e i tecnici di laboratorio biomedico devono comprendere i requisiti specifici di ciascuna modalità.

La patologia digitale si interseca inoltre con la \emph{telepatologia}, ovvero la pratica di trasmettere immagini patologiche attraverso una rete per consulenza o diagnosi a distanza. La telepatologia precede la moderna WSI --- i primi sistemi utilizzavano microscopi robotici controllati via linea telefonica --- ma la WSI l'ha trasformata in un servizio pratico e scalabile. La telepatologia per sezioni al congelatore, in cui una WSI di una sezione intraoperatoria viene trasmessa a un patologo remoto per una diagnosi immediata, è oggi in uso clinico routinario presso molte istituzioni. La pandemia di COVID-19 ha accelerato l'adozione dei flussi di lavoro di refertazione a distanza, dimostrando che la diagnosi primaria di alta qualità da WSI è fattibile e sicura quando sono in atto adeguate misure di garanzia della qualità \cite{song2023,giaretto2021}.

\section{Whole Slide Imaging: Principi e Strumenti}

\subsection{Funzionamento degli Scanner per Vetrini Interi}

Uno scanner per vetrini interi è, nella sua essenza, un microscopio automatizzato accoppiato a una fotocamera digitale ad alta risoluzione e controllato da un software che acquisisce sistematicamente l'intera superficie di un vetrino. La sfida fondamentale è di natura dimensionale: una sezione istologica standard può misurare $15 \times 10$\,mm o più, eppure le informazioni diagnostiche significative esistono a livello cellulare e subcellulare, richiedendo una risoluzione ottica di circa $0{,}25$\,\textmu m per pixel a ingrandimento $40\times$. Acquisire queste informazioni in una singola esposizione non è fattibile con la tecnologia di sensori attuale; gli scanner acquisiscono pertanto un gran numero di tessere (tile) di immagine sovrapposte o adiacenti e le uniscono computazionalmente per formare la WSI finale.

Il percorso ottico di uno scanner in campo chiaro è molto simile a quello di un microscopio convenzionale a luce trasmessa. La luce bianca proveniente da una sorgente LED o alogena attraversa un condensatore, poi la sezione di tessuto colorata, e infine entra nell'obiettivo. L'obiettivo proietta un'immagine ingrandita su una fotocamera a scansione lineare (negli scanner a scorrimento lineare) o su una fotocamera ad area (negli scanner a tessere). I sistemi a scansione lineare spostano il vetrino in modo continuo sotto una fotocamera fissa, acquisendo una striscia di dati immagine alla volta; i sistemi a tessere spostano il vetrino secondo uno schema a griglia, acquisendo campi visivi discreti. Entrambi gli approcci producono una qualità d'immagine equivalente se correttamente calibrati, ma differiscono nelle caratteristiche di throughput e nella suscettibilità agli artefatti da movimento.

La \emph{scansione in fluorescenza} segue gli stessi principi geometrici, ma sostituisce la sorgente di luce bianca trasmessa con una o più sorgenti di luce di eccitazione --- tipicamente LED ad alta potenza o linee laser --- e inserisce nel percorso ottico appropriati filtri di eccitazione ed emissione. Poiché l'emissione di fluorescenza è di ordini di grandezza più debole rispetto alla trasmissione in campo chiaro, gli scanner a fluorescenza richiedono fotocamere più sensibili (spesso sensori CMOS scientifici raffreddati o CCD), tempi di esposizione più lunghi e un'attenta gestione del fotosbiancamento. La WSI a fluorescenza multicanale, in cui lo stesso vetrino viene acquisito sequenzialmente attraverso più set di filtri, produce uno stack di immagini co-registrate che possono essere unite in un'immagine composita o analizzate canale per canale. Questa capacità è essenziale per i pannelli di immunofluorescenza multiplexata, sempre più utilizzati nella ricerca traslazionale e, più recentemente, nella diagnostica clinica \cite{bera2019}.

\subsection{Risoluzione e Ingrandimento}

La risoluzione di una WSI è determinata principalmente dall'apertura numerica (NA) dell'obiettivo e dalle dimensioni dei pixel del sensore della fotocamera. In pratica, gli scanner sono caratterizzati dall'ingrandimento ottico equivalente del loro obiettivo: gli obiettivi $20\times$ producono una dimensione del pixel di circa $0{,}50$\,\textmu m, mentre gli obiettivi $40\times$ producono circa $0{,}25$\,\textmu m per pixel. La scelta dell'ingrandimento di scansione comporta un compromesso tra il dettaglio dell'immagine e le dimensioni del file. Per la diagnosi routinaria con H\&E, la scansione a $20\times$ è generalmente considerata sufficiente; per le applicazioni che richiedono un fine dettaglio nucleare, il conteggio delle figure mitotiche o la localizzazione subcellulare dei marcatori immunoistochimici, è preferita la scansione a $40\times$. Alcuni scanner offrono un obiettivo a immersione in olio $60\times$ o $100\times$ per applicazioni specialistiche come l'ematologia o la microbiologia.

Le dimensioni dei file risultanti sono considerevoli. Una singola WSI a $20\times$ di una sezione tissutale standard occupa tipicamente da 500\,MB a 2\,GB di dati non compressi; a $40\times$, questo valore quadruplica. La compressione senza perdita o quasi senza perdita (JPEG\,2000, LZW) riduce i requisiti di archiviazione di un fattore da tre a cinque, ma anche i file compressi superano routinariamente i 500\,MB. Un laboratorio diagnostico impegnato che scansiona 200 vetrini al giorno genera oltre 100\,GB di nuovi dati immagine quotidianamente, ponendo richieste significative sull'infrastruttura di archiviazione, sulla larghezza di banda della rete e sui sistemi di archiviazione a lungo termine. Questi vincoli tecnici non sono preoccupazioni periferiche per il tecnico di laboratorio biomedico: influenzano direttamente la progettazione del flusso di lavoro, i protocolli di garanzia della qualità e la scelta della piattaforma di gestione delle immagini.

\subsection{Principali Produttori di Scanner}

Il mercato commerciale degli scanner WSI è dominato da un ristretto numero di produttori, ciascuno dei quali offre strumenti con caratteristiche tecniche, ecosistemi software e formati di file proprietari distinti.

\textbf{Aperio (Leica Biosystems)} produce le serie di scanner ScanScope e GT, che utilizzano un design ottico a scansione lineare e producono immagini nel formato SVS (ScanScope Virtual Slide). Gli scanner Aperio sono stati ampiamente impiegati sia in ambito clinico sia in quello della ricerca e sono stati tra i primi a ricevere l'autorizzazione regolatoria per la diagnosi primaria negli Stati Uniti. Il visualizzatore Aperio ImageScope e la piattaforma eSlide Manager sono comunemente presenti nei laboratori clinici.

\textbf{Hamamatsu Photonics} produce la serie NanoZoomer, che utilizza una fotocamera a scansione lineare con integrazione a ritardo di tempo (TDI) e produce immagini nel formato NDPI (NanoZoomer Digital Pathology Image). Gli scanner Hamamatsu sono particolarmente diffusi nei centri accademici giapponesi ed europei e sono noti per l'elevato throughput e la qualità costante delle immagini.

\textbf{Philips Digital Pathology} (ora parte di Philips Ultrasound \& Image Guided Therapy) produce la soluzione IntelliSite Pathology Solution, che utilizza un approccio di scansione a tessere e produce immagini nel formato iSyntax. Il sistema Philips è stata la prima piattaforma WSI a ricevere la marcatura CE per la diagnosi primaria in Europa ed è stata oggetto di diversi ampi studi di validazione clinica.

\textbf{3DHISTECH} (Budapest, Ungheria) produce la serie Pannoramic di scanner, che producono immagini nel formato MRXS. Gli strumenti 3DHISTECH sono ampiamente utilizzati in ambito di ricerca e didattico e offrono una gamma di opzioni di throughput, dai sistemi a vetrino singolo ai sistemi a carosello ad alta capacità.

\subsection{Software Open-Source per la Visualizzazione e l'Analisi}

Accanto alle piattaforme commerciali, si è sviluppato un vivace ecosistema di software open-source a supporto della visualizzazione, dell'annotazione e dell'analisi delle WSI. Due strumenti rivestono particolare importanza per il tecnico di laboratorio biomedico e il ricercatore in patologia.

\textbf{QuPath} (Quantitative Pathology \& Bioimage Analysis) è un'applicazione multipiattaforma basata su Java, sviluppata presso la Queen's University Belfast e successivamente presso l'Università di Edimburgo \cite{bankhead2017}. QuPath fornisce un visualizzatore WSI completo con supporto per tutti i principali formati proprietari (tramite la libreria OpenSlide), insieme a un'ampia suite di strumenti di analisi delle immagini che include il rilevamento cellulare, la classificazione tissutale, la quantificazione della fluorescenza e la programmazione tramite un'API basata su Groovy. La combinazione di accessibilità, estensibilità e rigore statistico ne ha fatto lo strumento open-source più utilizzato nella ricerca in patologia digitale \cite{bankhead2022}. QuPath è disponibile gratuitamente all'indirizzo \url{https://qupath.github.io} ed è attivamente mantenuto con rilasci regolari.

\textbf{OpenSlide} è una libreria C (con binding per Python, Java e altri linguaggi) che fornisce un'interfaccia di programmazione unificata per la lettura di file WSI in più formati proprietari, tra cui SVS, NDPI, MRXS, SCN e altri. Sviluppata presso la Carnegie Mellon University, OpenSlide astrae i dettagli specifici del formato della struttura dei file di ciascun produttore, consentendo agli sviluppatori di scrivere codice di analisi delle immagini che funziona su diverse piattaforme di scanner senza modifiche. È una dipendenza fondamentale di QuPath e di molti altri strumenti open-source, e la comprensione delle sue capacità e limitazioni è importante per chiunque sviluppi pipeline di analisi personalizzate.

\section{Formati delle Immagini Digitali, Risoluzione e Metadati}

\subsection{Struttura dei Pixel e Rappresentazione del Colore}

Al livello più fondamentale, un'immagine digitale è un array rettangolare di \emph{pixel} (elementi immagine), ciascuno dei quali memorizza un valore numerico che rappresenta il colore o l'intensità della luce in quella posizione. Nell'istologia standard in campo chiaro, le immagini vengono acquisite nel modello di colore RGB: ogni pixel è rappresentato da tre valori interi a 8 bit, uno per ciascuno dei canali di colore rosso, verde e blu, con ogni valore compreso tra 0 (nessun contributo da quel canale) e 255 (contributo massimo). La combinazione dei tre valori di canale definisce un colore da una tavolozza di $256^3 = 16{,}777{,}216$ colori possibili, sufficiente a rappresentare l'intera gamma di colori riscontrati nell'H\&E e nella maggior parte delle colorazioni speciali.

La convenzione degli 8 bit per canale è profondamente radicata nell'ecosistema WSI e nelle librerie di elaborazione delle immagini che operano sui dati WSI. È tuttavia importante essere consapevoli dei suoi limiti. Le immagini a fluorescenza, in particolare quelle acquisite con fotocamere di livello scientifico, vengono spesso acquisite a una profondità di 12 o 16 bit per canale, fornendo rispettivamente 4.096 o 65.536 livelli di intensità. Questo maggiore intervallo dinamico è importante per le misurazioni quantitative della fluorescenza, in cui deve essere preservata la relazione tra l'intensità del pixel e la concentrazione del fluoroforo. Quando le immagini a fluorescenza a 16 bit vengono convertite a 8 bit per la visualizzazione o l'archiviazione, le informazioni vengono irrimediabilmente perse a meno che la conversione non venga eseguita con attenzione e con un'appropriata scalatura. I tecnici di laboratorio biomedico coinvolti nella scansione a fluorescenza devono comprendere questa distinzione e assicurarsi che le pipeline di analisi operino alla profondità di bit appropriata.

\subsection{Piramidi di Immagini e Formati a Tessere}

Una WSI non può essere caricata interamente in memoria su una workstation standard; una scansione a $40\times$ di un grande campione di resezione può contenere $100{,}000 \times 80{,}000$ pixel o più, richiedendo diversi gigabyte di RAM per essere mantenuta non compressa. Per consentire una navigazione e una visualizzazione efficienti a più livelli di zoom, i file WSI vengono archiviati come \emph{piramidi di immagini}: una serie di versioni progressivamente sottocampionate dell'immagine a piena risoluzione, ciascuna con la metà della risoluzione lineare del livello precedente. Un'applicazione di visualizzazione carica solo il livello della piramide appropriato al fattore di zoom corrente e, all'interno di quel livello, solo le tessere che rientrano nel viewport corrente. Questa architettura consente una navigazione fluida e reattiva di immagini gigapixel su hardware con memoria limitata.

La piramide di immagini è tipicamente archiviata in un formato \emph{a tessere}, in cui ogni livello della piramide è suddiviso in una griglia regolare di piccole tessere (comunemente $256 \times 256$ o $512 \times 512$ pixel), ciascuna compressa in modo indipendente. La compressione indipendente delle tessere consente a un visualizzatore di decomprimere e visualizzare una singola tessera senza leggere l'intero file, riducendo drasticamente l'overhead di I/O, ovvero il tempo, i cicli della CPU e la memoria che viene utilizzata in operazioni di input/output che non contribuiscono al raggiungimento dell'operazione. I principali formati WSI proprietari --- SVS (Aperio), NDPI (Hamamatsu), MRXS (3DHISTECH) e iSyntax (Philips) --- implementano tutti varianti di questa architettura a piramide di tessere, sebbene differiscano nei loro specifici schemi di compressione, nelle dimensioni delle tessere e nelle strutture dei contenitori di file.

\textbf{OME-TIFF} (Open Microscopy Environment Tagged Image File Format) è un formato aperto e indipendente dal produttore che ha acquisito una posizione di rilievo come standard per lo scambio di WSI. OME-TIFF archivia i dati immagine in un contenitore TIFF standard con un blocco di metadati XML strutturato che codifica informazioni sulle dimensioni dell'immagine, le identità dei canali, le dimensioni fisiche dei pixel e i parametri di acquisizione in un vocabolario standardizzato. La variante BigTIFF di OME-TIFF supporta file di dimensioni superiori a 4\,GB, il che è essenziale per le WSI. L'adozione di OME-TIFF come formato di scambio comune è incoraggiata dagli enti regolatori e dalle organizzazioni di standardizzazione come mezzo per ridurre la dipendenza dai produttori e migliorare l'accessibilità dei dati a lungo termine \cite{song2023}.

\subsection{Concetti di Risoluzione: Micron per Pixel}

La risoluzione fisica di una WSI è espressa in \emph{micron per pixel} (\textmu m/px), che mette in relazione le dimensioni in pixel dell'immagine con le dimensioni fisiche del tessuto. Una scansione a $20\times$ ha tipicamente una risoluzione di circa $0{,}50$\,\textmu m/px, il che significa che ogni pixel corrisponde a un quadrato di tessuto di $0{,}5 \times 0{,}5$\,\textmu m. A $40\times$, questo valore si dimezza a circa $0{,}25$\,\textmu m/px. Il valore dei micron per pixel è memorizzato nei metadati dell'immagine ed è essenziale per qualsiasi misurazione quantitativa: le aree cellulari, i diametri nucleari, le distanze intercellulari e le dimensioni delle ghiandole sono tutti privi di significato senza la conoscenza della scala fisica.

È importante distinguere tra \emph{risoluzione ottica} (la caratteristica più piccola che l'obiettivo può risolvere, determinata dal criterio di Rayleigh e dalla NA dell'obiettivo) e \emph{risoluzione di campionamento digitale} (la dimensione del pixel). La teoria del campionamento di Nyquist richiede che la dimensione del pixel non sia superiore alla metà del limite di risoluzione ottica per evitare artefatti di aliasing. In pratica, la maggior parte degli scanner WSI è progettata in modo che la dimensione del pixel sia leggermente inferiore al limite di risoluzione ottica, garantendo che l'immagine digitale rappresenti fedelmente tutte le informazioni che l'ottica è in grado di acquisire. Il sovracampionamento oltre questo punto aumenta le dimensioni del file senza migliorare il contenuto informativo diagnostico.

\subsection{Metadati Incorporati nei File WSI}

I file WSI contengono ricchi metadati oltre ai dati pixel stessi. Questi metadati includono tipicamente: il modello e il numero di serie dello scanner; la data e l'ora di acquisizione; l'ingrandimento dell'obiettivo e l'apertura numerica; la dimensione fisica del pixel in micron; l'algoritmo di compressione e le impostazioni di qualità; l'immagine dell'etichetta del vetrino (una fotografia a bassa risoluzione dell'etichetta del vetrino di vetro); e, in alcuni formati, un'immagine macro che mostra l'intero vetrino a risoluzione molto bassa per scopi di navigazione. Nelle scansioni a fluorescenza, i metadati includono inoltre le lunghezze d'onda di eccitazione ed emissione per ciascun canale, i tempi di esposizione e le impostazioni di guadagno.

L'integrità dei metadati è fondamentale sia per le applicazioni cliniche sia per quelle di ricerca. In un contesto clinico, i metadati dello scanner forniscono una traccia di audit che collega l'immagine digitale al vetrino fisico e alla cartella del paziente. In un contesto di ricerca, i metadati consentono la riproducibilità: un ricercatore che analizza una coorte di WSI deve essere in grado di verificare che tutte le immagini siano state acquisite in condizioni comparabili. I tecnici di laboratorio biomedico responsabili della scansione devono pertanto assicurarsi che i registri di calibrazione dello scanner siano mantenuti, che le etichette dei vetrini siano correttamente associate agli identificativi del paziente prima della scansione e che tutti i campi di metadati richiesti dal sistema di gestione della qualità del laboratorio siano correttamente compilati. Gli errori nei metadati --- in particolare le discordanze degli identificativi del paziente --- costituiscono un grave rischio per la sicurezza del paziente e devono essere trattati con lo stesso rigore di qualsiasi altro errore pre-analitico in laboratorio.

\section{Integrazione con i Sistemi Informativi di Laboratorio}

\subsection{Sistemi Informativi di Laboratorio e LIMS}

Un \emph{Sistema Informativo di Laboratorio} (in inglese Laboratory Information System, LIS) è la piattaforma software che gestisce i dati amministrativi e clinici associati ai campioni patologici durante il loro ciclo di vita in laboratorio. Il LIS riceve le richieste dai team clinici (elettronicamente tramite sistemi di inserimento degli ordini o, in alcuni contesti, su moduli di richiesta cartacei), assegna numeri di accettazione univoci ai campioni, ne traccia il progresso attraverso il flusso di lavoro di laboratorio, archivia il referto del patologo e trasmette il referto finale al clinico richiedente. In un laboratorio di istopatologia, il LIS è il registro autorevole di ogni campione: collega i dati demografici del paziente, la storia clinica, la descrizione macroscopica, gli identificativi di blocco e livello, le richieste di colorazione e il referto diagnostico in un unico registro di caso coerente.

Un \emph{Laboratory Information Management System} (LIMS) è un concetto correlato ma distinto, più comunemente utilizzato in ambito di ricerca e industriale, che enfatizza il tracciamento dei campioni, la catena di custodia e la gestione dei dati piuttosto che la refertazione clinica. In pratica, i termini LIS e LIMS vengono talvolta usati in modo intercambiabile nella letteratura patologica, e molte piattaforme moderne sfumano la distinzione combinando funzionalità di refertazione clinica con capacità di gestione dei dati di ricerca. Ai fini del presente capitolo, il termine LIS verrà utilizzato per riferirsi al sistema di laboratorio clinico.

L'integrazione della patologia digitale con il LIS è un prerequisito per un'operatività clinica sicura ed efficiente. Quando un vetrino di vetro viene scansionato, la WSI risultante deve essere collegata in modo inequivocabile al caso corretto nel LIS. Questo collegamento viene tipicamente realizzato codificando il numero di accettazione del LIS e l'identificativo di blocco/livello in un codice a barre sull'etichetta del vetrino, che lo scanner legge automaticamente all'inizio di ogni scansione. Il software dello scanner trasmette quindi questo identificativo al sistema di gestione delle immagini (talvolta denominato sistema di gestione della patologia digitale o DPMS), che crea un record che collega il file WSI al caso nel LIS. Il patologo, quando apre un caso per la refertazione, visualizza la WSI insieme alle informazioni cliniche del LIS in un visualizzatore integrato, senza dover navigare tra sistemi separati.

\subsection{DICOM per la Patologia}

\emph{DICOM} (Digital Imaging and Communications in Medicine) è lo standard internazionale per l'archiviazione e la trasmissione di immagini mediche e delle informazioni associate. Originariamente sviluppato per la radiologia, DICOM è stato esteso alla patologia attraverso il Supplemento 145 (Whole Slide Microscopy Image) e i supplementi successivi. DICOM per la patologia definisce un formato di file standardizzato per i dati WSI che codifica sia i dati pixel (in una struttura a piramide di tessere) sia un insieme completo di attributi di metadati in un dizionario di dati standardizzato, inclusi i dati demografici del paziente, gli identificativi del campione, le informazioni sulla colorazione e i parametri di acquisizione.

L'adozione di DICOM come formato standard per le WSI cliniche offre diversi vantaggi importanti. In primo luogo, consente l'interoperabilità tra sistemi di produttori diversi: una WSI conforme a DICOM prodotta da uno scanner può, in linea di principio, essere visualizzata e refertata utilizzando qualsiasi visualizzatore conforme a DICOM, senza necessità di conversione del formato. In secondo luogo, si integra naturalmente con l'infrastruttura DICOM esistente dei reparti di radiologia ospedaliera, inclusi i PACS (Picture Archiving and Communication Systems) e i VNA (Vendor Neutral Archives), consentendo potenzialmente l'archiviazione e la gestione delle immagini patologiche all'interno della stessa piattaforma di imaging aziendale delle immagini radiologiche. In terzo luogo, i meccanismi maturi di sicurezza e traccia di audit di DICOM supportano i requisiti di governance della pratica clinica \cite{song2023}.

\subsection{Standard HL7 e Integrazione con l'Informatica Ospedaliera}

\emph{HL7} (Health Level Seven) è un insieme di standard internazionali per lo scambio di informazioni sanitarie cliniche e amministrative tra applicazioni software. La messaggistica HL7 versione 2 (HL7v2) è lo standard più ampiamente impiegato per l'integrazione del LIS negli ambienti informatici ospedalieri esistenti: i messaggi di ordine elettronico (ORM) vengono inviati dal sistema di inserimento degli ordini dell'ospedale al LIS quando viene richiesto un campione, e i messaggi di risultato (ORU) vengono inviati dal LIS al sistema clinico quando un referto viene autorizzato. HL7 FHIR (Fast Healthcare Interoperability Resources) è uno standard più recente basato su servizi web, sempre più adottato per le nuove integrazioni, in particolare nelle applicazioni sanitarie cloud e mobili.

Per la patologia digitale, la messaggistica HL7 fornisce il meccanismo attraverso il quale le informazioni sul campione e sul paziente fluiscono dai sistemi amministrativi dell'ospedale nel LIS e, da lì, nel sistema di gestione della patologia digitale. Un'interfaccia HL7 correttamente configurata garantisce che, quando un campione arriva in laboratorio, i dati demografici del paziente e le informazioni cliniche associate siano già presenti nel LIS, riducendo il rischio di errori di trascrizione. Consente inoltre la trasmissione automatica dei referti patologici alla cartella clinica elettronica del clinico richiedente, chiudendo il ciclo tra laboratorio e assistenza clinica. I tecnici di laboratorio biomedico non sono tipicamente responsabili della configurazione delle interfacce HL7 --- questo è il dominio degli specialisti IT e degli amministratori del LIS --- ma devono comprendere i flussi di informazioni che queste interfacce supportano al fine di riconoscere e segnalare i guasti di integrazione quando si verificano.

\section{Il Ruolo del Tecnico nel Flusso di Lavoro Digitale}

\subsection{Gestione dei Campioni e Preparazione dei Vetrini}

La qualità di un'immagine di patologia digitale è in ultima analisi determinata dalla qualità del vetrino di vetro da cui deriva. Nessuna quantità di post-elaborazione computazionale può recuperare le informazioni diagnostiche perse durante la fissazione, il processamento, l'inclusione, il sezionamento o la colorazione del tessuto. Le responsabilità del tecnico di laboratorio biomedico nella fase pre-analitica del flusso di lavoro digitale sono pertanto identiche a quelle in un laboratorio convenzionale: garantire un'adeguata fissazione, il corretto orientamento del tessuto durante l'inclusione, uno spessore di sezione costante e il rispetto dei protocolli di colorazione validati. Queste variabili pre-analitiche influenzano direttamente l'aspetto dell'immagine digitale e, di conseguenza, le prestazioni di qualsiasi algoritmo di analisi automatica delle immagini applicato ad essa.

Ulteriori considerazioni emergono specificamente nel contesto della scansione digitale. I vetrini devono essere privi di bolle d'aria, pieghe e galleggianti, che sono artefatti comuni nell'istologia convenzionale ma sono particolarmente problematici per gli scanner automatizzati, che potrebbero non riuscire a mettere a fuoco correttamente nelle regioni interessate o potrebbero contrassegnare il vetrino come insoddisfacente. La qualità del coprioggetto è fondamentale: un mezzo di montaggio non uniforme, un eccesso di mezzo che si estende oltre il bordo del coprioggetto o coprioggetti non completamente aderenti possono tutti causare problemi di messa a fuoco. L'uso di vetrini caricati positivamente e di adesivi appropriati riduce il rischio di distacco della sezione durante la colorazione, che comporterebbe una perdita di tessuto visibile nell'immagine digitale. Le etichette dei vetrini devono essere stampate in modo chiaro e duraturo, poiché il lettore di codici a barre dello scanner deve essere in grado di leggere l'etichetta in modo affidabile; le etichette scritte a mano o stampate su carta termica sbiadita sono una fonte comune di errori di scansione.

\subsection{Operazioni di Scansione e Controllo della Qualità}

Il tecnico di laboratorio biomedico che opera lo scanner è responsabile del corretto caricamento dei vetrini, dell'avvio dei lavori di scansione, del monitoraggio del processo di scansione e dell'esecuzione dei controlli di qualità sulle scansioni completate prima che vengano rilasciate per la refertazione. La maggior parte degli scanner moderni esegue una mappa di messa a fuoco pre-scansione automatizzata, in cui lo scanner campiona la messa a fuoco in una griglia di punti sul vetrino per costruire un modello della topografia del vetrino; questa mappa di messa a fuoco viene poi utilizzata per mantenere la messa a fuoco durante la scansione principale. Il tecnico di laboratorio biomedico deve avere familiarità con gli indicatori di qualità della messa a fuoco dello scanner e deve essere in grado di identificare le scansioni in cui la mappa di messa a fuoco ha fallito o in cui la messa a fuoco è stata persa durante la scansione.

Il controllo della qualità delle WSI completate deve seguire un protocollo documentato. Come minimo, questo dovrebbe includere: la verifica che la WSI sia correttamente collegata al caso previsto nel LIS; l'ispezione visiva dell'immagine in miniatura per artefatti grossolani (pieghe del tessuto, bolle d'aria, irregolarità di colorazione); la valutazione della qualità della messa a fuoco sull'area del tessuto, inclusi i bordi del tessuto dove i problemi di messa a fuoco sono più comuni; e la conferma che la piramide di immagini sia completa e che tutte le tessere siano state generate con successo. Molti laboratori utilizzano strumenti automatizzati di valutazione della qualità --- integrati nel software dello scanner o forniti da piattaforme di terze parti --- che segnalano potenziali problemi di qualità per la revisione umana. Il tecnico di laboratorio biomedico deve comprendere i criteri utilizzati da questi strumenti e deve essere in grado di formulare un giudizio informato su se un vetrino segnalato richieda una nuova scansione o possa essere rilasciato con un'annotazione di qualità.

\subsection{Gestione dei Dati e Governance}

Il tecnico di laboratorio biomedico svolge un ruolo importante nella gestione dei dati WSI durante il loro ciclo di vita. Ciò include garantire che i file di scansione siano correttamente denominati e archiviati nelle posizioni designate sul sistema di gestione delle immagini, che le procedure di backup e archiviazione siano seguite e che le politiche di conservazione dei dati siano rispettate. Nel Regno Unito, la conservazione del materiale patologico --- incluse le immagini digitali --- è disciplinata dall'Human Tissue Act 2004 e dai relativi codici di pratica, nonché dalle politiche di gestione dei documenti del NHS. Le immagini digitali dei vetrini diagnostici sono considerate parte della cartella del paziente e devono essere conservate per lo stesso periodo del corrispondente referto cartaceo o elettronico, tipicamente un minimo di 30 anni per i pazienti adulti.

La sicurezza dei dati e la governance delle informazioni sono responsabilità aggiuntive che ricadono in parte nell'ambito del tecnico di laboratorio biomedico. I file WSI contengono, o sono collegati a, informazioni identificative del paziente e sono pertanto soggetti ai requisiti del Regolamento Generale sulla Protezione dei Dati del Regno Unito (UK GDPR) e del Data Protection Act 2018. L'accesso al sistema di gestione delle immagini deve essere controllato tramite autorizzazioni di accesso basate sui ruoli, e qualsiasi trasferimento di dati WSI al di fuori dell'organizzazione --- ad esempio, per la garanzia della qualità esterna, la ricerca o la telepatologia --- deve essere condotto in conformità con un accordo di condivisione dei dati e, ove applicabile, con il consenso appropriato del paziente o una base giuridica per il trattamento. Il tecnico di laboratorio biomedico deve avere familiarità con le politiche di governance delle informazioni della propria organizzazione e deve segnalare qualsiasi preoccupazione relativa alla sicurezza dei dati al responsabile della governance delle informazioni competente.

\section{Attività Pratiche}

Le seguenti due attività pratiche sono progettate per consolidare i contenuti teorici di questo capitolo attraverso un coinvolgimento diretto con gli strumenti e i flussi di lavoro della patologia digitale. Sono adatte per essere completate in una sessione di laboratorio o di sala computer supervisionata della durata di circa due ore.

\subsection*{Attività Pratica 1: Esplorazione delle Whole Slide Images con QuPath}

\textbf{Obiettivo:} Acquisire familiarità con il visualizzatore di patologia digitale QuPath e comprendere le caratteristiche principali di una whole slide image.

\textbf{Materiali:} Un computer con QuPath (versione 0.5 o successiva) installato; alcune (sono file enormi, scaricatene pochi per volta) scaricabili ai seguenti link (\href{https://openslide.cs.cmu.edu/download/openslide-testdata/Aperio/CMU-1.svs}{singola WSI in H\&E}, \href{https://zenodo.org/records/18302140}{WSI a fluorescenza}, \href{https://openslide.cs.cmu.edu/download/openslide-testdata/}{tutto il testdata di OpenSlide})

).

\textbf{Procedura:}
\begin{enumerate}
  \item Avviare QuPath e aprire un file WSI fornito utilizzando \texttt{File > Open}. Osservare la vista in miniatura e i controlli di navigazione della piramide di immagini.
  \item Navigare l'immagine a più livelli di zoom utilizzando la rotella di scorrimento e il cursore dello zoom. Notare il cambiamento nel dettaglio dell'immagine mentre si passa dalla vista macro al livello cellulare.
  \item Utilizzare la scheda \texttt{Image} nel pannello di sinistra per ispezionare i metadati dell'immagine, inclusa la dimensione del pixel in micron, il numero di livelli della piramide e le dimensioni dell'immagine in pixel.
  \item Utilizzare lo strumento di misurazione (\texttt{Analyse > Measurements}) per misurare il diametro di un nucleo rappresentativo e la larghezza del lume di una ghiandola. Registrare le misurazioni e convertirle da pixel a micron utilizzando la dimensione del pixel dai metadati.
  \item Utilizzare gli strumenti di annotazione per disegnare una regione di interesse attorno a un'area di interesse (ad esempio, un nido tumorale o un aggregato linfoide). Notare l'area dell'annotazione in micron quadrati come riportato da QuPath.
  \item Se è disponibile una WSI a fluorescenza, aprirla e utilizzare il pannello \texttt{Brightness/Contrast} per regolare le impostazioni di visualizzazione per ciascun canale in modo indipendente. Osservare come l'immagine composita cambia quando i singoli canali vengono attivati e disattivati.
  \item Redigere una breve relazione (di circa 300 parole) che descriva l'immagine esaminata, le misurazioni effettuate e gli eventuali problemi di qualità osservati.
\end{enumerate}

\textbf{Risultati di apprendimento:} Al termine di questa attività, gli studenti dovrebbero essere in grado di navigare una WSI in QuPath, interpretare i metadati dell'immagine, eseguire misurazioni di base e identificare i problemi comuni di qualità dell'immagine.

\subsection*{Attività Pratica 2: Mappatura del Flusso di Lavoro dal Campione al Vetrino Digitale}

\textbf{Obiettivo:} Comprendere il flusso di lavoro completo della patologia digitale dalla ricezione del campione alla disponibilità dell'immagine per la refertazione, e identificare i principali punti di controllo della qualità e i passaggi di informazioni in ciascuna fase.

\textbf{Materiali:} Lavagna o carta di grandi dimensioni; pennarelli colorati; una copia delle procedure operative standard del laboratorio per la ricezione dei campioni, il processamento tissutale, il sezionamento, la colorazione, la scansione e il controllo della qualità delle immagini (fornita dal coordinatore del corso).

\textbf{Procedura:}
\begin{enumerate}
  \item Lavorando in piccoli gruppi (3--4 studenti), costruire un diagramma di flusso del processo del flusso di lavoro completo della patologia digitale, dal momento in cui un campione chirurgico arriva in laboratorio al momento in cui la WSI è disponibile per la refertazione da parte del patologo. Includere tutti i passaggi principali, i punti di decisione (ad esempio, ``La qualità della scansione è accettabile?'') e i flussi di informazioni (ad esempio, assegnazione del numero di accettazione del LIS, etichettatura con codice a barre).
  \item Per ogni fase del flusso di lavoro, identificare: (a) il membro del personale responsabile; (b) i criteri di qualità chiave che devono essere soddisfatti prima di procedere alla fase successiva; (c) le conseguenze di un fallimento in questa fase per la qualità finale dell'immagine e per la sicurezza del paziente.
  \item Identificare almeno tre punti nel flusso di lavoro in cui potrebbe verificarsi un incidente di sicurezza del paziente (ad esempio, errata etichettatura del campione, mancata corrispondenza tra scansione e caso, problema di messa a fuoco non rilevato al controllo della qualità) e descrivere i controlli in atto per prevenire ciascun incidente.
  \item Confrontare il proprio diagramma di flusso con quelli prodotti dagli altri gruppi e discutere le eventuali differenze. Considerare come il flusso di lavoro potrebbe differire tra un grande centro terziario e un ospedale distrettuale generale più piccolo.
\end{enumerate}

\textbf{Risultati di apprendimento:} Al termine di questa attività, gli studenti dovrebbero essere in grado di descrivere il flusso di lavoro completo della patologia digitale, identificare i principali punti di controllo della qualità e articolare le implicazioni per la sicurezza del paziente dei fallimenti in ciascuna fase.

\section{Riepilogo}

Il presente capitolo ha introdotto il campo della patologia digitale, tracciandone le origini nelle limitazioni della microscopia convenzionale su vetrino e delineandone lo sviluppo in una disciplina clinica e di ricerca matura. I seguenti punti chiave devono essere ritenuti:

\begin{itemize}
  \item La patologia digitale comprende l'acquisizione, la gestione, la condivisione e l'interpretazione delle informazioni patologiche in forma digitale, con il whole slide imaging come fondamento tecnologico.
  \item La patologia computazionale è una sotto-disciplina della patologia digitale che applica metodi computazionali quantitativi, inclusi l'apprendimento automatico e l'apprendimento profondo, per estrarre informazioni dalle immagini digitali \cite{campanella2019,bera2019}.
  \item Gli scanner per vetrini interi digitalizzano i vetrini di vetro acquisendo un gran numero di tessere di immagine e unendole in una piramide di immagini gigapixel. Le modalità di scansione in campo chiaro e a fluorescenza servono diverse applicazioni diagnostiche e di ricerca.
  \item I principali produttori commerciali di scanner includono Aperio/Leica, Hamamatsu, Philips e 3DHISTECH, ciascuno dei quali produce immagini in formati proprietari. Gli strumenti open-source, in particolare QuPath \cite{bankhead2017} e OpenSlide, forniscono un accesso indipendente dal produttore a queste immagini.
  \item I file WSI archiviano i dati pixel come piramidi di immagini a tessere, tipicamente in formato RGB a 8 bit per le immagini in campo chiaro. La risoluzione fisica è espressa in micron per pixel ed è essenziale per le misurazioni quantitative. OME-TIFF è un importante standard aperto per lo scambio di WSI.
  \item La patologia digitale si integra con l'infrastruttura informatica ospedaliera attraverso i sistemi LIS/LIMS, gli standard DICOM per l'archiviazione e la trasmissione delle immagini e la messaggistica HL7 per lo scambio di dati clinici.
  \item Il tecnico di laboratorio biomedico svolge un ruolo centrale nel flusso di lavoro digitale, con responsabilità che spaziano dalla preparazione dei vetrini alle operazioni di scansione, dal controllo della qualità delle immagini alla gestione dei dati e alla governance.
\end{itemize}

\section{Domande di Verifica}

\begin{enumerate}

  \item \label{ch1:q1} Quale delle seguenti affermazioni descrive meglio il principale vantaggio della patologia digitale rispetto alla microscopia convenzionale su vetrino per la consulenza a distanza?
  \begin{enumerate}[label=\Alph*.]
    \item Le immagini digitali hanno una risoluzione ottica superiore rispetto ai vetrini di vetro osservati al microscopio.
    \item Le whole slide images possono essere trasmesse elettronicamente, consentendo l'accesso remoto senza il trasporto fisico dei vetrini.
    \item La patologia digitale elimina la necessità della fissazione e del processamento tissutale.
    \item L'analisi computazionale delle immagini digitali è sempre più accurata della valutazione visiva di un patologo.
  \end{enumerate}

  \item \label{ch1:q2} Una whole slide image acquisita a ingrandimento $20\times$ ha una dimensione del pixel di $0{,}50$\,\textmu m. Qual è la dimensione approssimativa del pixel di una scansione a $40\times$ dello stesso tessuto?
  \begin{enumerate}[label=\Alph*.]
    \item $1{,}00$\,\textmu m
    \item $0{,}50$\,\textmu m
    \item $0{,}25$\,\textmu m
    \item $0{,}10$\,\textmu m
  \end{enumerate}

  \item \label{ch1:q3} Quale dei seguenti formati di file è uno standard aperto e indipendente dal produttore specificamente progettato per lo scambio di whole slide images?
  \begin{enumerate}[label=\Alph*.]
    \item SVS
    \item NDPI
    \item OME-TIFF
    \item MRXS
  \end{enumerate}

  \item \label{ch1:q4} Nel modello di colore RGB utilizzato per le WSI standard in campo chiaro, ogni canale di colore è rappresentato da un intero a 8 bit. Qual è il numero totale di colori possibili che possono essere rappresentati?
  \begin{enumerate}[label=\Alph*.]
    \item 256
    \item 65.536
    \item 16.777.216
    \item 4.294.967.296
  \end{enumerate}

  \item \label{ch1:q5} Quale libreria software open-source fornisce un'interfaccia di programmazione unificata per la lettura di file WSI in più formati proprietari, inclusi SVS, NDPI e MRXS?
  \begin{enumerate}[label=\Alph*.]
    \item QuPath
    \item ImageJ
    \item OpenSlide
    \item ITK
  \end{enumerate}

  \item \label{ch1:q6} Un tecnico di laboratorio biomedico nota che una WSI completata mostra una grande regione di tessuto sfocato al centro della sezione. Qual è la causa più probabile di questo artefatto?
  \begin{enumerate}[label=\Alph*.]
    \item Il tessuto è stato sovrafissato in formalina.
    \item Lo scanner non è riuscito a ottenere la corretta messa a fuoco in quella regione durante la scansione.
    \item La piramide di immagini non è stata completamente generata.
    \item Il vetrino è stato scansionato a un ingrandimento errato.
  \end{enumerate}

  \item \label{ch1:q7} Quale standard internazionale disciplina l'archiviazione e la trasmissione delle whole slide images nell'infrastruttura informatica ospedaliera, consentendo l'interoperabilità con i sistemi PACS radiologici?
  \begin{enumerate}[label=\Alph*.]
    \item HL7 FHIR
    \item DICOM
    \item OME-XML
    \item IHE XDS
  \end{enumerate}

  \item \label{ch1:q8} Qual è lo scopo principale della struttura a piramide di immagini utilizzata nei formati di file WSI?
  \begin{enumerate}[label=\Alph*.]
    \item Archiviare più canali di colorazione in un unico file.
    \item Consentire la visualizzazione e la navigazione efficienti di immagini gigapixel caricando solo il livello di risoluzione e le tessere necessari per la vista corrente.
    \item Comprimere i dati immagine utilizzando un algoritmo senza perdita.
    \item Incorporare i metadati del paziente insieme ai dati pixel.
  \end{enumerate}

  \item \label{ch1:q9} Quale delle seguenti affermazioni distingue meglio la \emph{patologia computazionale} dalla \emph{patologia digitale}?
  \begin{enumerate}[label=\Alph*.]
    \item La patologia computazionale utilizza vetrini di vetro, mentre la patologia digitale utilizza whole slide images.
    \item La patologia computazionale applica metodi computazionali quantitativi per estrarre informazioni dalle immagini digitali, mentre la patologia digitale è la disciplina più ampia che si occupa delle immagini digitali di vetrini.
    \item La patologia computazionale è utilizzata solo nella ricerca, mentre la patologia digitale è utilizzata solo nella pratica clinica.
    \item La patologia computazionale richiede l'imaging a fluorescenza, mentre la patologia digitale utilizza l'imaging in campo chiaro.
  \end{enumerate}

  \item \label{ch1:q10} In base alla legislazione del Regno Unito e alla politica di gestione dei documenti del NHS, per quanto tempo approssimativamente devono essere conservate le immagini digitali dei vetrini di istopatologia diagnostica per i pazienti adulti?
  \begin{enumerate}[label=\Alph*.]
    \item 7 anni
    \item 10 anni
    \item 20 anni
    \item 30 anni
  \end{enumerate}

\end{enumerate}

\begin{figure}[htbp]
  \centering
  \usetikzlibrary{arrows.meta, positioning, calc}

\begin{tikzpicture}[
    node distance = 0.55cm,
    every node/.style = {font=\small\sffamily},
    proc/.style = {
        rectangle,
        rounded corners=4pt,
        fill={rgb,255:red,46;green,134;blue,171},
        text=white,
        minimum width=1.55cm,
        minimum height=0.75cm,
        align=center,
        font=\small\sffamily\bfseries
    },
    arr/.style = {
        -{Stealth[length=5pt,width=4pt]},
        line width=0.8pt,
        color={rgb,255:red,46;green,134;blue,171}
    }
]

%% Title
\node[font=\large\sffamily\bfseries,
      text={rgb,255:red,30;green,60;blue,90}]
    (title) at (0,0) {Digital Pathology Workflow};

%% Nodes placed in a row below the title
\node[proc, below=0.55cm of title, xshift=-8.1cm] (specimen)  {Specimen};
\node[proc, right=0.55cm of specimen]              (tissue)    {Tissue\\Processing};
\node[proc, right=0.55cm of tissue]                (slide)     {Glass\\Slide};
\node[proc, right=0.55cm of slide]                 (scanning)  {Scanning};
\node[proc, right=0.55cm of scanning]              (wsi)       {WSI};
\node[proc, right=0.55cm of wsi]                   (analysis)  {Image\\Analysis};
\node[proc, right=0.55cm of analysis]              (report)    {Report};

%% Arrows
\draw[arr] (specimen) -- (tissue);
\draw[arr] (tissue)   -- (slide);
\draw[arr] (slide)    -- (scanning);
\draw[arr] (scanning) -- (wsi);
\draw[arr] (wsi)      -- (analysis);
\draw[arr] (analysis) -- (report);

%% Subtle baseline rule
\draw[line width=0.4pt, color=gray!40]
    ($(specimen.south west)+(-0.1,-0.25)$) --
    ($(report.south east)+(0.1,-0.25)$);

\end{tikzpicture}

  \caption{Panoramica del flusso di lavoro della patologia digitale. Un campione chirurgico o bioptico viene sottoposto al processamento tissutale standard e sezionato su un vetrino di vetro. Il vetrino colorato viene caricato in uno scanner per vetrini interi, che produce una whole slide image (WSI) gigapixel archiviata su un server di gestione delle immagini. La WSI è quindi disponibile per l'analisi computazionale delle immagini e per la refertazione da parte del patologo, con il referto diagnostico finale trasmesso al team clinico tramite il sistema informativo di laboratorio (LIS).}
  \label{fig:wsi_workflow}
\end{figure}

%% ============================================================
%%  Capitolo 2 – Flusso di Lavoro Digitale e Controllo Qualità
%%  Dispensa del Corso di Patologia Post-Laurea
%%  Questo file è incluso con \input{} da main.tex; NON aggiungere il preambolo qui.
%% ============================================================

\chapter{Flusso di Lavoro Digitale e Controllo Qualità}
\label{ch:2}

% ─────────────────────────────────────────────────────────────
\section{Introduzione}
\label{sec:ch2-intro}

La transizione dalla microscopia ottica convenzionale all'imaging di vetrini interi (WSI) ha
modificato in modo sostanziale la pratica della patologia diagnostica. Anziché un singolo
patologo che esamina un vetrino al microscopio da banco, il moderno laboratorio di patologia
digitale genera immagini digitali ad alta risoluzione che possono essere archiviate, trasmesse e
analizzate computazionalmente. Questo cambiamento introduce sia straordinarie opportunità sia
nuove categorie di rischio. Una diagnosi errata derivante da una scansione sfocata, da una
sezione mal colorata o da un campione mal etichettato non è meno grave di una diagnosi errata su
un vetrino tradizionale; nel dominio digitale, tuttavia, tali errori possono propagarsi
silenziosamente attraverso un intero lotto di casi prima di essere rilevati. Il controllo qualità
(CQ) non è quindi un complemento facoltativo della patologia digitale --- è la disciplina
fondamentale che rende l'intera attività clinicamente sicura e scientificamente riproducibile
\cite{bankhead2022}.

Il flusso di lavoro della patologia digitale può essere concettualizzato come una pipeline lineare
con diverse fasi discrete: ricezione del campione e descrizione macroscopica, processazione
tissutale e preparazione del vetrino, scansione, CQ dell'immagine, gestione delle immagini e,
infine, refertazione diagnostica o analisi computazionale. Ogni fase introduce potenziali fonti
di variabilità ed errore. I fattori pre-analitici --- le condizioni in cui il tessuto viene
fissato, processato, sezionato e colorato --- esercitano la maggiore influenza sulla qualità
finale dell'immagine, eppure sono spesso i meno suscettibili di correzione automatizzata. I
fattori post-analitici, come la deriva della calibrazione dello scanner e le differenze di resa
cromatica tra dispositivi, possono essere affrontati computazionalmente, ma solo se vengono prima
identificati attraverso procedure sistematiche di CQ.

Per gli specializzandi in patologia, una comprensione approfondita del flusso di lavoro digitale
è essenziale per due ragioni. In primo luogo, come futuri patologi diagnostici, saranno
responsabili dell'accettazione o del rifiuto delle immagini digitali come idonee alla diagnosi, e
devono essere in grado di articolare le basi tecniche di tale giudizio. In secondo luogo, con il
progredire del campo verso la diagnosi assistita dall'intelligenza artificiale (IA), la qualità
dei dati di addestramento è direttamente determinata dalla qualità dei vetrini e delle scansioni
sottostanti. \emph{Spazzatura in entrata, spazzatura in uscita} non è solo un aforisma
informatico; nella patologia computazionale è un principio di governance clinica. Le sezioni
seguenti ripercorrono il flusso di lavoro digitale dalla ricezione del campione all'immagine
archiviata, con particolare enfasi sui punti di controllo qualità che salvaguardano l'integrità
diagnostica in ogni fase.

L'importanza di un solido sistema di CQ è sottolineata dal crescente corpus di evidenze che
collega la qualità dell'immagine all'accuratezza diagnostica e alle prestazioni dei modelli di
IA. Studi hanno dimostrato che anche una lieve degradazione della messa a fuoco --- impercettibile
all'osservatore umano in un'ispezione superficiale --- può ridurre significativamente la
sensibilità degli algoritmi di apprendimento profondo per il rilevamento e la classificazione dei
tumori \cite{bankhead2022}. Un approccio sistematico alla gestione della qualità, integrato nelle
procedure operative standard del laboratorio e supportato da tecnologie appropriate, è quindi un
prerequisito per qualsiasi servizio di patologia digitale che aspiri all'eccellenza clinica o
scientifica.

% ─────────────────────────────────────────────────────────────
\section{Fase Pre-analitica: Preparazione del Campione e Protocolli di Scansione}
\label{sec:ch2-preanalytical}

\subsection{Fissazione e Processazione Tissutale}

La fissazione adeguata è il singolo fattore determinante più importante della qualità istologica.
Il fissativo standard in anatomia patologica chirurgica è la formalina neutra tamponata al 10\%
(FNT), una soluzione di formaldeide al 3,7--4\% in soluzione salina tamponata con fosfati a
pH~7,0. La formalina reticola le proteine attraverso ponti metilenici, preservando l'architettura
cellulare e prevenendo l'autolisi. Il tempo di fissazione raccomandato per la maggior parte dei
campioni chirurgici è di 6--72 ore; la sotto-fissazione porta a scarso dettaglio nucleare e a
immunoistochimica (IIC) inaffidabile, mentre la sovra-fissazione causa una reticolazione eccessiva
che può mascherare gli epitopi antigenici e ridurre la sensibilità dei saggi molecolari. Per le
biopsie con ago-cannula, il College of American Pathologists (CAP) e il Royal College of
Pathologists (RCPath) raccomandano generalmente un minimo di 6 ore e un massimo di 48 ore.

Dopo la fissazione, il tessuto viene disidratato attraverso una serie graduata di alcoli (70\%,
95\%, 100\%), chiarificato in xilene o in un sostituto dello xilene, e infiltrato con paraffina
fusa sotto vuoto. Questo processo, eseguito in un processore tissutale automatizzato nell'arco di
8--12 ore, rende il tessuto sufficientemente rigido per il sezionamento sottile. La qualità della
processazione dipende in modo critico dalla completezza della disidratazione; l'acqua residua nel
blocco tissutale causa una scarsa infiltrazione della paraffina, con conseguenti sezioni che si
sbriciolano o si lacerano durante la microtomia. L'inclusione --- l'orientamento del tessuto
processato in uno stampo di paraffina fresca --- richiede competenza tecnica, poiché il piano di
sezione deve essere scelto per massimizzare le informazioni diagnostiche. Un orientamento errato
è un errore pre-analitico che non può essere corretto in nessuna fase successiva.

La scelta del programma del processatore tissutale è una variabile di qualità importante. I
programmi di processazione rapida, che utilizzano temperature elevate ed energia a microonde per
accelerare la disidratazione e l'infiltrazione, possono ridurre il tempo di processazione da 12
ore a sole 2 ore. Tuttavia, la processazione rapida può compromettere la morfologia tissutale se
non viene attentamente validata, in particolare per i tessuti adiposi e i campioni di grandi
dimensioni. I laboratori devono validare i propri programmi di processazione per ciascun tipo di
tessuto e monitorare la qualità della processazione attraverso una valutazione regolare della
qualità delle sezioni e delle prestazioni dell'IIC.

\subsection{Sezionamento, Colorazione e Vetrino Coprioggetto}

La microtomia è il processo di taglio di sezioni sottili (tipicamente 3--5~\textmu m) dal blocco
di paraffina mediante un microtomo rotativo. Lo spessore della sezione è una variabile critica:
le sezioni troppo spesse producono nuclei sovrapposti e scarso dettaglio della cromatina, mentre
le sezioni troppo sottili possono non contenere tessuto sufficiente per la diagnosi. Le sezioni
vengono distese su un bagno d'acqua tiepida per eliminare le pieghe, montate su vetrini di vetro
a carica positiva e asciugate in forno a 60\textdegree C per favorire l'adesione. Le macchine
coprivetri automatizzate applicano uno strato sottile di mezzo di montaggio e un vetrino
coprioggetto; la qualità del coprioggetto influisce direttamente sulla qualità della scansione,
poiché le bolle d'aria o lo spessore non uniforme del mezzo di montaggio creano artefatti ottici
che degradano la messa a fuoco.

La colorazione con ematossilina ed eosina (E\&E) rimane il cardine della diagnosi istopatologica.
L'ematossilina, un colorante basico, colora gli acidi nucleici di blu-viola; l'eosina, un
colorante acido, colora le proteine citoplasmatiche e la matrice extracellulare di rosa. Il
preciso equilibrio tra questi due coloranti --- il rapporto E\&E --- varia tra i laboratori e
persino tra lotti di reagenti, producendo la variabilità di colorazione inter-laboratorio che
rappresenta una sfida importante per la patologia computazionale. La colorazione IIC utilizza
anticorpi coniugati a enzimi o fluorofori per rilevare proteine specifiche nelle sezioni tissutali.
I protocolli IIC sono altamente sensibili alle variabili pre-analitiche: il recupero antigenico
(indotto dal calore o enzimatico), la concentrazione dell'anticorpo, il tempo di incubazione e il
sistema di rivelazione influenzano tutti l'intensità finale della colorazione e il fondo. Le
piattaforme di colorazione automatizzata (ad es., Leica BOND, Ventana BenchMark) migliorano la
riproducibilità ma non eliminano la variabilità inter-corsa.

Il coprioggetto è l'ultima fase della preparazione del vetrino ed è spesso sottovalutato come
fonte di problemi di qualità. Il mezzo di montaggio deve essere applicato in quantità sufficiente
a riempire lo spazio tra il tessuto e il coprioggetto senza traboccare, e il coprioggetto deve
essere applicato senza intrappolare bolle d'aria. I coprivetri automatizzati riducono il rischio
di intrappolamento di bolle d'aria ma richiedono manutenzione e calibrazione regolari. L'indice
di rifrazione del mezzo di montaggio deve corrispondere a quello del vetrino coprioggetto
(circa 1,515) per minimizzare le aberrazioni ottiche; l'uso di un mezzo di montaggio errato può
produrre un alone o un effetto shimmer nell'immagine scansionata che degrada la qualità della
messa a fuoco.

\subsection{Preparazione del Vetrino per la Scansione}

Prima che un vetrino venga caricato in uno scanner per vetrini interi, deve soddisfare determinati
criteri fisici. Il vetrino deve essere pulito e privo di polvere, impronte digitali e mezzo di
montaggio residuo sulla superficie del vetro. I codici a barre o i codici QR stampati
sull'etichetta del vetrino devono essere leggibili dal lettore di etichette dello scanner; le
etichette illeggibili sono una causa comune di errori di identificazione del campione. I vetrini
con spessore tissutale eccessivo, bordi del coprioggetto sollevati o coprioggetti rotti possono
inceppare il meccanismo di trasporto del vetrino dello scanner o produrre errori di messa a fuoco.
Un'ispezione visiva pre-scansione da parte del tecnico di laboratorio è quindi un passaggio
essenziale di CQ che deve essere eseguito sistematicamente utilizzando una lista di controllo
standardizzata.

\subsection{Parametri di Scansione}

Gli scanner per vetrini interi digitalizzano i vetrini istologici acquisendo un mosaico di
tessere di immagine sovrapposte che vengono successivamente unite in un'unica immagine di grandi
dimensioni. I parametri di scansione principali sono:

\begin{itemize}
    \item \textbf{Ingrandimento (obiettivo):} La maggior parte degli scanner offre obiettivi
    20$\times$ e 40$\times$, corrispondenti a dimensioni del pixel di circa 0,5~\textmu m e
    0,25~\textmu m rispettivamente. La scelta dell'ingrandimento dipende dal compito diagnostico;
    il 20$\times$ è sufficiente per la maggior parte dell'istologia di routine, mentre il
    40$\times$ è preferito per la citologia e il dettaglio nucleare fine.

    \item \textbf{Qualità della messa a fuoco:} Gli scanner utilizzano algoritmi di messa a fuoco
    automatica per mantenere la messa a fuoco sull'intera superficie del vetrino. La maggior parte
    dei sistemi esegue una mappa di messa a fuoco pre-scansione, campionando i punti di messa a
    fuoco a intervalli regolari sul tessuto, e poi interpola i valori di messa a fuoco tra i punti
    campionati durante la scansione principale. Il tessuto con spessore non uniforme, pieghe o
    bolle d'aria può vanificare l'algoritmo di messa a fuoco automatica, producendo regioni
    sfocate nell'immagine finale.

    \item \textbf{Velocità di scansione:} Velocità di scansione più elevate riducono i tempi di
    elaborazione ma possono compromettere la qualità dell'immagine, in particolare in termini di
    precisione della messa a fuoco e di unione delle tessere. I laboratori devono validare la
    relazione tra velocità di scansione e qualità dell'immagine per il loro specifico modello di
    scanner e i tipi di tessuto.

    \item \textbf{Z-stacking:} Alcuni scanner possono acquisire più piani focali (z-stack) a
    diverse profondità all'interno della sezione tissutale, producendo un dataset tridimensionale
    da cui è possibile calcolare un'immagine a profondità di campo estesa. Lo z-stacking è
    particolarmente utile per i preparati citologici, dove le cellule a diverse profondità
    potrebbero altrimenti risultare sfocate, ma aumenta sostanzialmente il tempo di scansione e
    le dimensioni del file.
\end{itemize}

La calibrazione dello scanner --- inclusi il bilanciamento del bianco, l'uniformità
dell'illuminazione e il profilo colore --- deve essere eseguita regolarmente secondo il programma
del produttore e verificata con vetrini di riferimento. La deriva della calibrazione è una fonte
sottile ma importante di variabilità cromatica inter-scansione che può confondere sia
l'interpretazione umana sia gli algoritmi di IA \cite{bankhead2022}. I laboratori devono
mantenere un registro degli eventi di calibrazione dello scanner e monitorare lo stato della
calibrazione come parte del loro programma di CQ di routine.

% ─────────────────────────────────────────────────────────────
\section{Strategie di Campionamento e Descrizione Macroscopica}
\label{sec:ch2-sampling}

\subsection{Descrizione Macroscopica}

L'esame macroscopico è il primo passaggio analitico eseguito dal patologo o dall'assistente di
patologia alla ricezione di un campione chirurgico. Una descrizione macroscopica sistematica
documenta il tipo di campione, le dimensioni, il peso, le caratteristiche della superficie esterna
e l'aspetto della superficie di taglio. Per i campioni oncologici, la descrizione deve includere
le dimensioni del tumore, la localizzazione, la distanza dai margini chirurgici e la presenza di
linfonodi o altre strutture anatomiche. Queste informazioni non sono meramente amministrative;
guidano direttamente la selezione dei blocchi tissutali per l'esame istologico e determinano se
il campionamento è adeguato per la stadiazione e la valutazione dei margini.

La descrizione macroscopica deve seguire un modulo standardizzato appropriato al tipo di campione.
I moduli standardizzati, come quelli pubblicati dal Royal College of Pathologists o dal College
of American Pathologists, garantiscono che tutte le caratteristiche diagnosticamente rilevanti
siano registrate e che la descrizione sia riproducibile tra osservatori diversi. Nel contesto
della patologia digitale, le immagini macroscopiche --- fotografie della superficie di taglio del
campione --- sono sempre più incorporate nel fascicolo digitale del caso, fornendo un riferimento
visivo che integra le immagini microscopiche e facilita la correlazione tra i reperti macroscopici
e microscopici.

La documentazione dei reperti macroscopici nell'era della patologia digitale va oltre la
descrizione scritta. La fotografia macroscopica, con fotocamere calibrate e carte di riferimento
cromatico, fornisce una documentazione visiva permanente del campione che può essere esaminata
insieme alle immagini microscopiche. Alcuni laboratori stanno iniziando a esplorare la scansione
tridimensionale dei campioni macroscopici mediante luce strutturata o tomografia computerizzata,
che potrebbe fornire un riferimento volumetrico per il campionamento microscopico.

\subsection{Strategie di Campionamento per Diversi Tipi di Campione}

\paragraph{Biopsie con ago-cannula.}
Le biopsie con ago-cannula sono piccoli cilindri tissutali, tipicamente lunghi 1--2~cm e con un
diametro di 1--2~mm, ottenuti mediante biopsia percutanea con ago di organi solidi (mammella,
prostata, rene, fegato). Poiché il volume tissutale è limitato, tutti i cilindri vengono
solitamente inviati per l'esame istologico. I cilindri vengono inclusi in un unico blocco e si
tagliano più livelli (tipicamente 3--5) per massimizzare la probabilità di campionare la lesione.
Nel flusso di lavoro digitale, la piccola area tissutale significa che la scansione a 20$\times$
è solitamente sufficiente e i tempi di scansione sono brevi.

\paragraph{Campioni da escissione e resezione.}
I campioni da resezione (ad es., colectomia, mastectomia, nefrectomia) sono di grandi dimensioni
e richiedono un campionamento selettivo. Il patologo deve identificare le aree diagnosticamente
più informative --- il centro del tumore, il fronte invasivo, i margini chirurgici e le aree
macroscopicamente anomale --- e inviare blocchi rappresentativi. I protocolli di campionamento
sono definiti da linee guida nazionali e internazionali e devono essere documentati nella
descrizione macroscopica. Il sotto-campionamento è una fonte significativa di errore diagnostico;
ad esempio, il mancato campionamento del punto più profondo di un tumore colorettale può portare
a una sottostadiazione della categoria T.

\paragraph{Campioni citologici.}
I preparati citologici --- inclusi gli strisci da agoaspirato con ago sottile (FNAC), i preparati
citologici in fase liquida (LBC) e i campioni di espettorato o urina --- presentano sfide di
campionamento uniche. La cellularità del preparato è molto variabile e le cellule diagnostiche
possono essere scarsamente distribuite. Nel flusso di lavoro digitale, i vetrini citologici
vengono tipicamente scansionati a 40$\times$ per catturare il dettaglio nucleare fine, e l'intera
superficie del vetrino deve essere scansionata anziché una regione selezionata.

\subsection{Come i Reperti Macroscopici Guidano il Campionamento Microscopico}

La relazione tra reperti macroscopici e microscopici è bidirezionale. I reperti macroscopici
guidano la selezione dei blocchi tissutali, come descritto sopra. Al contrario, i reperti
microscopici possono richiedere un campionamento aggiuntivo del campione macroscopico --- ad
esempio, se le sezioni iniziali mostrano tumore a un margine, devono essere inviati blocchi
aggiuntivi da quel margine. Nel flusso di lavoro della patologia digitale, questo processo
iterativo è facilitato dalla possibilità di annotare le immagini digitali e collegare le
annotazioni a blocchi specifici nella descrizione macroscopica, creando un registro spazialmente
coerente dell'intero campione.

% ─────────────────────────────────────────────────────────────
\section{Controllo Qualità: Messa a Fuoco, Colorazione e Artefatti}
\label{sec:ch2-qc}

\subsection{Panoramica del CQ nella Patologia Digitale}

Il controllo qualità nella patologia digitale comprende tutte le procedure progettate per
rilevare, documentare e correggere le carenze nella qualità dell'immagine prima che le immagini
vengano utilizzate per la diagnosi o l'analisi. A differenza della microscopia convenzionale, in
cui il patologo può regolare la messa a fuoco e l'illuminazione in tempo reale, la patologia
digitale richiede che tutti i problemi di qualità vengano identificati e risolti prima che
l'immagine venga rilasciata per l'uso. Un sistema di CQ deve quindi operare a più livelli: nella
fase di preparazione del vetrino (CQ pre-scansione), al momento della scansione (CQ durante la
scansione) e dopo la scansione (CQ post-scansione). Ogni livello affronta diverse categorie di
difetti e richiede competenze e strumenti diversi \cite{janowczyk2019}.

Le conseguenze di un CQ inadeguato sono significative. In un contesto diagnostico, un'immagine
sfocata o piena di artefatti può portare a una diagnosi mancata o errata. In un contesto di
ricerca o di IA, le immagini di scarsa qualità in un dataset di addestramento introducono rumore
che degrada le prestazioni del modello e può introdurre bias sistematici. Ad esempio, se le
immagini sfocate sono più comuni in una categoria diagnostica rispetto a un'altra, un modello
addestrato su tali dati potrebbe imparare ad associare la qualità della messa a fuoco alla
diagnosi piuttosto che alla morfologia tissutale. Un CQ rigoroso è quindi un prerequisito per
un'IA affidabile in patologia \cite{bankhead2022}.

\subsection{Valutazione della Qualità della Messa a Fuoco}

La qualità della messa a fuoco è la metrica di qualità dell'immagine più comunemente valutata
nella patologia digitale. Un'immagine ben a fuoco mostra membrane nucleari nitide, pattern di
cromatina distinti e confini citoplasmatici chiari. Un'immagine sfocata appare indistinta, con
perdita del dettaglio strutturale fine. La qualità della messa a fuoco può essere valutata
visivamente da un osservatore addestrato o computazionalmente utilizzando metriche di nitidezza
dell'immagine. Le metriche computazionali comuni includono la varianza del Laplaciano (una misura
del contenuto ad alta frequenza nell'immagine), il gradiente di Brenner e la funzione Tenengrad.
Queste metriche vengono calcolate sulle tessere dell'immagine e aggregate per produrre un punteggio
di messa a fuoco a livello di vetrino.

La valutazione automatizzata della messa a fuoco è implementata in diversi strumenti di CQ
commerciali e open source. La pipeline HistoQC, sviluppata da Janowczyk e colleghi, calcola una
serie di metriche di qualità dell'immagine --- incluse messa a fuoco, intensità della colorazione,
copertura tissutale e presenza di artefatti --- e genera un rapporto di qualità per ogni vetrino
\cite{janowczyk2019}. HistoQC utilizza una pipeline configurabile di moduli di elaborazione delle
immagini, consentendo ai laboratori di personalizzare i criteri di CQ per i loro specifici tipi
di tessuto e protocolli di colorazione. I vetrini che non superano il CQ vengono segnalati per
la revisione e, se necessario, riscansionati o ricolorati.

\subsection{Uniformità della Colorazione}

L'uniformità della colorazione si riferisce alla coerenza dell'intensità della colorazione
attraverso la sezione tissutale e tra i vetrini all'interno di un lotto. La colorazione non
uniforme può derivare da una distribuzione non uniforme dei reagenti durante la colorazione
automatizzata, dalla variazione dello spessore tissutale o da una deparaffinizzazione incompleta.
Nelle sezioni colorate con E\&E, la colorazione non uniforme si manifesta come aree di tessuto
pallido o sovracolorato che possono essere scambiate per alterazioni patologiche. Nell'IIC, la
colorazione non uniforme può produrre risultati falsi positivi o falsi negativi nelle aree della
sezione che sono sovra- o sotto-colorate.

L'uniformità della colorazione viene valutata misurando la densità ottica o i valori cromatici
delle regioni tissutali sull'intero vetrino e calcolando statistiche riassuntive (media,
deviazione standard, coefficiente di variazione). I vetrini di riferimento con caratteristiche di
colorazione note possono essere inclusi in ogni corsa di colorazione come controlli interni.

\subsection{Artefatti Comuni nella Patologia Digitale}

Gli artefatti sono caratteristiche dell'immagine che non corrispondono a strutture tissutali
genuine e che possono interferire con la diagnosi o l'analisi. Gli artefatti più comuni nella
patologia digitale sono:

\begin{itemize}
    \item \textbf{Bolle d'aria:} L'aria intrappolata sotto il coprioggetto appare come aree
    circolari o irregolari chiare che oscurano il tessuto sottostante. Le bolle d'aria sono
    causate da un'applicazione inadeguata del mezzo di montaggio o dalla presenza di aria nel
    mezzo di montaggio.

    \item \textbf{Pieghe tissutali:} Le pieghe si verificano quando la sezione non viene
    completamente distesa sul bagno d'acqua prima del montaggio. Appaiono come bande scure e
    ispessite nell'immagine e producono regioni sfocate perché la superficie tissutale non è
    complanare con il piano focale dello scanner.

    \item \textbf{Polvere e detriti:} Le particelle di polvere sulla superficie del vetrino o
    sull'ottica dello scanner appaiono come macchie o striature scure nell'immagine.

    \item \textbf{Regioni sfocate:} Come descritto sopra, le regioni sfocate derivano da un
    errore di messa a fuoco automatica, tipicamente causato da pieghe tissutali, bolle d'aria o
    coprioggetto non uniforme.

    \item \textbf{Problemi con il coprioggetto:} Un coprioggetto sollevato o rotto impedisce allo
    scanner di raggiungere la messa a fuoco sull'intero vetrino.

    \item \textbf{Segni di penna e inchiostro:} L'inchiostrazione dei margini chirurgici e i
    segni di penna sull'etichetta del vetrino possono trasferirsi sulla superficie tissutale,
    producendo artefatti colorati che possono essere erroneamente identificati come colorazione
    tissutale dagli algoritmi automatizzati.
\end{itemize}

\subsection{Strumenti di CQ Automatizzati: HistoQC}

HistoQC è una pipeline di controllo qualità open source e configurabile per le immagini di
vetrini interi, descritta da Janowczyk et al.\ \cite{janowczyk2019}. La pipeline accetta file WSI
in formati standard (SVS, NDPI, MRXS, ecc.) e applica una serie di moduli di analisi delle
immagini per generare un rapporto di qualità per vetrino. I moduli includono la segmentazione
tissutale (per distinguere il tessuto dallo sfondo), la valutazione della messa a fuoco (mediante
varianza del Laplaciano), la misurazione dell'intensità della colorazione, il rilevamento di
segni di penna e artefatti da inchiostro, e l'identificazione di regioni sfocate. L'output è un
rapporto strutturato in formato CSV, insieme a immagini miniatura annotate che evidenziano le
aree problematiche.

HistoQC può essere integrato nel sistema informativo di laboratorio (SIL) come parte di un flusso
di lavoro di CQ post-scansione automatizzato. I vetrini che non soddisfano uno o più criteri di
CQ vengono automaticamente segnalati e indirizzati a un tecnico o a un patologo per la revisione.
La pipeline è stata validata su grandi coorti di vetrini del The Cancer Genome Atlas (TCGA) e ha
dimostrato di identificare problemi di qualità associati a errori nelle analisi successive
\cite{janowczyk2019}.

\subsection{Documentazione e Gestione delle Scansioni Non Conformi}

Una scansione non conforme è qualsiasi scansione che non soddisfa i criteri di CQ definiti dal
laboratorio. Le scansioni non conformi devono essere documentate nel SIL con un registro del
motivo del fallimento, delle metriche di CQ non soddisfatte e dell'azione intrapresa (nuova
scansione, nuova colorazione, nuovo montaggio o accettazione con annotazione). Questa
documentazione è essenziale per scopi di audit e per identificare problemi di qualità sistematici.

% ─────────────────────────────────────────────────────────────
%  Figura del Flusso di Lavoro CQ
% ─────────────────────────────────────────────────────────────
\begin{figure}[htbp]
    \centering
    \usetikzlibrary{arrows.meta, shapes.geometric, positioning, calc}

\begin{tikzpicture}[
    node distance = 0.6cm,
    every node/.style = {font=\small\sffamily},
    proc/.style = {
        rectangle, rounded corners=4pt,
        fill={rgb,255:red,46;green,134;blue,171},
        text=white, minimum width=3.4cm, minimum height=0.72cm,
        align=center, font=\small\sffamily\bfseries
    },
    decision/.style = {
        diamond, aspect=2.6,
        fill={rgb,255:red,255;green,165;blue,0},
        text=white, minimum width=3.4cm, minimum height=0.9cm,
        align=center, font=\small\sffamily\bfseries,
        inner sep=1pt
    },
    pass/.style = {
        rectangle, rounded corners=4pt,
        fill={rgb,255:red,56;green,161;blue,105},
        text=white, minimum width=3.0cm, minimum height=0.72cm,
        align=center, font=\small\sffamily\bfseries
    },
    fail/.style = {
        rectangle, rounded corners=4pt,
        fill={rgb,255:red,229;green,62;blue,62},
        text=white, minimum width=3.0cm, minimum height=0.72cm,
        align=center, font=\small\sffamily\bfseries
    },
    arr/.style = {
        -{Stealth[length=5pt,width=4pt]},
        line width=0.8pt, color=gray!70
    },
    lbl/.style = {font=\footnotesize\sffamily, text=gray!80}
]

%% Main vertical chain
\node[proc]     (recv)     {Slide Received};
\node[proc,     below=0.6cm of recv]     (prescan)  {Pre-scan QC Check};
\node[proc,     below=0.6cm of prescan]  (scan)     {Scanning};
\node[proc,     below=0.6cm of scan]     (postscan) {Post-scan QC Assessment};
\node[decision, below=0.7cm of postscan] (decide)   {Pass QC?};

%% Outcome nodes
\node[pass, below left=0.9cm and 1.5cm of decide]  (archive) {Archive \& Analyse};
\node[fail, below right=0.9cm and 1.5cm of decide] (rescan)  {Re-scan or Reject};

%% Arrows — main chain
\draw[arr] (recv)     -- (prescan);
\draw[arr] (prescan)  -- (scan);
\draw[arr] (scan)     -- (postscan);
\draw[arr] (postscan) -- (decide);

%% Arrows — decision branches
\draw[arr] (decide.south west) -- (archive.north)
    node[lbl, midway, left=2pt] {Yes};
\draw[arr] (decide.south east) -- (rescan.north)
    node[lbl, midway, right=2pt] {No};

%% Loop-back dashed arrow for Re-scan
\draw[arr, dashed, color=gray!50]
    (rescan.east) -- ++(0.6,0) |- (scan.east);

\end{tikzpicture}

    \caption{Panoramica del flusso di lavoro per il controllo qualità nella patologia digitale.
    I vetrini progrediscono dall'ispezione pre-scansione attraverso la scansione e il CQ
    post-scansione automatizzato (HistoQC). I vetrini che superano tutti i criteri di CQ vengono
    rilasciati per la refertazione diagnostica o l'analisi computazionale. I vetrini che non
    superano il CQ vengono indirizzati attraverso un percorso di rimedio (nuova scansione, nuova
    colorazione o nuovo montaggio) prima di rientrare nella pipeline di CQ. Le scansioni non
    conformi che non possono essere remediate vengono documentate e segnalate al patologo
    responsabile.}
    \label{fig:qc_workflow}
\end{figure}

% ─────────────────────────────────────────────────────────────
\section{Normalizzazione della Colorazione e Standardizzazione Cromatica}
\label{sec:ch2-normalisation}

\subsection{Fonti di Variabilità della Colorazione}

La variabilità della colorazione è una delle sfide più pervasive nella patologia digitale, che
influenza sia l'interpretazione umana sia l'analisi computazionale. Anche all'interno di un
singolo laboratorio, l'intensità della colorazione e il bilanciamento cromatico possono variare
tra lotti di reagenti, tra corse di colorazione e nel tempo con l'invecchiamento dei reagenti.
Tra laboratori diversi, la variabilità è amplificata dalle differenze nei protocolli di
colorazione, nei fornitori di reagenti, nelle piattaforme di colorazione automatizzata e nelle
pratiche locali. Quando i vetrini vengono scansionati su modelli di scanner diversi o persino su
unità diverse dello stesso modello, viene introdotta un'ulteriore variabilità cromatica dovuta
alle differenze negli spettri di illuminazione, nei sensori della fotocamera e nella calibrazione
cromatica. L'effetto cumulativo è che lo stesso tessuto, colorato e scansionato in due laboratori
diversi, può produrre immagini con caratteristiche cromatiche sostanzialmente diverse, anche se
la biologia sottostante è identica \cite{tellez2019}.

Questa variabilità ha importanti conseguenze pratiche. Per gli osservatori umani, significa che i
patologi devono ricalibrare mentalmente la loro percezione cromatica quando esaminano vetrini di
laboratori non familiari. Per gli algoritmi di IA, significa che un modello addestrato su vetrini
di un laboratorio può avere prestazioni scarse quando applicato a vetrini di un altro, un
fenomeno noto come \emph{domain shift}. La normalizzazione della colorazione --- la trasformazione
computazionale dei colori del vetrino verso uno standard di riferimento comune --- è la strategia
principale per affrontare questo problema \cite{tellez2019}.

\subsection{Deconvoluzione Cromatica: Il Metodo di Ruifrok e Johnston}

La deconvoluzione cromatica è una tecnica matematica per separare i contributi dei singoli
coloranti al colore di ogni pixel in un'immagine E\&E. Il metodo, descritto da Ruifrok e Johnston
\cite{ruifrok2001}, si basa sulla legge di Beer-Lambert dell'assorbimento ottico. Ogni colorante
assorbe la luce a lunghezze d'onda caratteristiche, e il colore osservato di un pixel è il
risultato dell'assorbimento combinato di tutti i coloranti presenti. Se gli spettri di
assorbimento (vettori di colorante) dei singoli coloranti sono noti, il contributo di ciascun
colorante a ogni pixel può essere calcolato risolvendo un sistema di equazioni lineari.

In pratica, i vettori di colorante per l'ematossilina e l'eosina vengono determinati
empiricamente da immagini di riferimento o dalla letteratura. L'algoritmo di deconvoluzione
cromatica trasforma l'immagine RGB in una rappresentazione separata per colorante, producendo
immagini separate per il canale dell'ematossilina (colorazione nucleare) e per il canale
dell'eosina (colorazione citoplasmatica). Il metodo di Ruifrok e Johnston è stato implementato in
piattaforme di analisi delle immagini ampiamente utilizzate, tra cui ImageJ/Fiji e QuPath
\cite{bankhead2017}.

La formulazione matematica della deconvoluzione cromatica è la seguente. Sia $\mathbf{I}$
l'immagine RGB e $\mathbf{M}$ la matrice di colorante $3 \times 3$ le cui colonne sono i vettori
di colorante (spettri di densità ottica) dei singoli coloranti. L'immagine di densità ottica
$\mathbf{OD} = -\log_{10}(\mathbf{I}/\mathbf{I}_0)$, dove $\mathbf{I}_0$ è l'intensità dello
sfondo, può essere decomposta come $\mathbf{OD} = \mathbf{M} \cdot \mathbf{C}$, dove
$\mathbf{C}$ è la matrice delle concentrazioni di colorante. Le concentrazioni di colorante
vengono recuperate con $\mathbf{C} = \mathbf{M}^{-1} \cdot \mathbf{OD}$ \cite{ruifrok2001}.

\subsection{Algoritmi di Normalizzazione della Colorazione}

Gli algoritmi di normalizzazione della colorazione trasformano la distribuzione cromatica di
un'immagine sorgente per farla corrispondere a quella di un'immagine di riferimento (target),
riducendo così la variabilità cromatica inter-vetrino. I metodi classici più utilizzati sono:

\begin{itemize}
    \item \textbf{Normalizzazione di Macenko:} Questo metodo stima la matrice di colorante
    dell'immagine sorgente mediante decomposizione ai valori singolari (SVD) della matrice di
    densità ottica, quindi riscala le concentrazioni di colorante per farle corrispondere a quelle
    dell'immagine di riferimento.

    \item \textbf{Normalizzazione di Vahadane:} Questo metodo utilizza la fattorizzazione a
    matrice non negativa sparsa (SNMF) per stimare la matrice di colorante, più robusta agli
    outlier e più accurata nella separazione dei coloranti rispetto ai metodi basati su SVD.

    \item \textbf{Normalizzazione di Reinhard:} Questo metodo trasferisce le statistiche
    cromatiche (media e deviazione standard) dell'immagine di riferimento all'immagine sorgente
    nello spazio colore LAB. È il metodo di normalizzazione più semplice e veloce.
\end{itemize}

Gli approcci di apprendimento profondo alla normalizzazione della colorazione, incluse le reti
generative avversariali (GAN) e le CycleGAN, hanno dimostrato di produrre immagini normalizzate
visivamente più realistiche rispetto ai metodi classici, ma richiedono grandi dataset di
addestramento e sono computazionalmente costosi \cite{tellez2019}.

\subsection{Importanza per le Applicazioni di IA}

L'importanza della normalizzazione della colorazione per le applicazioni di IA non può essere
sopravvalutata. I modelli di apprendimento profondo per la patologia computazionale vengono
addestrati su grandi dataset di WSI annotate, e le loro prestazioni sono altamente sensibili alle
caratteristiche cromatiche dei dati di addestramento. La normalizzazione della colorazione,
applicata in modo coerente sia ai dati di addestramento sia a quelli di test, riduce questo
\emph{domain shift} indotto dal colore e migliora la generalizzabilità dei modelli di IA tra
laboratori e scanner \cite{tellez2019}.

% ─────────────────────────────────────────────────────────────
\section{Collaborazione tra Tecnico e Patologo}
\label{sec:ch2-collaboration}

\subsection{Ruoli Complementari}

Il laboratorio di patologia digitale è un ambiente collaborativo in cui i tecnici di laboratorio
e i patologi svolgono ruoli complementari e interdipendenti. I tecnici sono responsabili della
preparazione fisica dei vetrini --- fissazione, processazione, inclusione, sezionamento,
colorazione e coprioggetto --- e del funzionamento delle apparecchiature di scansione. Sono la
prima linea del controllo qualità, eseguendo l'ispezione visiva dei vetrini prima della
scansione e monitorando la qualità della scansione in tempo reale. I patologi sono responsabili
dell'interpretazione diagnostica delle immagini digitali e dell'accettazione o del rifiuto finale
delle immagini come idonee alla diagnosi.

La transizione alla patologia digitale ha, per certi versi, aumentato l'importanza del ruolo del
tecnico. Nella patologia tradizionale su vetrino, il patologo può compensare le carenze minori
nella preparazione del vetrino regolando la messa a fuoco o l'illuminazione del microscopio.
Nella patologia digitale, il patologo dipende interamente dalla qualità dell'immagine scansionata;
non c'è possibilità di riesaminare il vetrino fisico durante la sessione diagnostica. Ciò pone
un onere maggiore sul tecnico per garantire che i vetrini siano della massima qualità possibile
prima della scansione.

\subsection{Flusso di Comunicazione e Segnalazione dei Problemi di Qualità}

Il flusso di comunicazione in un laboratorio di patologia digitale segue tipicamente un percorso
strutturato. Quando un tecnico identifica un problema di qualità --- ad esempio, un vetrino con
una piega tissutale, una bolla d'aria o un artefatto di colorazione --- segnala il problema nel
SIL e, se necessario, contatta direttamente il patologo responsabile. La segnalazione deve
includere una descrizione del problema, i vetrini interessati e la valutazione del tecnico sulla
probabilità che il problema influenzi l'interpretazione diagnostica.

Gli strumenti di CQ automatizzati come HistoQC \cite{janowczyk2019} possono pre-filtrare i
vetrini e indirizzare solo quelli con problemi di qualità significativi alla revisione umana,
riducendo il carico sia sui tecnici sia sui patologi. Tuttavia, gli strumenti automatizzati non
possono sostituire il giudizio umano per i casi limite, e la decisione finale sull'idoneità
diagnostica di un vetrino deve sempre spettare al patologo responsabile.

\subsection{Procedure di Escalation}

Le procedure di escalation definiscono il percorso per la gestione dei problemi di qualità che
non possono essere risolti a livello del tecnico. Un tipico percorso di escalation potrebbe
essere: (1) il tecnico identifica il problema di qualità e tenta la rimediazione (ad es., nuova
scansione, nuovo montaggio); (2) se la rimediazione non ha successo, il tecnico segnala il
problema al biologo senior o al responsabile del laboratorio; (3) se il problema non può essere
risolto a livello del biologo senior, viene escalato al patologo responsabile; (4) se il patologo
determina che il problema di qualità influisce sull'interpretazione diagnostica e non può essere
risolto, il caso viene escalato al team clinico con un referto che indica la limitazione di
qualità.

\subsection{Standard di Documentazione}

La documentazione è la spina dorsale della gestione della qualità nel laboratorio di patologia
digitale. Ogni azione correlata alla qualità --- dall'ispezione visiva pre-scansione alla
decisione del patologo di accettare un'immagine non ottimale --- deve essere registrata in un
formato standardizzato recuperabile per scopi di audit. Gli standard di documentazione devono
specificare: le informazioni da registrare (identificatore del vetrino, data e ora, natura del
problema, azione intrapresa, esito); il formato del registro (inserimento dati strutturato nel
SIL, commento in testo libero o entrambi); e il periodo di conservazione dei registri.

% ─────────────────────────────────────────────────────────────
\section{Documentazione e Tracciabilità}
\label{sec:ch2-documentation}

\subsection{Catena di Custodia}

La catena di custodia si riferisce al registro documentato del possesso, della gestione e della
localizzazione di un campione dal momento della raccolta a quello dello smaltimento. In patologia,
il mantenimento di una catena di custodia ininterrotta è un requisito legale ed etico, poiché
garantisce che il campione esaminato dal patologo sia lo stesso campione prelevato dal paziente.
Nel contesto della patologia digitale, la catena di custodia si estende all'immagine digitale:
l'immagine deve essere dimostrabilmente collegata al vetrino fisico, che deve essere
dimostrabilmente collegato al campione del paziente.

La catena di custodia è mantenuta attraverso una combinazione di etichettatura fisica (codici a
barre, codici QR o tag RFID su contenitori di campioni, cassette e vetrini) e registrazioni
elettroniche nel SIL. Ogni trasferimento di custodia deve essere registrato con un timestamp e
l'identità della persona responsabile. Le conseguenze di un'interruzione della catena di custodia
possono essere gravi: un campione etichettato in modo errato o confuso con un altro campione può
portare a una diagnosi errata, con conseguenze potenzialmente catastrofiche per il paziente.

\subsection{Codici a Barre e Identificazione Digitale}

Il codice a barre è la tecnologia principale per l'identificazione dei campioni nel moderno
laboratorio di patologia. I codici a barre vengono stampati su contenitori di campioni, cassette
e vetrini al momento dell'etichettatura, e vengono letti in ogni fase del flusso di lavoro per
verificare l'identità del campione e registrare il trasferimento di custodia. I codici a barre
bidimensionali (codici QR, codici DataMatrix) sono preferiti a quelli unidimensionali perché
possono codificare più informazioni in uno spazio più piccolo e sono più resistenti ai danni.

Nel flusso di lavoro della patologia digitale, il codice a barre del vetrino viene letto dallo
scanner al momento del caricamento, e l'immagine digitale risultante viene automaticamente
collegata al corrispondente registro del caso nel SIL. Questo collegamento automatizzato elimina
il rischio di errori di trascrizione manuale.

\subsection{Audit Trail Digitali}

Un audit trail digitale è un registro cronologico di tutte le azioni eseguite su un oggetto
digitale --- in questo contesto, un file WSI. L'audit trail registra chi ha avuto accesso
all'immagine, quando, da quale dispositivo e quale azione è stata intrapresa (visualizzazione,
annotazione, refertazione, esportazione, eliminazione). Gli audit trail digitali sono essenziali
per la conformità normativa, per scopi medicolegali e per l'indagine di incidenti di qualità.

\subsection{Requisiti Normativi}

I laboratori di patologia digitale sono soggetti a una serie di requisiti normativi e di
accreditamento che disciplinano la qualità dei loro processi e della loro documentazione. Gli
standard più rilevanti sono:

\begin{itemize}
    \item \textbf{ISO 15189:} Lo standard internazionale per i laboratori medici, ISO 15189
    specifica i requisiti per la qualità e la competenza nei test di laboratorio medico. Richiede
    ai laboratori di stabilire e mantenere un sistema di gestione della qualità, di validare i
    propri metodi analitici, di partecipare a programmi di valutazione esterna della qualità (VEQ)
    e di mantenere una documentazione completa di tutte le attività correlate alla qualità.

    \item \textbf{Accreditamento CAP:} Il College of American Pathologists (CAP) Laboratory
    Accreditation Program fornisce l'accreditamento per i laboratori di patologia negli Stati
    Uniti e a livello internazionale. Il CAP ha pubblicato liste di controllo specifiche per la
    patologia digitale, che coprono la validazione dello scanner, la gestione delle immagini, le
    procedure di CQ e la refertazione diagnostica.

    \item \textbf{IVDR (UE):} Nell'Unione Europea, i sistemi di patologia digitale utilizzati per
    la diagnosi primaria sono regolamentati come dispositivi medici diagnostici in vitro (IVD) ai
    sensi del Regolamento sui Dispositivi Medico-Diagnostici In Vitro (IVDR, UE 2017/746).
\end{itemize}

% ─────────────────────────────────────────────────────────────
\section{Attività Pratiche}
\label{sec:ch2-practical}

\subsection{Attività A: Esercizio di Riconoscimento degli Artefatti}

\textbf{Obiettivo:} Sviluppare la capacità di identificare e classificare gli artefatti comuni
nelle immagini di patologia digitale e valutare il loro impatto sulla qualità diagnostica.

\textbf{Materiali:} Un set curato di 20 immagini di vetrini interi (o tessere di immagine
rappresentative) tratte da un dataset pubblicamente disponibile (ad es., TCGA) o dall'archivio
del laboratorio, inclusi esempi dei seguenti artefatti: bolle d'aria, pieghe tissutali, regioni
sfocate, polvere e detriti, segni di penna e difetti del coprioggetto. Deve essere fornita una
guida di riferimento per la classificazione degli artefatti.

\textbf{Procedura:}
\begin{enumerate}
    \item Gli studenti esaminano ogni immagine in modo indipendente e registrano la loro
    valutazione di: (a) il tipo/i di artefatto presente; (b) la proporzione dell'area tissutale
    interessata; (c) se l'artefatto è probabile che influenzi l'interpretazione diagnostica; e
    (d) l'azione di rimediazione raccomandata (nuova scansione, nuova colorazione, nuovo montaggio
    o accettazione con annotazione).

    \item Gli studenti confrontano le loro valutazioni con quelle di un patologo di riferimento
    (fornite come chiave delle risposte) e calcolano il loro tasso di accordo per ogni categoria
    di artefatto.

    \item In piccoli gruppi, gli studenti discutono i casi in cui la loro valutazione differiva
    da quella di riferimento e identificano le ragioni del disaccordo.

    \item Gli studenti utilizzano la pipeline HistoQC \cite{janowczyk2019} per eseguire il CQ
    automatizzato su un sottoinsieme delle immagini e confrontano i risultati del CQ automatizzato
    con le loro valutazioni manuali.
\end{enumerate}

\textbf{Risultati di apprendimento:} Al termine di questa attività, gli studenti dovrebbero
essere in grado di identificare e classificare i principali tipi di artefatti nelle immagini di
patologia digitale, valutarne la rilevanza diagnostica e applicare strumenti di CQ automatizzati
a supporto della valutazione della qualità.

\subsection{Attività B: Studio di Caso --- Flusso di Lavoro Digitale dalla Biopsia alla Diagnosi}

\textbf{Obiettivo:} Tracciare il flusso di lavoro completo della patologia digitale per un caso
clinico, dalla ricezione del campione alla refertazione diagnostica, identificando i punti di
controllo qualità e i potenziali punti di errore in ogni fase.

\textbf{Materiali:} Un caso clinico anonimizzato (fornito dal coordinatore del corso) che include:
il modulo di richiesta, la descrizione macroscopica, le immagini WSI delle sezioni istologiche
(E\&E e IIC), i risultati del CQ automatizzato da HistoQC e il referto diagnostico finale.

\textbf{Procedura:}
\begin{enumerate}
    \item Gli studenti esaminano il caso clinico e costruiscono una mappa del flusso di lavoro
    che tracci il percorso del campione dalla ricezione alla diagnosi, identificando ogni fase,
    il personale responsabile e i criteri di qualità applicabili.

    \item Gli studenti esaminano i risultati del CQ di HistoQC e identificano eventuali problemi
    di qualità nelle immagini WSI. Per ogni problema identificato, propongono un'azione di
    rimediazione e valutano l'impatto potenziale sulla diagnosi.

    \item Gli studenti esaminano le immagini WSI in QuPath e confrontano le loro osservazioni
    diagnostiche con il referto finale. Identificano le caratteristiche morfologiche chiave che
    supportano la diagnosi e valutano se eventuali problemi di qualità dell'immagine hanno
    influenzato la loro capacità di identificare queste caratteristiche.

    \item Gli studenti redigono un breve rapporto (circa 500 parole) che descrive il flusso di
    lavoro del caso, i problemi di qualità identificati, le azioni di rimediazione proposte e
    le lezioni apprese per la pratica futura.
\end{enumerate}

\textbf{Risultati di apprendimento:} Al termine di questa attività, gli studenti dovrebbero
essere in grado di tracciare il flusso di lavoro completo della patologia digitale per un caso
clinico, identificare i punti di controllo qualità e i potenziali punti di errore, e applicare
strumenti di CQ automatizzati e manuali per valutare la qualità dell'immagine.

% ─────────────────────────────────────────────────────────────
\section{Riepilogo}
\label{sec:ch2-summary}

Questo capitolo ha esaminato il flusso di lavoro della patologia digitale e il controllo qualità,
dalla preparazione del campione all'immagine archiviata. I seguenti punti chiave devono essere
ricordati:

\begin{itemize}
    \item Il flusso di lavoro della patologia digitale è una pipeline lineare con diverse fasi
    discrete, ognuna delle quali introduce potenziali fonti di variabilità ed errore. Il controllo
    qualità deve operare a ogni fase per garantire l'integrità diagnostica.

    \item I fattori pre-analitici --- fissazione, processazione, sezionamento, colorazione e
    coprioggetto --- esercitano la maggiore influenza sulla qualità finale dell'immagine e non
    possono essere corretti computazionalmente dopo la scansione.

    \item La valutazione della qualità della messa a fuoco, dell'uniformità della colorazione e
    degli artefatti è essenziale per garantire che le immagini siano idonee alla diagnosi e
    all'analisi computazionale. Strumenti automatizzati come HistoQC \cite{janowczyk2019}
    supportano ma non sostituiscono il giudizio umano.

    \item La normalizzazione della colorazione --- mediante metodi come Macenko, Vahadane o
    Reinhard --- riduce la variabilità cromatica inter-laboratorio e migliora la
    generalizzabilità dei modelli di IA \cite{tellez2019}. La deconvoluzione cromatica secondo
    Ruifrok e Johnston \cite{ruifrok2001} consente l'analisi quantitativa dei singoli coloranti.

    \item La collaborazione tra tecnico e patologo è fondamentale per un servizio di patologia
    digitale di alta qualità. I tecnici sono la prima linea del controllo qualità; i patologi
    sono responsabili dell'accettazione finale delle immagini come idonee alla diagnosi.

    \item La documentazione e la tracciabilità --- incluse la catena di custodia, i codici a
    barre, gli audit trail digitali e la conformità agli standard normativi (ISO 15189, CAP,
    IVDR) --- sono requisiti non negoziabili per un servizio di patologia digitale clinicamente
    sicuro e legalmente conforme.
\end{itemize}

% ─────────────────────────────────────────────────────────────
\section{Domande di Revisione}
\label{sec:ch2-questions}

\begin{enumerate}

    \item \label{ch2:q1} Quale fissativo è più comunemente utilizzato per la processazione
    tissutale istologica di routine nella patologia digitale?
    \begin{enumerate}[label=\Alph*.]
        \item Alcool etilico al 70\%
        \item Glutaraldeide al 2,5\%
        \item Formalina neutra tamponata al 10\%
        \item Soluzione di Bouin
    \end{enumerate}

    \item \label{ch2:q2} Qual è lo scopo principale della normalizzazione della colorazione nella
    patologia digitale?
    \begin{enumerate}[label=\Alph*.]
        \item Aumentare la risoluzione delle immagini WSI
        \item Ridurre la variazione nell'aspetto della colorazione causata da laboratori,
        protocolli e scanner diversi
        \item Eliminare la necessità di colorazione IIC
        \item Convertire le immagini a colori in immagini in scala di grigi per l'analisi
    \end{enumerate}

    \item \label{ch2:q3} Quale dei seguenti artefatti è più probabile che causi una regione di
    un'immagine di vetrino intero che appare sfocata?
    \begin{enumerate}[label=\Alph*.]
        \item Sovra-colorazione con ematossilina
        \item Contaminazione batterica del campione
        \item Spessore tissutale non uniforme
        \item Uso di un obiettivo a basso ingrandimento
    \end{enumerate}

    \item \label{ch2:q4} Il metodo di deconvoluzione cromatica sviluppato da Ruifrok e Johnston
    viene utilizzato per:
    \begin{enumerate}[label=\Alph*.]
        \item Aumentare la risoluzione delle immagini WSI mediante interpolazione
        \item Separare i contributi dei singoli coloranti (ad es., ematossilina e DAB) da
        un'immagine RGB
        \item Convertire le immagini WSI dal formato SVS al formato OME-TIFF
        \item Rilevare automaticamente le regioni tumorali nelle immagini WSI
    \end{enumerate}

    \item \label{ch2:q5} Nel contesto del controllo qualità nella patologia digitale, HistoQC è:
    \begin{enumerate}[label=\Alph*.]
        \item Un sistema commerciale di gestione delle immagini per la patologia digitale
        \item Uno strumento open source per la valutazione automatizzata della qualità delle
        immagini di vetrini interi
        \item Un protocollo di colorazione standardizzato per l'E\&E
        \item Un formato di file per l'archiviazione di immagini WSI
    \end{enumerate}

    \item \label{ch2:q6} Quale dei seguenti descrive meglio la fase pre-analitica nella patologia
    digitale?
    \begin{enumerate}[label=\Alph*.]
        \item Solo la fase di scansione del vetrino
        \item Solo la fase di refertazione diagnostica
        \item Tutte le fasi dalla ricezione del campione alla scansione del vetrino
        \item Solo la fase di analisi computazionale dell'immagine
    \end{enumerate}

    \item \label{ch2:q7} Qual è il motivo principale per cui il tempo di fissazione deve essere
    attentamente controllato nella patologia digitale?
    \begin{enumerate}[label=\Alph*.]
        \item Il tempo di fissazione influenza la dimensione del file WSI
        \item Il tempo di fissazione influenza la morfologia tissutale, la preservazione degli
        antigeni e le analisi IIC e molecolari successive
        \item Il tempo di fissazione determina il numero di livelli della piramide di immagine
        \item Il tempo di fissazione influenza la velocità di scansione dello scanner
    \end{enumerate}

    \item \label{ch2:q8} Quale dei seguenti è il primo passo più appropriato quando si scopre che
    una WSI scansionata presenta significative regioni sfocate?
    \begin{enumerate}[label=\Alph*.]
        \item Accettare l'immagine con un'annotazione di qualità e procedere con la diagnosi
        \item Eliminare l'immagine e richiedere un nuovo campione al chirurgo
        \item Riscansionare il vetrino dopo aver verificato la calibrazione della messa a fuoco
        dello scanner
        \item Applicare un filtro di nitidezza computazionale all'immagine
    \end{enumerate}

    \item \label{ch2:q9} In un laboratorio di patologia digitale, la documentazione della catena
    di custodia serve a:
    \begin{enumerate}[label=\Alph*.]
        \item Ridurre il tempo di scansione dei vetrini
        \item Garantire la tracciabilità dei campioni e dei vetrini attraverso il flusso di lavoro
        \item Migliorare la qualità della colorazione E\&E
        \item Aumentare la velocità di refertazione diagnostica
    \end{enumerate}

    \item \label{ch2:q10} Quale standard internazionale è più rilevante per la gestione della
    qualità in un laboratorio medico che esegue patologia digitale?
    \begin{enumerate}[label=\Alph*.]
        \item ISO 9001
        \item ISO 15189
        \item ISO 27001
        \item ISO 13485
    \end{enumerate}

\end{enumerate}

% ============================================================
%  Capitolo 3 – Fondamenti dell'Intelligenza Artificiale in Patologia
%  Dispensa del Corso di Patologia Post-Laurea
%  Questo file è incluso con \input{} dal documento principale; NON aggiungere
%  \documentclass, \begin{document}, o \end{document} qui.
% ============================================================

\chapter{Fondamenti dell'Intelligenza Artificiale in Patologia}
\label{ch:3}

% ============================================================
\section{Introduzione}
\label{sec:ch3_intro}
% ============================================================

La patologia è sempre stata una disciplina fondata sul riconoscimento di pattern. Dalle prime descrizioni microscopiche dell'architettura cellulare all'analisi contemporanea di immagini di vetrini interi (WSI) che si estendono su miliardi di pixel, il compito fondamentale del patologo è estrarre un segnale significativo da dati visivi e molecolari complessi. L'intelligenza artificiale (IA) sta ora ridefinendo le modalità con cui questo compito viene svolto, offrendo strumenti computazionali in grado di elaborare immagini a una scala e a una velocità ben oltre la capacità umana, di identificare sottili caratteristiche morfologiche che possono sfuggire all'occhio non assistito, e di integrare flussi di dati multimodali — istologia, genomica, proteomica e metadati clinici — in modelli predittivi unificati \cite{bera2019}. Le implicazioni per l'accuratezza diagnostica, l'efficienza del flusso di lavoro e, in ultima analisi, gli esiti dei pazienti sono profonde.

La storia dell'IA in patologia è più lunga di quanto spesso si riconosca. I primi tentativi di automatizzare il conteggio cellulare e l'analisi morfometrica risalgono agli anni Sessanta e Settanta, quando analizzatori di immagini analogici venivano utilizzati per quantificare l'area nucleare e l'intensità della colorazione nei preparati citologici. L'avvento della patologia digitale negli anni Novanta e Duemila — guidato da scanner per vetrini interi in grado di digitalizzare vetrini di vetro ad alta risoluzione — ha creato il substrato di dati su cui si basano i moderni metodi computazionali. I sistemi di analisi delle immagini basati su regole, in grado di rilevare nuclei positivamente colorati o di segmentare compartimenti tissutali secondo criteri espliciti di colore e forma, sono diventati strumenti di laboratorio di routine negli anni Duemila \cite{bankhead2017}. Il cambiamento trasformativo dell'ultimo decennio è stato l'ascesa dell'apprendimento automatico (ML) e, in particolare, dell'apprendimento profondo (DL), che ha permesso ai sistemi di apprendere rappresentazioni complesse e gerarchiche delle immagini patologiche direttamente da dati annotati, senza la necessità di regole elaborate manualmente \cite{komura2018}.

Il ritmo del progresso è stato notevole. Studi fondamentali pubblicati tra il 2016 e il 2023 hanno dimostrato che i modelli di apprendimento profondo potevano eguagliare o superare le prestazioni diagnostiche di patologi certificati in compiti che spaziavano dal rilevamento di metastasi linfonodali nel cancro al seno alla classificazione dell'adenocarcinoma prostatico e alla predizione di sottotipi molecolari da sole sezioni colorate con ematossilina ed eosina (E\&E) \cite{song2023}. Questi risultati hanno catalizzato investimenti sostanziali sia dal settore pubblico sia dall'industria delle tecnologie mediche, e un numero crescente di strumenti diagnostici basati sull'IA ha ricevuto l'autorizzazione normativa da organismi come la Food and Drug Administration (FDA) degli Stati Uniti e l'Agenzia Europea per i Medicinali (EMA).

Questo capitolo fornisce le basi concettuali e tecniche che gli specializzandi in patologia necessitano per interagire criticamente con gli strumenti basati sull'IA nella loro pratica professionale. Inizia con una precisa tassonomia dell'IA e delle sue sotto-discipline, poi traccia l'arco storico dai sistemi simbolici basati su regole attraverso l'apprendimento automatico classico fino all'apprendimento profondo moderno. Il ragionamento probabilistico bayesiano viene introdotto come quadro complementare per il supporto decisionale diagnostico. Il capitolo si conclude esaminando il ruolo specifico e indispensabile del tecnico di laboratorio biomedico nei flussi di lavoro assistiti dall'IA, e fornisce un'attività pratica utilizzando la piattaforma open source di analisi delle immagini QuPath \cite{bankhead2017} per illustrare l'IA simbolica in azione.

% ============================================================
\section{Definizioni e Tassonomia dell'Intelligenza Artificiale}
\label{sec:ch3_taxonomy}
% ============================================================

\subsection{Che cos'è l'Intelligenza Artificiale?}

L'intelligenza artificiale è un ampio campo dell'informatica che si occupa della progettazione di sistemi che esibiscono comportamenti tipicamente associati all'intelligenza umana: ragionamento, apprendimento, risoluzione di problemi, percezione e comprensione del linguaggio. Il termine fu coniato da John McCarthy alla Conferenza di Dartmouth del 1956, dove fu proposto che ``ogni aspetto dell'apprendimento o qualsiasi altra caratteristica dell'intelligenza può in linea di principio essere descritta con tale precisione da consentire a una macchina di simularla.'' Nell'uso contemporaneo, l'IA comprende un ampio spettro di approcci, dai semplici programmi basati su regole che seguono istruzioni logiche esplicite alle sofisticate reti neurali che apprendono pattern statistici da vasti dataset \cite{rashidi2019}.

Una distinzione fondamentale nel campo è tra \textbf{IA ristretta} (detta anche IA debole) e \textbf{IA generale} (detta anche IA forte o intelligenza artificiale generale, AGI). L'IA ristretta si riferisce a sistemi progettati e addestrati per svolgere un compito specifico e ben definito — come classificare un'immagine istologica come maligna o benigna, o rilevare figure mitotiche in una regione di interesse definita. Tutti i sistemi di IA attualmente dispiegati in patologia clinica sono IA ristretta: sono altamente capaci all'interno del loro dominio definito ma non possono trasferire le loro conoscenze a compiti non correlati senza riaddestrare. L'IA generale, al contrario, denoterebbe un sistema capace di svolgere qualsiasi compito intellettuale che un essere umano può svolgere, con un ragionamento flessibile tra domini. Nessun sistema del genere esiste oggi; l'AGI rimane un'aspirazione di ricerca a lungo termine piuttosto che una realtà clinica.

\subsection{La Tassonomia dell'IA: IA, Apprendimento Automatico e Apprendimento Profondo}

La relazione tra IA, apprendimento automatico e apprendimento profondo è meglio compresa come un insieme di discipline annidate, ciascuna più specifica dell'ultima (Figura~\ref{fig:ai_taxonomy}). \textbf{L'apprendimento automatico} (ML) è un sotto-campo dell'IA in cui i sistemi imparano a svolgere compiti identificando pattern nei dati, piuttosto che seguendo regole esplicitamente programmate. \textbf{L'apprendimento profondo} (DL) è un sotto-campo del ML che utilizza reti neurali artificiali con molti strati (da qui ``profondo'') per apprendere rappresentazioni gerarchiche dei dati. Tutto l'apprendimento profondo è apprendimento automatico, e tutto l'apprendimento automatico è IA, ma non tutta l'IA è apprendimento automatico, e non tutto l'apprendimento automatico è apprendimento profondo.

\begin{figure}[htbp]
    \centering
    \usetikzlibrary{positioning, calc, fit}

\begin{tikzpicture}[
    every node/.style={font=\small\sffamily},
    exbox/.style={
        rectangle, rounded corners=3pt,
        fill=white, draw=gray!50, line width=0.5pt,
        minimum width=2.6cm, minimum height=0.6cm,
        align=center, font=\footnotesize\sffamily
    },
    colhead/.style={font=\small\sffamily\bfseries, align=center}
]

%% ── Nested rectangles ──────────────────────────────────────────────
%% Outermost: Artificial Intelligence (light blue)
\fill[rounded corners=8pt,
      fill={rgb,255:red,190;green,225;blue,245}]
    (-5.4, 2.6) rectangle (5.4, -1.0);

%% Middle: Machine Learning (medium blue)
\fill[rounded corners=6pt,
      fill={rgb,255:red,100;green,170;blue,220}]
    (-3.6, 2.1) rectangle (3.6, -0.6);

%% Innermost: Deep Learning (dark blue)
\fill[rounded corners=4pt,
      fill={rgb,255:red,30;green,100;blue,170}]
    (-1.8, 1.6) rectangle (1.8, -0.1);

%% Labels inside each region
\node[text={rgb,255:red,220;green,240;blue,255},
      font=\normalsize\sffamily\bfseries]
    at (0, 0.75) {Deep Learning};

\node[text=white, font=\small\sffamily\bfseries]
    at (-2.7, 1.85) {Machine Learning};

\node[text={rgb,255:red,20;green,60;blue,110},
      font=\small\sffamily\bfseries]
    at (-4.1, 2.3) {Artificial Intelligence};

%% ── Example columns below ──────────────────────────────────────────
\def\rowA{-1.65}
\def\rowB{-2.35}
\def\rowC{-3.05}

%% Column headers
\node[colhead, text={rgb,255:red,20;green,60;blue,110}]
    at (-3.8, -1.25) {AI Examples};
\node[colhead, text={rgb,255:red,30;green,90;blue,160}]
    at (0,    -1.25) {ML Examples};
\node[colhead, text={rgb,255:red,30;green,100;blue,170}]
    at (3.8,  -1.25) {DL Examples};

%% AI column
\node[exbox, fill={rgb,255:red,220;green,238;blue,250},
      draw={rgb,255:red,100;green,170;blue,220}]
    at (-3.8, \rowA) {Expert Systems};
\node[exbox, fill={rgb,255:red,220;green,238;blue,250},
      draw={rgb,255:red,100;green,170;blue,220}]
    at (-3.8, \rowB) {Rule-based Systems};

%% ML column
\node[exbox, fill={rgb,255:red,200;green,228;blue,248},
      draw={rgb,255:red,70;green,140;blue,200}]
    at (0, \rowA) {Random Forests};
\node[exbox, fill={rgb,255:red,200;green,228;blue,248},
      draw={rgb,255:red,70;green,140;blue,200}]
    at (0, \rowB) {SVMs};
\node[exbox, fill={rgb,255:red,200;green,228;blue,248},
      draw={rgb,255:red,70;green,140;blue,200}]
    at (0, \rowC) {Bayesian Networks};

%% DL column
\node[exbox, fill={rgb,255:red,180;green,215;blue,245},
      draw={rgb,255:red,30;green,100;blue,170}]
    at (3.8, \rowA) {CNNs};
\node[exbox, fill={rgb,255:red,180;green,215;blue,245},
      draw={rgb,255:red,30;green,100;blue,170}]
    at (3.8, \rowB) {Transformers};
\node[exbox, fill={rgb,255:red,180;green,215;blue,245},
      draw={rgb,255:red,30;green,100;blue,170}]
    at (3.8, \rowC) {Foundation Models};

%% Thin separator line
\draw[gray!30, line width=0.4pt]
    (-5.4, -1.1) -- (5.4, -1.1);

\end{tikzpicture}

    \caption{Tassonomia dell'intelligenza artificiale come discipline annidate. L'apprendimento automatico (ML) è un sottoinsieme dell'IA; l'apprendimento profondo (DL) è un sottoinsieme del ML. Ogni dominio interno impiega approcci più orientati ai dati e all'apprendimento delle rappresentazioni rispetto al suo genitore. Sono indicati esempi di metodi all'interno di ciascun dominio.}
    \label{fig:ai_taxonomy}
\end{figure}

\subsection{Terminologia Chiave}

Precise use of terminology is essential for critical engagement with the AI literature. The following definitions are used consistently throughout this handbook:

\begin{description}
    \item[\textbf{Algoritmo}] Una sequenza finita e ordinata di istruzioni computazionali che trasforma un input in un output. Un algoritmo è la procedura; un modello è l'artefatto prodotto applicando tale procedura ai dati.

    \item[\textbf{Modello}] Una funzione matematica, parametrizzata da valori appresi dai dati, che mappa gli input (ad es., i valori dei pixel di un'immagine istologica) sugli output (ad es., una probabilità di malignià). Il modello codifica la conoscenza estratta durante l'addestramento.

    \item[\textbf{Addestramento}] Il processo mediante il quale i parametri di un modello vengono regolati, tipicamente minimizzando una funzione di perdita su un dataset etichettato, in modo che le previsioni del modello diventino sempre più accurate. L'addestramento è computazionalmente intensivo e viene eseguito una volta (o aggiornato periodicamente) prima del dispiegamento.

    \item[\textbf{Inferenza}] L'applicazione di un modello addestrato a nuovi dati non visti in precedenza per generare previsioni. L'inferenza è ciò che avviene nel dispiegamento clinico: il modello è fisso e i nuovi casi vengono passati attraverso di esso per ottenere output.

    \item[\textbf{Etichetta}] L'annotazione ground truth associata a un esempio di addestramento. Nell'apprendimento supervisionato, le etichette sono fornite da esperti umani (ad es., un patologo che annota una regione come ``carcinoma invasivo'' o ``epitelio normale''). La qualità e la coerenza delle etichette determinano direttamente la qualità del modello addestrato.

    \item[\textbf{Caratteristica}] Una proprietà o caratteristica misurabile dei dati di input utilizzata da un modello per fare previsioni. Nel ML classico, le caratteristiche sono progettate manualmente da esperti del dominio (ad es., area nucleare, texture della cromatina, rapporto ghiandola-stroma). Nell'apprendimento profondo, le caratteristiche vengono apprese automaticamente dai dati grezzi.
\end{description}

La comprensione di questi termini consente al professionista di leggere criticamente i lavori sull'IA, di interrogare le affermazioni dei fornitori sugli strumenti diagnostici basati sull'IA e di partecipare in modo significativo alla validazione e alla governance dei sistemi di IA all'interno del laboratorio.

% ============================================================
\section{IA Simbolica e Sistemi Basati su Regole}
\label{sec:ch3_symbolic}
% ============================================================

\subsection{Sistemi Esperti e Ragionamento Basato su Regole}

I primi sistemi di IA pratici in medicina erano \textbf{sistemi esperti}: programmi che codificavano la conoscenza di specialisti umani come un insieme strutturato di regole logiche, tipicamente nella forma \textit{SE} [condizione] \textit{ALLORA} [conclusione]. Il sistema MYCIN, sviluppato all'Università di Stanford negli anni Settanta per raccomandare la terapia antibiotica per le infezioni batteriche, è l'esempio canonico. I sistemi esperti sono composti da due componenti: una \textbf{base di conoscenza} contenente le regole specifiche del dominio, e un \textbf{motore di inferenza} che applica tali regole a un dato insieme di input per giungere a una conclusione. In patologia, i sistemi basati su regole sono stati implementati in software di analisi delle immagini per svolgere compiti come la segmentazione cellulare, la quantificazione della colorazione e la misurazione morfometrica.

Gli \textbf{alberi decisionali} sono un formalismo strettamente correlato in cui una serie di domande binarie sulle caratteristiche di input sono disposte in una struttura ad albero gerarchica. In ogni nodo interno, l'algoritmo verifica una caratteristica rispetto a una soglia; il risultato determina quale ramo viene seguito, fino a raggiungere un nodo foglia che specifica la classe o il valore previsto. Gli alberi decisionali sono intrinsecamente interpretabili: il percorso dalla radice alla foglia costituisce una spiegazione leggibile dall'uomo della previsione. Erano ampiamente utilizzati nei primi sistemi di diagnosi assistita da computer (CAD) e rimangono preziosi come componenti di metodi ensemble come le foreste casuali (Sezione~\ref{sec:ch3_ml}).

\subsection{La Valutazione del Ki-67 come Esempio di Analisi delle Immagini Basata su Regole}

Il Ki-67 è una proteina nucleare espressa nelle cellule in proliferazione e viene valutata di routine mediante immunoistochimica (IHC) nel cancro al seno per stimare l'indice di proliferazione — la proporzione di cellule tumorali che ciclano attivamente attraverso il ciclo cellulare. L'indice di proliferazione Ki-67 ha significato prognostico e predittivo: un'elevata espressione di Ki-67 è associata a un comportamento tumorale aggressivo e può influenzare le decisioni riguardanti la chemioterapia adiuvante \cite{bankhead2022}.

In un sistema di analisi delle immagini basato su regole, la valutazione del Ki-67 procede attraverso una sequenza deterministica di passaggi. In primo luogo, l'immagine viene convertita dallo spazio colore RGB standard a una rappresentazione con separazione delle colorazioni mediante deconvoluzione di Beer-Lambert, isolando il canale della diaminobenzidina (DAB) (marrone, indicante la positività al Ki-67) dal canale dell'ematossilina (blu, indicante tutti i nuclei). In secondo luogo, un algoritmo di segmentazione nucleare — tipicamente basato sulla sogliatura dell'intensità seguita dalla trasformata watershed — identifica i singoli nuclei come oggetti discreti. In terzo luogo, per ogni nucleo segmentato viene calcolata la densità ottica media del DAB. In quarto luogo, viene applicata una soglia: i nuclei con densità ottica media del DAB superiore a un valore definito (ad es., 0,2 unità di assorbanza) sono classificati come \textbf{positivi}; quelli al di sotto della soglia sono classificati come \textbf{negativi}. Infine, l'indice di proliferazione viene calcolato come:

\begin{equation}
    \text{Indice Ki-67} = \frac{N_{\text{positive}}}{N_{\text{positive}} + N_{\text{negative}}} \times 100\%
\end{equation}

dove $N_{\text{positive}}$ e $N_{\text{negative}}$ sono rispettivamente i conteggi dei nuclei DAB-positivi e DAB-negativi. L'intera procedura è implementata in piattaforme come QuPath \cite{bankhead2017}, dove la soglia e i parametri di segmentazione sono impostati dall'utente e applicati uniformemente all'immagine.

\subsection{Punti di Forza e Limitazioni dei Sistemi Basati su Regole}

I sistemi basati su regole offrono diversi vantaggi importanti. Sono \textbf{interpretabili}: ogni decisione può essere ricondotta a una regola specifica, rendendo il comportamento del sistema trasparente e verificabile. Sono \textbf{deterministici}: dato lo stesso input e le stesse regole, l'output è sempre identico, il che è essenziale per la riproducibilità in un contesto diagnostico. Sono anche \textbf{veloci} e computazionalmente poco costosi, non richiedendo dati di addestramento né hardware GPU. Per compiti ben definiti e vincolati — come il conteggio dei nuclei DAB-positivi in un preparato IHC colorato uniformemente — i sistemi basati su regole possono raggiungere prestazioni eccellenti e sono del tutto adeguati allo scopo.

Tuttavia, i sistemi basati su regole presentano limitazioni fondamentali che diventano evidenti quando il compito è più complesso o i dati più variabili. Sono \textbf{fragili}: piccole variazioni nel protocollo di colorazione, nella processazione tissutale o nella calibrazione dello scanner possono spostare la distribuzione della densità ottica del DAB, causando la classificazione errata dei nuclei da parte della soglia fissa. Sono \textbf{difficili da generalizzare}: un insieme di regole ottimizzato per il protocollo di colorazione Ki-67 di un laboratorio può avere prestazioni scadenti quando applicato a vetrini di un'istituzione diversa. \textbf{Non possono apprendere}: se le regole sono errate o incomplete, il sistema fallirà sistematicamente, e il miglioramento richiede la revisione manuale della base di conoscenza da parte di un esperto del dominio. Inoltre, i sistemi basati su regole sono poco adatti a compiti che richiedono l'integrazione di informazioni contestuali — ad esempio, distinguere i linfociti infiltranti il tumore dalle cellule tumorali in un campo densamente cellulare — perché tali distinzioni dipendono da relazioni spaziali e dal contesto morfologico che sono difficili da codificare come regole esplicite \cite{rashidi2019}.

Queste limitazioni hanno motivato lo sviluppo di approcci di apprendimento automatico, che possono adattare il loro comportamento alle proprietà statistiche dei dati di addestramento e generalizzare in modo più robusto in condizioni variabili.

% ============================================================
\section{Apprendimento Automatico: Approcci Supervisionati e Non Supervisionati}
\label{sec:ch3_ml}
% ============================================================

\subsection{Il Paradigma dell'Apprendimento Automatico}

L'apprendimento automatico rappresenta un cambiamento fondamentale nel modo in cui i sistemi computazionali vengono progettati. Invece di specificare le regole che governano il comportamento di un sistema, lo sviluppatore specifica un'\textbf{architettura del modello} (la forma matematica della funzione da apprendere) e un \textbf{algoritmo di apprendimento} (la procedura per regolare i parametri del modello), e poi espone il sistema ai dati da cui inferisce automaticamente le regole appropriate. Questo approccio basato sui dati consente ai sistemi ML di catturare relazioni complesse e non lineari in dati ad alta dimensionalità che sarebbero intrattabili da codificare manualmente \cite{komura2018}.

Il concetto centrale nel ML è il \textbf{vettore di caratteristiche}: una rappresentazione numerica di un esempio di input, costruita misurando un insieme di proprietà predefinite. Nell'analisi delle immagini istopatologiche, le caratteristiche possono includere area nucleare, perimetro, circolarità, intensità media della cromatina, descrittori di texture (ad es., caratteristiche di Haralick derivate dalla matrice di co-occorrenza dei livelli di grigio) e statistiche spaziali che descrivono la disposizione delle cellule all'interno di una regione tissutale. La qualità della rappresentazione delle caratteristiche è critica: se le caratteristiche non catturano le informazioni rilevanti per il compito, nessun algoritmo di apprendimento può compensare questa carenza. L'ingegneria delle caratteristiche — il processo di progettazione di caratteristiche informative — era storicamente un importante collo di bottiglia nel ML applicato alla patologia, richiedendo una profonda competenza nel dominio e un notevole sforzo manuale.

\subsection{Apprendimento Supervisionato}

Nell'\textbf{apprendimento supervisionato}, il modello viene addestrato su un dataset di coppie input-output $\{(\mathbf{x}_i, y_i)\}_{i=1}^{N}$, dove $\mathbf{x}_i$ è il vettore di caratteristiche dell'$i$-esimo esempio e $y_i$ è la sua etichetta. L'algoritmo di apprendimento regola i parametri del modello per minimizzare una \textbf{funzione di perdita} $\mathcal{L}(\hat{y}, y)$ che quantifica la discrepanza tra le previsioni del modello $\hat{y}$ e le etichette vere $y$. Esistono due tipi principali di compiti nell'apprendimento supervisionato:

\begin{itemize}
    \item \textbf{Classificazione}: L'output $y$ è una categoria discreta (ad es., ``maligno'' vs.\ ``benigno''; ``Grado 1'', ``Grado 2'', o ``Grado 3''). Le funzioni di perdita comuni includono la cross-entropy loss. Le prestazioni vengono valutate utilizzando metriche come accuratezza, sensibilità, specificità, area sotto la curva ROC (AUROC) e il punteggio F1.

    \item \textbf{Regressione}: L'output $y$ è un valore continuo (ad es., indice Ki-67 previsto, percentuale di cellularità tumorale o tempo di sopravvivenza del paziente). Le funzioni di perdita comuni includono l'errore quadratico medio (MSE) e l'errore assoluto medio (MAE).
\end{itemize}

La qualità dell'apprendimento supervisionato è fondamentalmente vincolata dalla qualità e dalla quantità dei dati di addestramento etichettati. Ottenere annotazioni di esperti per grandi dataset istopatologici è costoso e dispendioso in termini di tempo, e la variabilità inter-osservatore tra i patologi introduce rumore nelle etichette che può degradare le prestazioni del modello \cite{bera2019}.

\subsection{Apprendimento Non Supervisionato e Semi-Supervisionato}

L'\textbf{apprendimento non supervisionato} opera su dati non etichettati, cercando di scoprire strutture intrinseche — come cluster, varietà o fattori generativi — senza riferimento a categorie predefinite. Gli algoritmi di \textbf{clustering}, come $k$-means e il clustering agglomerativo gerarchico, partizionano i dati in gruppi di esempi simili. In patologia, il clustering non supervisionato è stato applicato per identificare sottotipi tumorali morfologicamente distinti, per stratificare i pazienti in base alla composizione del microambiente tissutale e per scoprire nuovi fenotipi istologici associati ad alterazioni molecolari. I metodi di riduzione della dimensionalità come l'analisi delle componenti principali (PCA) e la proiezione e approssimazione di varietà uniforme (UMAP) sono frequentemente utilizzati in combinazione con il clustering per visualizzare spazi di caratteristiche ad alta dimensionalità.

L'\textbf{apprendimento semi-supervisionato} occupa lo spazio tra gli approcci supervisionati e non supervisionati, sfruttando un piccolo dataset etichettato insieme a un dataset non etichettato molto più grande. Questo è particolarmente rilevante in patologia, dove le WSI non etichettate sono abbondanti ma le annotazioni di esperti sono scarse. I metodi semi-supervisionati utilizzano la struttura dei dati non etichettati — ad esempio, l'assunzione che i punti vicini nello spazio delle caratteristiche debbano avere la stessa etichetta — per regolarizzare il modello e migliorare la generalizzazione oltre quanto sarebbe ottenibile con i soli dati etichettati.

\subsection{Algoritmi Classici di Apprendimento Automatico}

Diversi algoritmi classici di ML rimangono ampiamente utilizzati nella ricerca in patologia e nelle applicazioni cliniche:

La \textbf{regressione logistica} è un modello lineare per la classificazione binaria che stima la probabilità di appartenenza a una classe come funzione sigmoide di una combinazione lineare delle caratteristiche di input. Nonostante la sua semplicità, la regressione logistica è robusta, interpretabile e funziona bene quando il confine decisionale è approssimativamente lineare nello spazio delle caratteristiche. Viene frequentemente utilizzata come modello di riferimento e per compiti in cui l'interpretabilità è fondamentale.

Le \textbf{macchine a vettori di supporto} (SVM) trovano l'iperpiano nello spazio delle caratteristiche che separa massimamente le due classi, massimizzando il margine tra gli esempi di addestramento più vicini (i vettori di supporto) e il confine decisionale. Le SVM possono essere estese alla classificazione non lineare attraverso il \textbf{kernel trick}, che mappa implicitamente i dati in uno spazio di dimensionalità superiore dove la separazione lineare è possibile. Le SVM erano l'approccio dominante nella patologia computazionale negli anni Duemila e all'inizio degli anni Duemiladieci, raggiungendo prestazioni elevate in compiti come il rilevamento delle mitosi e la classificazione tumorale.

Le \textbf{foreste casuali} sono metodi ensemble che costruiscono un gran numero di alberi decisionali, ciascuno addestrato su un campione bootstrap casuale dei dati di addestramento e utilizzando un sottoinsieme casuale di caratteristiche ad ogni divisione. La previsione finale si ottiene per voto di maggioranza (classificazione) o mediando (regressione) su tutti gli alberi. Le foreste casuali sono robuste all'overfitting, gestiscono bene gli spazi di caratteristiche ad alta dimensionalità e forniscono una misura naturale dell'importanza delle caratteristiche, rendendole preziose sia per la previsione che per l'analisi esplorativa \cite{rashidi2019}.

\subsection{Overfitting, Regolarizzazione e Validazione Incrociata}

Una sfida centrale nel ML è l'\textbf{overfitting}: la tendenza di un modello a memorizzare i dati di addestramento, inclusi il loro rumore e le loro idiosincrasie, piuttosto che apprendere il pattern generalizzabile sottostante. Un modello in overfitting raggiunge un'elevata accuratezza sul set di addestramento ma ha prestazioni scadenti su nuovi dati non visti. L'overfitting è più probabile quando il modello è complesso (molti parametri) rispetto alla dimensione del dataset di addestramento.

Le tecniche di \textbf{regolarizzazione} mitigano l'overfitting penalizzando la complessità del modello. La regolarizzazione L1 (Lasso) aggiunge una penale proporzionale alla somma dei valori assoluti dei parametri del modello, incoraggiando soluzioni sparse. La regolarizzazione L2 (Ridge) aggiunge una penale proporzionale alla somma dei parametri al quadrato, riducendo tutti i parametri verso zero. Il dropout, utilizzato nelle reti neurali, disattiva casualmente una proporzione di neuroni durante l'addestramento, prevenendo la co-adattazione.

La \textbf{validazione incrociata} è la metodologia standard per stimare le prestazioni di generalizzazione di un modello. Nella validazione incrociata a $k$ fold, il dataset viene partizionato in $k$ sottoinsiemi uguali (fold); il modello viene addestrato su $k-1$ fold e valutato sul fold trattenuto, e questo processo viene ripetuto $k$ volte in modo che ogni esempio serva come caso di test esattamente una volta. La media e la deviazione standard delle prestazioni tra i fold forniscono una stima imparziale delle prestazioni di generalizzazione. In patologia, è fondamentale che le divisioni della validazione incrociata vengano eseguite a livello di paziente (non a livello di immagine o patch) per prevenire la fuga di dati, in cui le informazioni dello stesso paziente compaiono sia nel set di addestramento che in quello di test.

% ============================================================
\section{Reti Bayesiane e Ragionamento Probabilistico}
\label{sec:ch3_bayesian}
% ============================================================

\subsection{Inferenza Bayesiana}

L'inferenza bayesiana fornisce un quadro matematico rigoroso per aggiornare le credenze alla luce delle prove. Si basa sul \textbf{teorema di Bayes}:

\begin{equation}
    P(H \mid E) = \frac{P(E \mid H) \cdot P(H)}{P(E)}
\end{equation}

dove $H$ è un'ipotesi (ad es., ``la lesione è maligna''), $E$ è la prova osservata (ad es., i risultati di un pannello di colorazioni immunoistochimiche), $P(H)$ è la \textbf{probabilità a priori} dell'ipotesi prima che la prova venga considerata, $P(E \mid H)$ è la \textbf{verosimiglianza} di osservare la prova dato che l'ipotesi è vera, $P(H \mid E)$ è la \textbf{probabilità a posteriori} dell'ipotesi dopo aver incorporato la prova, e $P(E)$ è la probabilità marginale della prova (una costante di normalizzazione). In patologia diagnostica, la probabilità a priori corrisponde alla probabilità pre-test di una diagnosi dato il contesto clinico (età del paziente, sesso, sede della biopsia, risultati dell'imaging), e la probabilità a posteriori è la probabilità aggiornata dopo che i risultati istologici e ancillari sono stati incorporati.

Il quadro bayesiano è particolarmente adatto al ragionamento diagnostico perché modella esplicitamente l'incertezza. Invece di produrre una singola diagnosi deterministica, un sistema bayesiano produce una distribuzione di probabilità sulle possibili diagnosi, quantificando il grado di fiducia in ciascuna. Questo è clinicamente prezioso: un sistema che riporta ``probabilità di carcinoma sieroso ad alto grado: 0,87; probabilità di carcinoma endometrioide: 0,11; probabilità di carcinoma a cellule chiare: 0,02'' trasmette informazioni più utili di una classificazione binaria, e consente al clinico di calibrare la propria risposta in modo appropriato.

\subsection{Reti Bayesiane}

Una \textbf{rete bayesiana} (BN) è un modello grafico probabilistico che rappresenta un insieme di variabili e le loro dipendenze condizionali come un \textbf{grafo aciclico diretto} (DAG). Ogni nodo nel grafo rappresenta una variabile casuale (ad es., un reperto clinico, una caratteristica istologica o una diagnosi), e ogni arco diretto rappresenta una dipendenza condizionale tra variabili. L'assenza di un arco tra due nodi codifica l'assunzione di indipendenza condizionale, che consente di fattorizzare la distribuzione di probabilità congiunta su tutte le variabili in un prodotto di distribuzioni di probabilità condizionali locali — una per ogni nodo dati i suoi genitori nel grafo.

In una rete bayesiana diagnostica per la patologia mammaria, ad esempio, i nodi potrebbero rappresentare variabili come l'età del paziente, il grado tumorale, lo stato del recettore degli estrogeni (ER), lo stato del recettore del progesterone (PR), lo stato HER2, l'indice Ki-67 e la presenza di invasione linfovascolare. Gli archi codificano le dipendenze biologiche e statistiche note: ad esempio, lo stato ER influenza lo stato PR, ed entrambi influenzano la probabilità di risposta alla terapia endocrina. Dati i valori osservati per un sottoinsieme di nodi (la prova), la rete può calcolare la distribuzione di probabilità a posteriori su tutti i nodi non osservati utilizzando algoritmi di inferenza esatta o approssimata.

\subsection{Vantaggi e Limitazioni nel Supporto Decisionale Diagnostico}

Le reti bayesiane offrono diversi vantaggi per il supporto decisionale diagnostico in patologia. In primo luogo, forniscono \textbf{quantificazione dell'incertezza}: l'output è una distribuzione di probabilità, non una stima puntuale, rendendo esplicita la fiducia del sistema. In secondo luogo, sono \textbf{interpretabili}: la struttura del grafo codifica la conoscenza del dominio in una forma che i clinici possono ispezionare e validare. In terzo luogo, possono \textbf{gestire i dati mancanti} in modo elegante: se un particolare risultato di test non è disponibile, la rete marginalizza sulla variabile mancante invece di fallire. In quarto luogo, possono \textbf{incorporare la conoscenza a priori}: la struttura e i parametri iniziali della rete possono essere specificati da esperti del dominio e poi raffinati utilizzando i dati, combinando i punti di forza degli approcci basati sulla conoscenza e basati sui dati.

Le principali limitazioni delle reti bayesiane sono la loro \textbf{scalabilità} e la loro \textbf{assunzione di indipendenza condizionale}. All'aumentare del numero di variabili, il numero di parametri necessari per specificare le tabelle di probabilità condizionale cresce esponenzialmente, rendendo l'inferenza esatta computazionalmente intrattabile per reti di grandi dimensioni. I metodi di inferenza approssimata (ad es., Monte Carlo a catena di Markov, inferenza variazionale) possono affrontare questo problema, ma a costo di una complessità aggiuntiva. Inoltre, le assunzioni di indipendenza condizionale codificate nella struttura del grafo potrebbero non valere in pratica, e la specificazione errata del grafo può portare a un'inferenza distorta. Nonostante queste limitazioni, le reti bayesiane sono state applicate con successo al supporto decisionale diagnostico in oncologia, radiologia e patologia, e rimangono un'area di ricerca attiva \cite{rashidi2019}.

% ============================================================
\section{Dalle Caratteristiche alle Rappresentazioni: Il Passaggio all'Apprendimento Profondo}
\label{sec:ch3_deeplearning}
% ============================================================

\subsection{Le Limitazioni delle Caratteristiche Costruite Manualmente}

Gli approcci classici di ML in patologia dipendono criticamente dalla qualità delle caratteristiche costruite manualmente utilizzate per rappresentare le immagini. La progettazione di caratteristiche informative richiede una profonda competenza nel dominio, è dispendiosa in termini di tempo ed è intrinsecamente limitata dalla conoscenza pregressa del progettista su ciò che è diagnosticamente rilevante. Le caratteristiche efficaci per un compito (ad es., la morfometria nucleare per la classificazione) possono essere non informative per un altro (ad es., la predizione dell'instabilità dei microsatelliti dalla morfologia E\&E). Inoltre, le caratteristiche costruite manualmente vengono tipicamente calcolate da singole cellule o piccole regioni di interesse, e possono non riuscire a catturare l'organizzazione spaziale del tessuto a più scale — la disposizione delle ghiandole, la composizione del microambiente tumorale, la relazione tra tumore e stroma — che è sempre più riconosciuta come diagnosticamente e prognosticamente importante \cite{bera2019}.

L'intuizione fondamentale dell'apprendimento profondo è che, dati dati sufficienti e risorse computazionali, una rete neurale può \textbf{apprendere} la rappresentazione ottimale delle caratteristiche per un dato compito direttamente dai dati grezzi dei pixel, senza alcuna ingegneria manuale delle caratteristiche. Le caratteristiche apprese non sono interpretabili allo stesso modo delle misurazioni morfometriche costruite manualmente, ma possono catturare pattern complessi e non lineari che sono invisibili all'ingegneria convenzionale delle caratteristiche. Questo passaggio dall'ingegneria delle caratteristiche all'\textbf{apprendimento delle rappresentazioni} è la caratteristica distintiva della rivoluzione dell'apprendimento profondo nella patologia computazionale \cite{song2023}.

\subsection{Reti Neurali Artificiali}

Una \textbf{rete neurale artificiale} (ANN) è un modello computazionale liberamente ispirato alla struttura dei circuiti neurali biologici. È composta da una serie di \textbf{strati}, ciascuno comprendente un insieme di unità computazionali chiamate \textbf{neuroni} (o nodi). Ogni neurone calcola una somma pesata dei suoi input, aggiunge un termine di bias e passa il risultato attraverso una \textbf{funzione di attivazione} non lineare per produrre il suo output. Le funzioni di attivazione comuni includono l'unità lineare rettificata (ReLU, $f(x) = \max(0, x)$), la funzione sigmoide ($f(x) = 1/(1+e^{-x})$) e la tangente iperbolica (tanh). La non linearità della funzione di attivazione è essenziale: senza di essa, una rete multistrato sarebbe matematicamente equivalente a una singola trasformazione lineare, indipendentemente dalla sua profondità.

Una rete con uno o più \textbf{strati nascosti} tra gli strati di input e output è chiamata \textbf{percettrone multistrato} (MLP). La \textbf{profondità} di una rete si riferisce al numero di strati; le reti profonde (con molti strati) possono rappresentare funzioni dell'input sempre più astratte e complesse. I parametri della rete — i pesi e i bias di tutti i neuroni — vengono appresi durante l'addestramento mediante \textbf{backpropagation}: il gradiente della funzione di perdita rispetto a ciascun parametro viene calcolato usando la regola della catena del calcolo, e i parametri vengono aggiornati nella direzione che riduce la perdita, utilizzando un algoritmo di ottimizzazione come la discesa del gradiente stocastica (SGD) o Adam.

\subsection{Reti Neurali Convoluzionali per l'Analisi delle Immagini}

Per i compiti di analisi delle immagini, l'architettura di rete neurale più importante è la \textbf{rete neurale convoluzionale} (CNN). Le CNN sfruttano la struttura spaziale delle immagini attraverso due operazioni chiave. Prima, la \textbf{convoluzione}: un piccolo filtro (kernel) di pesi apprendibili viene fatto scorrere sull'immagine, calcolando il prodotto scalare tra il filtro e ogni patch locale di pixel. Questo produce una \textbf{mappa delle caratteristiche} che codifica la presenza e la posizione di un particolare pattern (ad es., un bordo, una texture o un motivo più complesso) nell'immagine. Più filtri vengono applicati in parallelo, producendo più mappe delle caratteristiche. Seconda, il \textbf{pooling}: le mappe delle caratteristiche vengono sottocampionate spazialmente (ad es., prendendo il valore massimo in ogni regione $2 \times 2$ non sovrapposta), riducendo la risoluzione spaziale mantenendo le attivazioni più salienti e introducendo un grado di invarianza alla traslazione.

Impilando più strati convoluzionali e di pooling, una CNN costruisce una rappresentazione gerarchica dell'immagine: gli strati iniziali rilevano caratteristiche di basso livello come bordi e gradienti di colore; gli strati intermedi rilevano caratteristiche di medio livello come texture e forme; e gli strati profondi rilevano caratteristiche di alto livello, semanticamente significative, come strutture ghiandolari, pleomorfismo nucleare o figure mitotiche. Questo apprendimento gerarchico delle caratteristiche è ciò che rende le CNN così potenti per l'analisi delle immagini istopatologiche \cite{komura2018}. Architetture come ResNet, VGG, Inception ed EfficientNet, originariamente sviluppate per il riconoscimento di immagini naturali, sono state ampiamente adattate e ottimizzate per applicazioni in patologia, spesso utilizzando il \textbf{transfer learning}: inizializzando la rete con pesi pre-addestrati su un grande dataset di immagini naturali (ad es., ImageNet) e poi ottimizzando su un dataset più piccolo specifico per la patologia.

\subsection{Oltre le CNN: Transformer e Modelli Fondazionali}

Più recentemente, i \textbf{vision transformer} (ViT) e i \textbf{modelli fondazionali} sono emersi come potenti alternative alle CNN per la patologia computazionale. I transformer, originariamente sviluppati per l'elaborazione del linguaggio naturale, utilizzano un meccanismo di \textbf{auto-attenzione} che consente a ogni elemento dell'input (ad es., ogni patch di una WSI) di prestare attenzione a tutti gli altri elementi, catturando dipendenze spaziali a lungo raggio che sono difficili da modellare per le CNN. I modelli fondazionali — grandi reti neurali pre-addestrate su dataset massicci e diversificati mediante apprendimento auto-supervisionato — possono essere ottimizzati per un'ampia gamma di compiti a valle con relativamente pochi dati etichettati. Modelli come UNI, CONCH e PLIP, pre-addestrati su milioni di immagini patologiche, rappresentano la frontiera attuale del settore e stanno iniziando a dimostrare prestazioni che si avvicinano o superano quelle dei patologi specialisti in una serie di compiti diagnostici \cite{song2023}.

% ============================================================
\section{Il Ruolo del Tecnico nei Processi Assistiti dall'IA}
\label{sec:ch3_technician}
% ============================================================

\subsection{L'Imperativo dell'Uomo nel Ciclo}

Il dispiegamento dell'IA in patologia clinica non elimina la necessità di competenza umana; piuttosto, trasforma e per certi aspetti eleva il ruolo del tecnico di laboratorio biomedico. I sistemi di IA sono strumenti, e come tutti gli strumenti possono fallire — a volte in modi sottili, sistematici e difficili da rilevare senza conoscenza del dominio. Il concetto di \textbf{uomo nel ciclo} (HITL) si riferisce all'integrazione del giudizio umano in punti critici di un flusso di lavoro assistito dall'IA, garantendo che gli output del sistema vengano esaminati, validati e contestualizzati da un professionista competente prima che influenzino le decisioni cliniche \cite{bankhead2022}. Il tecnico che comprende i principi alla base di un sistema di IA è molto meglio attrezzato per identificarne le modalità di fallimento rispetto a chi lo tratta come una scatola nera.

\subsection{Annotazione dei Dati}

Le prestazioni di qualsiasi modello ML supervisionato sono fondamentalmente determinate dalla qualità dei suoi dati di addestramento, e la qualità dei dati di addestramento è fondamentalmente determinata dalla qualità delle sue annotazioni. L'\textbf{annotazione dei dati} — il processo di assegnazione di etichette agli esempi di addestramento — è un compito qualificato che richiede competenza nel dominio, attenzione ai dettagli e consapevolezza delle fonti di variabilità e ambiguità nelle immagini istopatologiche. I tecnici e i biologi di laboratorio biomedico sono sempre più coinvolti nei flussi di lavoro di annotazione, lavorando a fianco dei patologi per etichettare cellule, regioni di interesse, compartimenti tissutali e categorie diagnostiche in grandi dataset di immagini.

Le migliori pratiche nell'annotazione includono l'uso di linee guida di annotazione chiare e operativamente definite; sessioni di calibrazione regolari per valutare e minimizzare la variabilità inter-annotatore; l'uso di più annotatori indipendenti per i casi ambigui, con i disaccordi risolti per consenso o da un patologo senior; e il mantenimento di registri di provenienza dettagliati che documentano chi ha annotato ogni caso, quando e secondo quale versione delle linee guida. La qualità dell'annotazione dovrebbe essere monitorata continuamente, e i casi con elevato disaccordo inter-annotatore dovrebbero essere segnalati per la revisione. L'investimento in annotazioni di alta qualità non è semplicemente una questione tecnica: è il fondamento su cui si basa la validità clinica del modello di IA risultante.

\subsection{Controllo Qualità dei Dati di Addestramento e degli Output dell'IA}

Oltre all'annotazione, i tecnici svolgono un ruolo critico nel \textbf{controllo qualità} (QC) dei dati di addestramento. Le fonti comuni di problemi di qualità dei dati nella patologia digitale includono: artefatti di processazione tissutale (ad es., pieghe, lacerazioni, bolle d'aria, fissazione eccessiva o insufficiente); variabilità della colorazione (ad es., variazione da lotto a lotto nell'intensità dell'ematossilina ed eosina); artefatti di scansione (ad es., errori di messa a fuoco, artefatti di cucitura, polvere sul vetro dello scanner); ed errori di metadati (ad es., identificatori di caso errati, diagnosi etichettate erroneamente). Ciascuno di questi può introdurre bias sistematici in un modello addestrato. Una rigorosa pipeline di QC dovrebbe esaminare le immagini di addestramento per questi artefatti prima che entrino nel dataset di addestramento, e dovrebbe documentare i criteri di esclusione applicati.

Il controllo qualità degli output dell'IA nel dispiegamento clinico è ugualmente importante. I tecnici dovrebbero essere addestrati a riconoscere le modalità di fallimento caratteristiche dei sistemi di IA che utilizzano: ad esempio, un algoritmo di quantificazione del Ki-67 può sovrastimare i nuclei nelle regioni di piega tissutale, o può classificare erroneamente le cellule stromali come cellule tumorali nelle aree di desmoplasia. I controlli QC di routine — come l'ispezione visiva di un campione casuale di casi elaborati dall'IA, il confronto degli output dell'IA con i conteggi manuali su un sottoinsieme di casi e il monitoraggio delle distribuzioni degli output per spostamenti inattesi — sono salvaguardie essenziali. Qualsiasi discrepanza sistematica tra gli output dell'IA e la valutazione degli esperti dovrebbe essere segnalata e indagata.

\subsection{Validazione e Considerazioni Normative}

Prima che uno strumento di IA venga introdotto nella pratica clinica, deve essere rigorosamente validato su dati rappresentativi della popolazione di utilizzo prevista e che non sono stati utilizzati nell'addestramento. Gli studi di validazione dovrebbero valutare non solo le metriche di prestazione complessive (ad es., AUROC, sensibilità, specificità) ma anche le prestazioni tra sottogruppi clinicamente rilevanti (ad es., diversi sottotipi tumorali, diverse caratteristiche demografiche dei pazienti, diversi protocolli di colorazione) per identificare potenziali disparità. Il ruolo del tecnico nella validazione include la preparazione del dataset di validazione, l'esecuzione del sistema di IA sui casi di validazione, la registrazione degli output e il loro confronto con lo standard di riferimento stabilito dai patologi esperti.

I quadri normativi per i dispositivi diagnostici in vitro (IVD) basati sull'IA si stanno evolvendo rapidamente. Nell'Unione Europea, gli strumenti diagnostici basati sull'IA sono regolamentati dal Regolamento sui Dispositivi Diagnostici In Vitro (IVDR, UE 2017/746) e, ove applicabile, dall'AI Act (UE 2024/1689). Negli Stati Uniti, la FDA regola il software basato su IA/ML come dispositivo medico (SaMD) nell'ambito di un quadro che include disposizioni per i sistemi di apprendimento continuo. I tecnici dovrebbero essere consapevoli dello stato normativo degli strumenti di IA che utilizzano, delle condizioni in cui tali strumenti sono stati validati e degli obblighi del laboratorio riguardo alla sorveglianza post-commercializzazione e alla segnalazione degli incidenti \cite{bankhead2022}.

\subsection{Comprensione delle Limitazioni dell'IA}

Una competenza critica per qualsiasi professionista che lavora con l'IA è la capacità di riconoscere e articolare le limitazioni dei sistemi che utilizza. Le principali limitazioni di cui essere consapevoli includono:

\begin{itemize}
    \item \textbf{Spostamento della distribuzione}: I modelli di IA funzionano meglio su dati che assomigliano ai loro dati di addestramento. Quando la distribuzione degli input cambia — ad esempio, a causa di un cambiamento nel protocollo di colorazione, un nuovo modello di scanner o una diversa popolazione di pazienti — le prestazioni del modello possono degradarsi, a volte drasticamente, senza alcun segnale di avvertimento ovvio.

    \item \textbf{Apprendimento di scorciatoie}: I modelli possono imparare a sfruttare correlazioni spurie nei dati di addestramento (ad es., un particolare artefatto di colorazione che si trova a co-occorrere con una diagnosi nel set di addestramento) piuttosto che le vere caratteristiche diagnostiche. Tali scorciatoie potrebbero non essere evidenti dalle metriche di prestazione aggregate ma possono causare errori sistematici nel dispiegamento.

    \item \textbf{Calibrazione}: Le probabilità previste da un modello potrebbero non riflettere accuratamente la vera frequenza del risultato previsto. Un modello ben calibrato che prevede una probabilità del 70\% di malignià dovrebbe essere corretto circa il 70\% delle volte; un modello mal calibrato può essere sistematicamente troppo o troppo poco fiducioso.

    \item \textbf{Spiegabilità}: Molti modelli di apprendimento profondo sono difficili da interpretare: non è sempre possibile capire perché un modello ha fatto una particolare previsione. I metodi di spiegabilità come la mappatura dell'attivazione di classe pesata dal gradiente (Grad-CAM) e la visualizzazione dell'attenzione possono fornire una comprensione parziale, ma dovrebbero essere interpretati con cautela.
\end{itemize}

% ============================================================
\section{Attività Pratiche}
\label{sec:ch3_practical}
% ============================================================

\subsection{Attività 3.1: L'IA Simbolica in Azione --- Quantificazione del Ki-67 con QuPath}

\textbf{Obiettivo:} Dimostrare i principi dell'IA simbolica basata su regole implementando un algoritmo di rilevamento delle cellule positive al Ki-67 basato su soglia in QuPath \cite{bankhead2017}, ed esplorare la sensibilità del risultato alla scelta del parametro di soglia.

\textbf{Risultati di Apprendimento:} Al completamento di questa attività, gli studenti saranno in grado di: (1) spiegare i passaggi di una pipeline di analisi delle immagini basata su regole per la quantificazione IHC; (2) implementare un algoritmo di rilevamento delle cellule positive in QuPath utilizzando il classificatore basato su soglia integrato; (3) valutare criticamente l'effetto delle scelte dei parametri sull'indice Ki-67; e (4) identificare i tipi di errore prodotti dall'algoritmo e collegarli alle limitazioni dei sistemi basati su regole discusse nella Sezione~\ref{sec:ch3_symbolic}.

\textbf{Materiali Necessari:}
\begin{itemize}
    \item Un computer con QuPath (versione 0.4 o successiva) installato (disponibile gratuitamente da \url{https://qupath.github.io}).
    \item Un'immagine di vetrino intero di cancro al seno colorata con Ki-67 digitalizzata o una regione di interesse rappresentativa (fornita dal coordinatore del corso, o scaricabile dal repository di immagini campione di QuPath).
    \item La guida utente di QuPath (disponibile online).
\end{itemize}

\textbf{Procedura:}

\begin{enumerate}
    \item \textbf{Aprire l'immagine in QuPath.} Avviare QuPath e aprire la WSI Ki-67 o il tile di immagine fornito. Esaminare l'immagine a più ingrandimenti per familiarizzare con il pattern di colorazione: i nuclei DAB-positivi (marroni) sono cellule in proliferazione Ki-67-positive; i nuclei positivi all'ematossilina (blu) sono Ki-67-negativi.

    \item \textbf{Definire una regione di interesse (ROI).} Utilizzando lo strumento di annotazione rettangolo o poligono, disegnare una ROI che comprenda un'area rappresentativa del tumore invasivo, evitando necrosi, pieghe tissutali e aree di scarsa qualità della colorazione. Questo passaggio rispecchia la pratica del patologo di selezionare un ``hot spot'' o un'area rappresentativa per la valutazione del Ki-67.

    \item \textbf{Eseguire la stima del vettore di colorazione.} Navigare su \texttt{Analyze > Preprocessing > Estimate stain vectors}. QuPath stimerà automaticamente i vettori di colorazione dell'ematossilina e del DAB dall'immagine utilizzando un algoritmo di deconvoluzione del colore. Ispezionare i vettori stimati e regolarli se necessario per garantire una separazione accurata delle colorazioni.

    \item \textbf{Eseguire l'algoritmo di Rilevamento delle Cellule Positive.} Navigare su \texttt{Analyze > Cell detection > Positive cell detection}. Nella finestra di dialogo dei parametri, impostare quanto segue: \textit{Detection image}: ``Optical density sum''; \textit{Requested pixel size}: 0,5\,$\mu$m; \textit{Background radius}: 8\,$\mu$m; \textit{Median filter radius}: 0\,$\mu$m; \textit{Sigma}: 1,5\,$\mu$m; \textit{Minimum area}: 10\,$\mu$m$^2$; \textit{Maximum area}: 400\,$\mu$m$^2$. Per la soglia di positività, iniziare con una soglia OD media DAB di \textbf{0,2}. Fare clic su \textit{Run}.

    \item \textbf{Ispezionare i risultati.} QuPath sovrapporrà marcatori colorati su ogni nucleo rilevato: le cellule positive sono contrassegnate con un colore (tipicamente marrone o rosso) e le cellule negative con un altro (tipicamente blu). Aprire il pannello \textit{Annotations} per visualizzare le statistiche di riepilogo, incluso il conteggio totale delle cellule, il conteggio delle cellule positive e l'indice Ki-67 (espresso come percentuale).

    \item \textbf{Variare la soglia.} Ripetere l'analisi con soglie OD media DAB di \textbf{0,1} e \textbf{0,4}. Registrare l'indice Ki-67 ottenuto a ciascuna soglia. Notare come la classificazione dei nuclei borderline cambia al variare della soglia.

    \item \textbf{Identificare gli errori.} A ciascuna impostazione di soglia, ispezionare visivamente i risultati del rilevamento e identificare esempi di: (a) falsi positivi — cellule non tumorali (ad es., fibroblasti stromali, cellule endoteliali, linfociti) classificate come cellule tumorali Ki-67-positive; (b) falsi negativi — cellule tumorali Ki-67-positive mancate dall'algoritmo; (c) errori di segmentazione — nuclei che sono fusi (sotto-segmentati) o divisi (sovra-segmentati).
\end{enumerate}

\textbf{Domande di Discussione:}
\begin{enumerate}
    \item Come cambia l'indice Ki-67 al variare della soglia DAB? Quali sono le implicazioni cliniche di questa sensibilità?
    \item Quali tipi di errore sono più comuni a una soglia bassa rispetto a una soglia alta? Come si relazionano questi ai concetti di sensibilità e specificità?
    \item L'algoritmo non distingue tra cellule tumorali e cellule non tumorali. Come potrebbe questa limitazione influenzare la validità clinica dell'indice Ki-67 che produce? Come potrebbe essere affrontata questa limitazione utilizzando un approccio più sofisticato (ad es., basato sull'apprendimento automatico)?
    \item Confrontare l'indice Ki-67 prodotto dall'algoritmo con una stima visiva effettuata esaminando l'immagine manualmente. Quali sono le fonti di discrepanza?
    \item Riflettere sul ruolo del tecnico in questo flusso di lavoro. Quali decisioni richiedono il giudizio umano e quali sono le conseguenze degli errori in ogni punto decisionale?
\end{enumerate}

\textbf{Risultati Attesi e Interpretazione:} Gli studenti dovrebbero osservare che l'indice Ki-67 è sensibile alla scelta della soglia, con soglie più basse che producono indici più alti (più nuclei classificati come positivi) e soglie più alte che producono indici più bassi. La soglia ottimale non è universale ma dipende dal protocollo di colorazione, dallo scanner e dal tipo di tessuto. Questo esercizio illustra la limitazione fondamentale dei sistemi basati su regole: le loro prestazioni dipendono dalla corretta specificazione dei parametri che devono essere impostati da un esperto umano e potrebbero non generalizzarsi in diverse condizioni di laboratorio. Illustra anche il valore di tali sistemi per la quantificazione riproducibile e ad alto rendimento all'interno di un contesto definito e validato.

% ============================================================
\section{Sintesi}
\label{sec:ch3_summary}
% ============================================================

Questo capitolo ha fornito un'introduzione sistematica ai principi dell'intelligenza artificiale applicati alla patologia diagnostica. I seguenti punti chiave dovrebbero essere ricordati:

\begin{itemize}
    \item \textbf{L'IA è un campo ampio} che comprende qualsiasi sistema computazionale che esibisce un comportamento intelligente. In patologia clinica, tutte le attuali applicazioni di IA sono IA ristretta: altamente capaci all'interno di un dominio definito ma incapaci di generalizzare oltre di esso.

    \item \textbf{La tassonomia dell'IA} è gerarchica: l'apprendimento automatico è un sotto-campo dell'IA, e l'apprendimento profondo è un sotto-campo dell'apprendimento automatico. Comprendere questa gerarchia è essenziale per interpretare le affermazioni sugli strumenti diagnostici basati sull'IA.

    \item \textbf{L'IA simbolica e i sistemi basati su regole} codificano la conoscenza del dominio come regole logiche esplicite. Sono interpretabili, deterministici e computazionalmente efficienti, ma fragili e difficili da generalizzare. L'algoritmo basato su soglia Ki-67 in QuPath è un esempio paradigmatico.

    \item \textbf{L'apprendimento automatico} consente ai sistemi di apprendere dai dati piuttosto che da regole esplicite. L'apprendimento supervisionato (classificazione e regressione), l'apprendimento non supervisionato (clustering) e l'apprendimento semi-supervisionato hanno ciascuno applicazioni distinte in patologia. Gli algoritmi classici — regressione logistica, SVM, foreste casuali — rimangono ampiamente utilizzati. L'overfitting, la regolarizzazione e la validazione incrociata sono preoccupazioni metodologiche centrali.

    \item \textbf{Le reti bayesiane} forniscono un quadro rigoroso per il ragionamento diagnostico probabilistico, quantificando esplicitamente l'incertezza e incorporando la conoscenza clinica a priori. Sono particolarmente preziose per il supporto decisionale in scenari diagnostici complessi e multi-variabili.

    \item \textbf{L'apprendimento profondo} supera le limitazioni delle caratteristiche costruite manualmente apprendendo rappresentazioni gerarchiche direttamente dai dati grezzi delle immagini. Le reti neurali convoluzionali sono l'architettura dominante per l'analisi delle immagini istopatologiche, e i modelli fondazionali rappresentano la frontiera attuale del settore.

    \item \textbf{Il ruolo del tecnico} nella patologia assistita dall'IA è critico e multiforme: comprende l'annotazione dei dati, il controllo qualità dei dati di addestramento e degli output dell'IA, la validazione dei sistemi di IA e l'esercizio di un giudizio umano informato come salvaguardia contro gli errori dell'IA. La comprensione dei principi dell'IA è un prerequisito per svolgere questo ruolo in modo efficace.
\end{itemize}

% ============================================================
\section{Domande di Revisione}
\label{sec:ch3_questions}
% ============================================================

Le seguenti domande a risposta multipla sono progettate per consolidare la comprensione del materiale presentato in questo capitolo. Per ciascuna domanda, selezionare la singola risposta migliore.

\begin{enumerate}

    % Q1
    \item \label{ch3:q1} Quale delle seguenti affermazioni descrive meglio la relazione tra intelligenza artificiale, apprendimento automatico e apprendimento profondo?
    \begin{enumerate}[label=\Alph*.]
        \item Sono tre campi indipendenti e non correlati dell'informatica.
        \item L'apprendimento automatico è un sottoinsieme dell'apprendimento profondo, che è un sottoinsieme dell'intelligenza artificiale.
        \item L'apprendimento profondo è un sottoinsieme dell'apprendimento automatico, che è un sottoinsieme dell'intelligenza artificiale.
        \item L'intelligenza artificiale è un sottoinsieme dell'apprendimento automatico, che è un sottoinsieme dell'apprendimento profondo.
    \end{enumerate}

    % Q2
    \item \label{ch3:q2} Nel contesto dell'apprendimento automatico, quale delle seguenti definizioni descrive meglio un'\textit{etichetta}?
    \begin{enumerate}[label=\Alph*.]
        \item Una rappresentazione numerica dei dati di input utilizzata dal modello per effettuare previsioni.
        \item L'annotazione di riferimento assegnata a un esempio di addestramento da un esperto del dominio.
        \item La probabilità di output prodotta da un modello addestrato durante l'inferenza.
        \item Un parametro di regolarizzazione che controlla la complessità del modello.
    \end{enumerate}

    % Q3
    \item \label{ch3:q3} Un sistema di analisi delle immagini basato su regole per la valutazione del Ki-67 classifica un nucleo come positivo se la sua densità ottica media del DAB supera una soglia definita. Quale delle seguenti è la limitazione più significativa di questo approccio?
    \begin{enumerate}[label=\Alph*.]
        \item È computazionalmente troppo lento per l'uso routinario in laboratorio.
        \item Richiede un grande insieme di dati di addestramento etichettati per funzionare correttamente.
        \item La soglia fissa potrebbe non generalizzarsi tra diversi protocolli di colorazione o scanner.
        \item Non è in grado di distinguere tra colorazione DAB ed ematossilina.
    \end{enumerate}

    % Q4
    \item \label{ch3:q4} Nell'apprendimento automatico supervisionato, l'\textit{overfitting} si riferisce a quale dei seguenti fenomeni?
    \begin{enumerate}[label=\Alph*.]
        \item Il modello raggiunge un'elevata accuratezza sui dati di addestramento ma prestazioni scadenti su dati non visti.
        \item Il modello è troppo semplice per catturare il pattern sottostante nei dati di addestramento.
        \item L'insieme di dati di addestramento è troppo grande per le risorse computazionali disponibili.
        \item La funzione di perdita non converge durante il processo di addestramento.
    \end{enumerate}

    % Q5
    \item \label{ch3:q5} Il teorema di Bayes mette in relazione la probabilità a posteriori di un'ipotesi con quali delle seguenti quantità?
    \begin{enumerate}[label=\Alph*.]
        \item La probabilità a priori dell'ipotesi e la verosimiglianza dell'evidenza data l'ipotesi.
        \item La sensibilità e la specificità del test diagnostico.
        \item La media e la deviazione standard della distribuzione delle caratteristiche.
        \item Il numero di esempi di addestramento e il numero di parametri del modello.
    \end{enumerate}

    % Q6
    \item \label{ch3:q6} Quale delle seguenti descrive meglio la funzione di uno strato convoluzionale in una rete neurale convoluzionale?
    \begin{enumerate}[label=\Alph*.]
        \item Riduce le dimensioni spaziali delle mappe delle caratteristiche prendendo il valore massimo in ogni regione locale.
        \item Applica un filtro apprendibile sull'immagine di input per produrre una mappa delle caratteristiche che codifica la presenza di un particolare pattern.
        \item Calcola la somma pesata di tutte le caratteristiche di input e passa il risultato attraverso una funzione di attivazione.
        \item Disattiva casualmente una proporzione di neuroni durante l'addestramento per prevenire l'overfitting.
    \end{enumerate}

    % Q7
    \item \label{ch3:q7} Un classificatore a foresta casuale differisce da un singolo albero decisionale in quanto:
    \begin{enumerate}[label=\Alph*.]
        \item Utilizza una macchina a vettori di supporto come apprenditore base anziché un albero decisionale.
        \item Costruisce più alberi decisionali su sottoinsiemi casuali di dati e caratteristiche, e aggrega le loro previsioni.
        \item Applica la regolarizzazione L2 all'albero decisionale per prevenire l'overfitting.
        \item Richiede che le caratteristiche di input siano standardizzate a media zero e varianza unitaria.
    \end{enumerate}

    % Q8
    \item \label{ch3:q8} Nel contesto della validazione dei modelli di IA in patologia, lo \textit{spostamento della distribuzione} si riferisce a:
    \begin{enumerate}[label=\Alph*.]
        \item La trasformazione statistica applicata per normalizzare i valori di intensità dei pixel prima dell'addestramento del modello.
        \item Un cambiamento nelle proprietà statistiche dei dati di input tra le impostazioni di addestramento e di dispiegamento, che può degradare le prestazioni del modello.
        \item La ridistribuzione dei pesi del modello durante la backpropagation.
        \item Il processo di trasferimento di un modello addestrato su immagini naturali a un insieme di dati specifico per la patologia.
    \end{enumerate}

    % Q9
    \item \label{ch3:q9} Quale delle seguenti descrive meglio il ruolo del biologo o tecnico di laboratorio biomedico in un flusso di lavoro di quantificazione del Ki-67 assistito dall'IA?
    \begin{enumerate}[label=\Alph*.]
        \item Sostituire il patologo nell'interpretazione dell'indice Ki-67 e nelle decisioni terapeutiche.
        \item Eseguire l'annotazione dei dati, il controllo qualità degli output dell'IA e l'identificazione degli errori dell'algoritmo, agendo come elemento critico dell'uomo nel ciclo.
        \item Addestrare il modello di apprendimento profondo su nuovi insiemi di dati Ki-67 utilizzando risorse di calcolo GPU.
        \item Selezionare l'architettura di rete neurale più appropriata per il compito di quantificazione del Ki-67.
    \end{enumerate}

    % Q10
    \item \label{ch3:q10} Quale delle seguenti affermazioni sulle reti bayesiane è corretta?
    \begin{enumerate}[label=\Alph*.]
        \item Richiedono che tutte le variabili di input siano indipendenti l'una dall'altra.
        \item Rappresentano le variabili e le loro dipendenze condizionali come un grafo aciclico diretto, consentendo l'inferenza probabilistica in condizioni di incertezza.
        \item Sono un tipo di architettura di apprendimento profondo che utilizza meccanismi di attenzione per modellare le dipendenze a lungo raggio.
        \item Possono essere applicate solo a problemi di classificazione binaria con due possibili esiti diagnostici.
    \end{enumerate}

\end{enumerate}

\bigskip
\noindent\textit{Le risposte alle domande di revisione sono fornite nella Chiave delle Risposte alla fine del presente manuale.}

% ============================================================
%  Capitolo 4 – Apprendimento Profondo e Applicazioni Multimodali in Patologia
%  Postgraduate Pathology Course Handbook
%  This file is \input{} from the main document; do NOT add
%  \documentclass, \begin{document}, or \end{document} here.
% ============================================================

\chapter{Apprendimento Profondo e Applicazioni Multimodali in Patologia}
\label{ch:4}

% ============================================================
\section{Introduzione}
\label{sec:ch4_intro}
% ============================================================

L'ultimo decennio ha testimoniato una profonda trasformazione nella patologia computazionale, guidata in larga misura dalla rapida maturazione dell'apprendimento profondo---una famiglia di metodi di apprendimento automatico che apprendono rappresentazioni gerarchiche direttamente dai dati grezzi. A differenza delle pipeline classiche di analisi delle immagini, che richiedevano a esperti del dominio di costruire manualmente caratteristiche come descrittori della forma nucleare o istogrammi di colore, i modelli di apprendimento profondo scoprono autonomamente rappresentazioni discriminative da grandi dataset annotati. Questo cambiamento di paradigma ha consentito agli algoritmi di eguagliare o superare l'accuratezza diagnostica di patologi esperti su un numero crescente di compiti, dal rilevamento delle metastasi linfonodali nel cancro al seno alla classificazione dell'adenocarcinoma prostatico \cite{ehteshami2017,bulten2022}.

La convergenza di tre fattori abilitanti ha accelerato questo progresso. In primo luogo, l'adozione diffusa degli scanner per immagini di vetrini interi (WSI) ha generato enormi archivi di sezioni tissutali digitalizzate, fornendo i dati di addestramento richiesti dalle reti profonde. In secondo luogo, i progressi nell'hardware delle unità di elaborazione grafica (GPU) hanno reso computazionalmente fattibile addestrare modelli con centinaia di milioni di parametri in ore anziché settimane. In terzo luogo, il movimento per la scienza aperta ha prodotto dataset benchmark pubblicamente disponibili---come CAMELYON, PANDA e TCGA---che consentono un confronto rigoroso e riproducibile tra metodi concorrenti \cite{campanella2019,song2023}.

Questo capitolo fornisce un'introduzione sistematica alle architetture di apprendimento profondo che sono alla base della moderna patologia computazionale. Iniziamo con le reti neurali convoluzionali (CNN), il pilastro del riconoscimento delle immagini, e spieghiamo come i loro componenti architetturali---strati convoluzionali, pooling, batch normalisation e dropout---consentano un apprendimento robusto delle caratteristiche dalle immagini istopatologiche. Affrontiamo poi la sfida unica posta dalle WSI gigapixel e descriviamo come i framework di apprendimento multi-istanza (MIL) aggirino i vincoli di memoria GPU sfruttando solo la supervisione a livello di vetrino. Il capitolo prosegue con le architetture basate su transformer e i meccanismi di auto-attenzione, che hanno recentemente soppiantato le CNN come stato dell'arte in molti benchmark di patologia, e con i modelli linguistici di grandi dimensioni (LLM) che automatizzano la generazione di referti strutturati. Esaminiamo successivamente i modelli multimodali che fondono dati di immagini con informazioni genomiche, cliniche e radiologiche per produrre predizioni prognostiche più ricche. Il capitolo si conclude con una discussione sul ruolo del biologo nei flussi di lavoro assistiti dall'IA, due attività pratiche, una sintesi e dieci domande di revisione.

In tutto il capitolo, l'enfasi è sulla comprensione concettuale piuttosto che sulla derivazione matematica. I lettori che desiderano approfondire la matematica sottostante sono indirizzati alla letteratura primaria citata in ogni fase. Al termine di questo capitolo, gli studenti dovrebbero essere in grado di descrivere le principali architetture di apprendimento profondo utilizzate in patologia, spiegarne i punti di forza e le limitazioni, e articolare le considerazioni pratiche ed etiche che governano il loro impiego clinico.

% ============================================================
\section{Reti Neurali Convoluzionali per l'Analisi delle Immagini Istopatologiche}
\label{sec:ch4_cnn}
% ============================================================

\subsection{Architettura delle Reti Neurali Convoluzionali}

Una rete neurale convoluzionale è una rete neurale feed-forward specificamente progettata per elaborare dati con una topologia a griglia, come le immagini. La sua caratteristica distintiva è lo \emph{strato convoluzionale}, in cui un insieme di filtri apprendibili---chiamati anche \emph{kernel}---viene convolutato con l'input per produrre \emph{mappe delle caratteristiche}. Ogni kernel è una piccola matrice di pesi (tipicamente $3 \times 3$ o $5 \times 5$ pixel) che scorre sull'immagine di input con uno \emph{stride} definito (il numero di pixel di cui il kernel avanza ad ogni passo). Quando lo stride è uguale a uno, il kernel si sposta di un pixel alla volta, producendo una mappa delle caratteristiche di dimensioni spaziali approssimativamente uguali all'input; stride maggiori producono mappe delle caratteristiche più piccole e introducono un certo grado di sottocampionamento spaziale. Il \emph{padding}---l'aggiunta di zeri attorno al bordo dell'immagine---viene comunemente applicato per preservare le dimensioni spaziali e garantire che i pixel ai bordi ricevano la stessa copertura dei pixel centrali.

L'output di ogni operazione convoluzionale viene passato attraverso una \emph{funzione di attivazione} non lineare. L'Unità Lineare Rettificata (ReLU), definita come $f(x) = \max(0, x)$, è la funzione di attivazione più ampiamente utilizzata nelle CNN moderne. La ReLU introduce non-linearità evitando il problema del gradiente che svanisce, che affligge le attivazioni sigmoide e tangente iperbolica nelle reti profonde, poiché il suo gradiente è zero (per input negativi) o uno (per input positivi). Dopo le operazioni convoluzionali e di attivazione, gli \emph{strati di pooling} riducono le dimensioni spaziali delle mappe delle caratteristiche, diminuendo così il costo computazionale e conferendo un certo grado di invarianza alla traslazione. Il max-pooling, che conserva il valore massimo all'interno di ogni finestra di pooling, è la variante più comune; l'average-pooling, che calcola la media, è utilizzato anche in alcune architetture.

La \emph{batch normalisation} è una tecnica che normalizza le attivazioni di ogni strato su un mini-batch di esempi di addestramento, ricentrandole e ridimensionandole per avere media zero e varianza unitaria prima di applicare parametri di scala e traslazione apprendibili. La batch normalisation accelera l'addestramento, riduce la sensibilità alla scelta del tasso di apprendimento e agisce come un lieve regolarizzatore. Il \emph{dropout} è una strategia di regolarizzazione complementare in cui una frazione casuale di neuroni viene impostata a zero durante ogni passo di addestramento, prevenendo la co-adattazione delle caratteristiche e riducendo l'overfitting. Il dropout viene tipicamente applicato agli strati completamente connessi alla fine della rete, dove il rischio di overfitting è maggiore.

Gli strati finali di una CNN sono gli \emph{strati completamente connessi} (chiamati anche strati densi), in cui ogni neurone è connesso a ogni neurone dello strato precedente. Questi strati integrano le caratteristiche distribuite spazialmente estratte dallo stack convoluzionale in una rappresentazione globale adatta alla classificazione o alla regressione. Lo strato di output applica una funzione \emph{softmax} per la classificazione multi-classe, convertendo i punteggi grezzi (logit) in una distribuzione di probabilità sulle classi target. Uno schema dell'architettura CNN è mostrato nella Figura~\ref{fig:cnn_arch}.

\begin{figure}[htbp]
  \centering
  % figure4.tex  –  TikZ diagram: CNN architecture for histopathology
% Included via % figure4.tex  –  TikZ diagram: CNN architecture for histopathology
% Included via \input{../figures/figure4.tex} from chapter4.tex
% Requires: tikz, tikz libraries (shapes.geometric, arrows.meta, positioning, fit, backgrounds)

\begin{tikzpicture}[
    font=\small,
    >=Stealth,
    node distance=0.55cm and 0.9cm,
    % ---- block styles ----
    imgblock/.style={
        draw, fill=pink!40, rounded corners=2pt,
        minimum width=0.7cm, minimum height=1.6cm,
        align=center
    },
    convblock/.style={
        draw, fill=blue!25, rounded corners=2pt,
        minimum width=0.55cm, minimum height=1.4cm,
        align=center
    },
    poolblock/.style={
        draw, fill=green!30, rounded corners=2pt,
        minimum width=0.55cm, minimum height=1.0cm,
        align=center
    },
    fcblock/.style={
        draw, fill=orange!30, rounded corners=2pt,
        minimum width=0.55cm, minimum height=0.9cm,
        align=center
    },
    outblock/.style={
        draw, fill=red!25, rounded corners=4pt,
        minimum width=1.1cm, minimum height=0.7cm,
        align=center
    },
    lbl/.style={font=\scriptsize, align=center},
    arrow/.style={->, thick, gray!70}
]

%% ---- Input image ----
\node[imgblock] (img) {};
\node[lbl, below=0.15cm of img] {Input\\image\\patch};

%% ---- Conv block 1 (3 feature maps shown as stacked rectangles) ----
\foreach \i in {0,1,2}{
    \node[convblock, right=\i*0.12cm+0.9cm of img, yshift=\i*0.12cm] (conv1\i) {};
}
\node[lbl, below=0.15cm of conv10] {Conv 1\\(ReLU)};

%% ---- Pool 1 ----
\foreach \i in {0,1,2}{
    \node[poolblock, right=\i*0.12cm+0.55cm of conv10, yshift=\i*0.12cm] (pool1\i) {};
}
\node[lbl, below=0.15cm of pool10] {Pool 1\\(Max)};

%% ---- Conv block 2 ----
\foreach \i in {0,1,2,3}{
    \node[convblock, right=\i*0.12cm+0.55cm of pool10, yshift=\i*0.12cm] (conv2\i) {};
}
\node[lbl, below=0.15cm of conv20] {Conv 2\\(BN+ReLU)};

%% ---- Pool 2 ----
\foreach \i in {0,1,2,3}{
    \node[poolblock, right=\i*0.12cm+0.55cm of conv20, yshift=\i*0.12cm] (pool2\i) {};
}
\node[lbl, below=0.15cm of pool20] {Pool 2\\(Max)};

%% ---- Conv block 3 ----
\foreach \i in {0,...,4}{
    \node[convblock, right=\i*0.12cm+0.55cm of pool20, yshift=\i*0.12cm] (conv3\i) {};
}
\node[lbl, below=0.15cm of conv30] {Conv 3\\(BN+ReLU)};

%% ---- Flatten arrow label ----
\node[right=0.65cm of conv34, yshift=-0.3cm] (flat) {};

%% ---- FC layers ----
\node[fcblock, right=1.5cm of conv30] (fc1) {};
\node[lbl, below=0.15cm of fc1] {FC 1\\(Dropout)};

\node[fcblock, right=0.7cm of fc1] (fc2) {};
\node[lbl, below=0.15cm of fc2] {FC 2\\(Dropout)};

%% ---- Output ----
\node[outblock, right=0.7cm of fc2] (out) {Softmax\\output};
\node[lbl, below=0.15cm of out] {Class\\probs.};

%% ---- Arrows ----
\draw[arrow] (img.east)    -- (conv10.west);
\draw[arrow] (conv12.east) -- (pool10.west);
\draw[arrow] (pool12.east) -- (conv20.west);
\draw[arrow] (conv23.east) -- (pool20.west);
\draw[arrow] (pool23.east) -- (conv30.west);
\draw[arrow] (conv34.east) -- node[above, font=\scriptsize]{Flatten} (fc1.west);
\draw[arrow] (fc1.east)    -- (fc2.west);
\draw[arrow] (fc2.east)    -- (out.west);

%% ---- Brace annotations above ----
\draw[decorate, decoration={brace, amplitude=5pt, raise=4pt}, thick, blue!60]
    (conv10.north west) -- (pool12.north east)
    node[midway, above=8pt, font=\scriptsize, blue!80] {Feature extraction stage 1};

\draw[decorate, decoration={brace, amplitude=5pt, raise=4pt}, thick, blue!60]
    (conv20.north west) -- (pool23.north east)
    node[midway, above=12pt, font=\scriptsize, blue!80] {Feature extraction stage 2};

\draw[decorate, decoration={brace, amplitude=5pt, raise=3pt}, thick, blue!60]
    (conv30.north west) -- (conv34.north east)
    node[midway, above=14pt, font=\scriptsize, blue!80] {Deep feature maps};

\draw[decorate, decoration={brace, amplitude=5pt, raise=3pt}, thick, orange!70]
    (fc1.north west) -- (fc2.north east)
    node[midway, above=10pt, font=\scriptsize, orange!90] {Classification head};

\end{tikzpicture}
 from chapter4.tex
% Requires: tikz, tikz libraries (shapes.geometric, arrows.meta, positioning, fit, backgrounds)

\begin{tikzpicture}[
    font=\small,
    >=Stealth,
    node distance=0.55cm and 0.9cm,
    % ---- block styles ----
    imgblock/.style={
        draw, fill=pink!40, rounded corners=2pt,
        minimum width=0.7cm, minimum height=1.6cm,
        align=center
    },
    convblock/.style={
        draw, fill=blue!25, rounded corners=2pt,
        minimum width=0.55cm, minimum height=1.4cm,
        align=center
    },
    poolblock/.style={
        draw, fill=green!30, rounded corners=2pt,
        minimum width=0.55cm, minimum height=1.0cm,
        align=center
    },
    fcblock/.style={
        draw, fill=orange!30, rounded corners=2pt,
        minimum width=0.55cm, minimum height=0.9cm,
        align=center
    },
    outblock/.style={
        draw, fill=red!25, rounded corners=4pt,
        minimum width=1.1cm, minimum height=0.7cm,
        align=center
    },
    lbl/.style={font=\scriptsize, align=center},
    arrow/.style={->, thick, gray!70}
]

%% ---- Input image ----
\node[imgblock] (img) {};
\node[lbl, below=0.15cm of img] {Input\\image\\patch};

%% ---- Conv block 1 (3 feature maps shown as stacked rectangles) ----
\foreach \i in {0,1,2}{
    \node[convblock, right=\i*0.12cm+0.9cm of img, yshift=\i*0.12cm] (conv1\i) {};
}
\node[lbl, below=0.15cm of conv10] {Conv 1\\(ReLU)};

%% ---- Pool 1 ----
\foreach \i in {0,1,2}{
    \node[poolblock, right=\i*0.12cm+0.55cm of conv10, yshift=\i*0.12cm] (pool1\i) {};
}
\node[lbl, below=0.15cm of pool10] {Pool 1\\(Max)};

%% ---- Conv block 2 ----
\foreach \i in {0,1,2,3}{
    \node[convblock, right=\i*0.12cm+0.55cm of pool10, yshift=\i*0.12cm] (conv2\i) {};
}
\node[lbl, below=0.15cm of conv20] {Conv 2\\(BN+ReLU)};

%% ---- Pool 2 ----
\foreach \i in {0,1,2,3}{
    \node[poolblock, right=\i*0.12cm+0.55cm of conv20, yshift=\i*0.12cm] (pool2\i) {};
}
\node[lbl, below=0.15cm of pool20] {Pool 2\\(Max)};

%% ---- Conv block 3 ----
\foreach \i in {0,...,4}{
    \node[convblock, right=\i*0.12cm+0.55cm of pool20, yshift=\i*0.12cm] (conv3\i) {};
}
\node[lbl, below=0.15cm of conv30] {Conv 3\\(BN+ReLU)};

%% ---- Flatten arrow label ----
\node[right=0.65cm of conv34, yshift=-0.3cm] (flat) {};

%% ---- FC layers ----
\node[fcblock, right=1.5cm of conv30] (fc1) {};
\node[lbl, below=0.15cm of fc1] {FC 1\\(Dropout)};

\node[fcblock, right=0.7cm of fc1] (fc2) {};
\node[lbl, below=0.15cm of fc2] {FC 2\\(Dropout)};

%% ---- Output ----
\node[outblock, right=0.7cm of fc2] (out) {Softmax\\output};
\node[lbl, below=0.15cm of out] {Class\\probs.};

%% ---- Arrows ----
\draw[arrow] (img.east)    -- (conv10.west);
\draw[arrow] (conv12.east) -- (pool10.west);
\draw[arrow] (pool12.east) -- (conv20.west);
\draw[arrow] (conv23.east) -- (pool20.west);
\draw[arrow] (pool23.east) -- (conv30.west);
\draw[arrow] (conv34.east) -- node[above, font=\scriptsize]{Flatten} (fc1.west);
\draw[arrow] (fc1.east)    -- (fc2.west);
\draw[arrow] (fc2.east)    -- (out.west);

%% ---- Brace annotations above ----
\draw[decorate, decoration={brace, amplitude=5pt, raise=4pt}, thick, blue!60]
    (conv10.north west) -- (pool12.north east)
    node[midway, above=8pt, font=\scriptsize, blue!80] {Feature extraction stage 1};

\draw[decorate, decoration={brace, amplitude=5pt, raise=4pt}, thick, blue!60]
    (conv20.north west) -- (pool23.north east)
    node[midway, above=12pt, font=\scriptsize, blue!80] {Feature extraction stage 2};

\draw[decorate, decoration={brace, amplitude=5pt, raise=3pt}, thick, blue!60]
    (conv30.north west) -- (conv34.north east)
    node[midway, above=14pt, font=\scriptsize, blue!80] {Deep feature maps};

\draw[decorate, decoration={brace, amplitude=5pt, raise=3pt}, thick, orange!70]
    (fc1.north west) -- (fc2.north east)
    node[midway, above=10pt, font=\scriptsize, orange!90] {Classification head};

\end{tikzpicture}

  \caption{Architettura schematica di una rete neurale convoluzionale (CNN) applicata alla classificazione di immagini istopatologiche. Un patch dell'immagine di input viene elaborato attraverso successivi strati convoluzionali e di pooling che estraggono caratteristiche gerarchiche di crescente astrazione. Il vettore di caratteristiche risultante viene passato a strati completamente connessi e a uno strato di output softmax che produce le probabilità di classe. Conv: strato convoluzionale; Pool: strato di pooling; FC: strato completamente connesso.}
  \label{fig:cnn_arch}
\end{figure}

\subsection{Apprendimento Gerarchico delle Caratteristiche}

Un'intuizione fondamentale alla base del successo delle CNN è che esse apprendono rappresentazioni \emph{gerarchiche}. Negli strati iniziali, i kernel rispondono a caratteristiche di basso livello come bordi, gradienti di colore e texture semplici. Negli strati intermedi, le combinazioni di caratteristiche di basso livello danno origine a pattern più complessi---nuclei, lumi ghiandolari, fibre stromali. Negli strati più profondi, la rete codifica concetti semantici di alto livello come ``carcinoma invasivo'' o ``ghiandola benigna''. Questa gerarchia rispecchia il modo in cui un patologo costruisce una diagnosi: percependo prima le singole cellule, poi l'architettura tissutale, e infine integrando queste osservazioni in una categoria diagnostica. La natura gerarchica delle rappresentazioni CNN è stata confermata da tecniche di visualizzazione come la mappatura dell'attivazione di classe ponderata per gradiente (Grad-CAM), che evidenzia le regioni dell'immagine più influenti per una data predizione, fornendo un grado di interpretabilità essenziale per la fiducia clinica \cite{song2023}.

La profondità di una CNN---il numero di strati impilati---è un determinante critico della sua capacità rappresentativa. Architetture di riferimento come AlexNet (8 strati), VGG-16 (16 strati), ResNet-50 (50 strati) ed EfficientNet hanno progressivamente aumentato la profondità introducendo innovazioni architetturali---in particolare le \emph{connessioni residuali} in ResNet, che consentono ai gradienti di fluire direttamente attraverso connessioni di salto, permettendo così l'addestramento di reti molto profonde senza degradazione. Queste architetture, originariamente sviluppate per il riconoscimento di immagini naturali sul dataset ImageNet, sono state ampiamente adattate per l'istopatologia.

Il processo di apprendimento delle caratteristiche gerarchiche è guidato dalla \emph{backpropagation}---l'algoritmo con cui il gradiente della funzione di perdita rispetto a ciascun parametro della rete viene calcolato e utilizzato per aggiornare i pesi tramite discesa del gradiente. Durante l'addestramento, alla rete vengono presentati patch di immagini etichettati in mini-batch; le probabilità di classe predette vengono confrontate con le etichette vere utilizzando una funzione di perdita (tipicamente l'entropia incrociata), e il segnale di errore risultante viene propagato all'indietro attraverso la rete per aggiornare tutti i pesi. Nel corso di molte iterazioni di addestramento, i kernel in ogni strato convergono per rilevare le caratteristiche più discriminative per il compito target. Il numero di iterazioni di addestramento, il tasso di apprendimento e la dimensione del batch sono \emph{iperparametri} che devono essere attentamente regolati per ottenere buone prestazioni di generalizzazione.

\subsection{Transfer Learning e Fine-Tuning}

Addestrare una CNN profonda da zero richiede milioni di esempi etichettati e risorse computazionali sostanziali. In patologia, i grandi dataset annotati sono costosi da produrre perché l'annotazione da parte di esperti richiede molto tempo ed è soggetta a variabilità inter-osservatore. Il \emph{transfer learning} affronta questo vincolo inizializzando una CNN con pesi pre-addestrati su un grande dataset di uso generale---più comunemente ImageNet, che contiene 1,2 milioni di immagini naturali in 1.000 categorie---e adattando poi la rete al compito di patologia target. La logica è che le caratteristiche di basso e medio livello apprese da immagini naturali (bordi, texture, gradienti di colore) sono ampiamente trasferibili alle immagini istopatologiche, anche se i due domini differiscono sostanzialmente nell'aspetto.

Il \emph{fine-tuning} si riferisce al processo di continuazione dell'ottimizzazione basata sul gradiente sulla rete pre-addestrata utilizzando il dataset target. Due strategie sono comuni. Nell'\emph{estrazione di caratteristiche}, gli strati convoluzionali vengono congelati e solo gli strati completamente connessi finali vengono riaddestrati; questo è appropriato quando il dataset target è piccolo e i domini sorgente e target sono dissimili. Nel \emph{fine-tuning completo}, tutti gli strati vengono aggiornati con un piccolo tasso di apprendimento, consentendo alla rete di adattare le sue rappresentazioni al dominio target; questo è preferito quando il dataset target è sufficientemente grande. Strategie intermedie---congelare gli strati iniziali mentre si esegue il fine-tuning di quelli successivi---sono anch'esse ampiamente utilizzate. Il transfer learning ha dimostrato di migliorare sostanzialmente le prestazioni e ridurre i tempi di addestramento in istopatologia, in particolare per i tipi di tumore rari dove i dati annotati sono scarsi \cite{coudray2018}.

\subsection{Applicazioni in Patologia}

\paragraph{Classificazione del Cancro al Seno.}
Una delle prime e più clinicamente rilevanti applicazioni delle CNN in patologia è la classificazione dell'istologia del cancro al seno in carcinoma invasivo e carcinoma duttale \emph{in situ} (DCIS). Le CNN addestrate su sezioni colorate con ematossilina ed eosina (E\&E) hanno dimostrato prestazioni comparabili a quelle di patologi specialisti del seno in questo compito, con valori dell'area sotto la curva caratteristica operativa del ricevitore (AUC) superiori a 0,95 in diversi studi di validazione indipendenti. La capacità di automatizzare questa distinzione è clinicamente significativa perché influenza direttamente la gestione chirurgica: il carcinoma invasivo richiede la biopsia del linfonodo sentinella, mentre il DCIS no \cite{coudray2018}.

\paragraph{Rilevamento delle Metastasi Linfonodali: La Sfida CAMELYON.}
Le grandi sfide CAMELYON16 e CAMELYON17, organizzate sotto gli auspici del Gruppo di Analisi delle Immagini Diagnostiche del Centro Medico Universitario di Radboud, hanno fornito un benchmark di riferimento per il rilevamento basato su CNN delle metastasi del cancro al seno nelle sezioni dei linfonodi sentinella \cite{ehteshami2017}. I team partecipanti hanno addestrato CNN per classificare i patch di immagini come tumore o non-tumore, aggregando poi le predizioni a livello di patch per generare mappe di probabilità a livello di vetrino. L'algoritmo vincitore in CAMELYON16 ha raggiunto un AUC di 0,994 sul set di test, superando le prestazioni di un patologo che lavora sotto vincoli di tempo. L'analisi successiva ha rivelato che l'algoritmo era particolarmente abile nel rilevare cellule tumorali isolate e micrometastasi---lesioni facilmente trascurate nell'esame di routine---evidenziando il potenziale dell'IA come rete di sicurezza per i reperti a bassa prevalenza ma clinicamente importanti.

\paragraph{Classificazione di Gleason del Cancro alla Prostata: La Sfida PANDA.}
La sfida PANDA (Prostate cANcer graDe Assessment), riportata da Bulten e colleghi su \emph{Nature Medicine}, ha valutato gli algoritmi di apprendimento profondo per la classificazione di Gleason delle biopsie prostatiche con ago sottile \cite{bulten2022}. La classificazione di Gleason---l'assegnazione di un punteggio da 1 a 5 basato sull'architettura ghiandolare---è il fattore prognostico più importante nel cancro alla prostata, ma è soggetta a sostanziale variabilità inter-osservatore anche tra uropatologi esperti. La sfida PANDA ha dimostrato che gli algoritmi con le migliori prestazioni hanno raggiunto un accordo con uno standard di riferimento derivato da più patologi esperti comparabile all'accordo tra singoli patologi esperti, con valori di kappa ponderato quadratico superiori a 0,87. Questo risultato ha stabilito che l'apprendimento profondo può raggiungere prestazioni di livello clinico in un compito di classificazione complesso e multi-classe con conseguenze cliniche dirette.

% ============================================================
\section{Apprendimento Multi-Istanza per Immagini di Vetrini Interi}
\label{sec:ch4_mil}
% ============================================================

\subsection{La Sfida delle Immagini di Vetrini Interi Gigapixel}

Un'immagine di vetrino intero digitalizzata a ingrandimento $20\times$ contiene tipicamente tra $10^8$ e $10^9$ pixel, corrispondenti a una dimensione di file di diversi gigabyte in formato non compresso. Questa scala presenta una sfida computazionale fondamentale: nessuna GPU attuale dispone di memoria sufficiente per elaborare un'intera WSI come singolo tensore di input. La soluzione ingenua---sottocampionare l'immagine a una risoluzione compatibile con la memoria GPU---scarta il dettaglio a livello cellulare essenziale per molti compiti diagnostici. Un approccio alternativo, e quello più ampiamente adottato nel settore, è decomporre la WSI in un gran numero di \emph{patch} non sovrapposti o sovrapposti (chiamati anche \emph{tile}), tipicamente di dimensione $256 \times 256$ o $512 \times 512$ pixel, ed elaborare ogni patch indipendentemente attraverso un estrattore di caratteristiche CNN \cite{campanella2019}.

L'estrazione dei patch è preceduta dalla segmentazione del tessuto---l'identificazione delle regioni contenenti tessuto e l'esclusione del vetro di sfondo, degli artefatti e dei segni di penna. La segmentazione del tessuto viene comunemente eseguita utilizzando la sogliatura di Otsu nello spazio colore HSV o, sempre più spesso, utilizzando una CNN leggera addestrata a questo scopo. Il numero di patch estratti da una singola WSI può variare da poche centinaia a diverse decine di migliaia, a seconda dell'area del tessuto e della dimensione del patch. Ogni patch viene poi codificato in un vettore di caratteristiche compatto da un encoder CNN pre-addestrato, producendo un \emph{bag di caratteristiche} che rappresenta l'intero vetrino.

Una considerazione critica nell'analisi basata su patch è la scelta del livello di ingrandimento. Diverse caratteristiche diagnostiche sono meglio apprezzate a diversi ingrandimenti: la morfologia nucleare è più chiaramente visibile a $40\times$, l'architettura ghiandolare a $10\times$, e l'organizzazione a livello tissutale a $5\times$ o inferiore. Gli approcci multi-scala che estraggono patch a più livelli di ingrandimento e combinano le caratteristiche risultanti hanno dimostrato di superare gli approcci a scala singola nei compiti che richiedono l'integrazione di informazioni attraverso scale spaziali, come la classificazione di Gleason e la caratterizzazione del microambiente tumorale \cite{bulten2022}.

\subsection{Il Framework dell'Apprendimento Multi-Istanza}

L'\emph{apprendimento multi-istanza} (MIL) è un paradigma di apprendimento debolmente supervisionato idealmente adatto al problema dell'analisi WSI \cite{lu2020}. Nel framework MIL, ogni WSI viene trattata come un \emph{bag} di istanze (patch), e al bag viene assegnata una singola etichetta (ad esempio, tumore-positivo o tumore-negativo) senza specificare quali singole istanze siano responsabili dell'etichetta. Questo corrisponde precisamente alla realtà clinica: un patologo assegna una diagnosi all'intero vetrino, non ai singoli patch, e l'annotazione dei singoli patch su larga scala è proibitivamente costosa.

L'assunzione standard del MIL è che un bag sia positivo se e solo se almeno una delle sue istanze è positiva. Sotto questa assunzione, il problema di apprendimento consiste nell'addestrare un modello che possa classificare correttamente i bag utilizzando solo etichette a livello di bag. I primi approcci MIL aggregavano le predizioni a livello di patch utilizzando semplici operazioni di pooling---max-pooling (che assegna l'etichetta del bag in base al patch più sospetto) o mean-pooling (che media le predizioni su tutti i patch). Sebbene computazionalmente efficienti, questi approcci trattano i patch come indipendenti e ignorano le relazioni spaziali e contestuali tra di essi.

Il framework MIL è particolarmente potente perché consente l'\emph{apprendimento debolmente supervisionato} dagli archivi clinici di routine. Invece di richiedere l'annotazione esperta dei singoli patch---un processo che richiederebbe migliaia di ore per un grande dataset---i modelli MIL possono essere addestrati utilizzando solo le etichette diagnostiche registrate di routine nel sistema informativo di laboratorio. Questo riduce drasticamente l'onere di annotazione e consente la costruzione di dataset di addestramento di ordini di grandezza più grandi di quelli ottenibili con approcci completamente supervisionati. Campanella e colleghi hanno dimostrato questo principio su larga scala, addestrando un modello MIL su oltre 44.000 WSI del Memorial Sloan Kettering Cancer Center utilizzando solo etichette diagnostiche di routine \cite{campanella2019}.

\subsection{MIL Basato sull'Attenzione e CLAM}

Un significativo progresso nel MIL per la patologia è stato l'introduzione dell'\emph{aggregazione basata sull'attenzione} da parte di Lu e colleghi nel loro framework CLAM (Clustering-constrained Attention Multiple instance learning) \cite{lu2020}. Nel MIL basato sull'attenzione, una rete di attenzione apprendibile assegna un peso scalare a ogni patch, riflettendo la sua rilevanza stimata per la predizione a livello di vetrino. La rappresentazione del bag viene poi calcolata come la somma ponderata dei vettori di caratteristiche dei patch, con i patch ad alta attenzione che contribuiscono maggiormente alla predizione finale. Questo meccanismo è interpretabile: i pesi di attenzione possono essere visualizzati come una mappa di calore sovrapposta alla WSI, evidenziando le regioni che il modello considera più diagnosticamente informative.

CLAM estende il MIL basato sull'attenzione con due componenti aggiuntive. In primo luogo, impiega una perdita di clustering a livello di istanza che incoraggia i patch all'interno dello stesso cluster di attenzione ad avere rappresentazioni di caratteristiche simili, migliorando la qualità dei pesi di attenzione appresi. In secondo luogo, supporta la classificazione multi-classe attraverso una strategia uno-contro-tutti, consentendo la predizione simultanea di più categorie diagnostiche. CLAM è stato validato su grandi coorti del Cancer Genome Atlas (TCGA) e ha dimostrato forti prestazioni nei compiti di classificazione dei sottotipi attraverso più tipi di cancro, incluso adenocarcinoma polmonare versus carcinoma a cellule squamose, e sottotipi di carcinoma a cellule renali \cite{lu2020}.

\subsection{Applicazioni di Livello Clinico}

Il potenziale clinico dell'analisi WSI basata su MIL è stato dimostrato in una gamma di tipi di cancro e compiti diagnostici. Oltre al fondamentale studio di Campanella, i modelli MIL sono stati applicati al rilevamento dell'instabilità dei microsatelliti nel cancro colorettale, alla predizione dello stato di mutazione BRCA1/2 nel cancro al seno, e all'identificazione di alterazioni genomiche azionabili da immagini E\&E---compiti che tradizionalmente richiedono costosi test molecolari. Consentendo la profilazione molecolare virtuale da vetrini E\&E di routine, i modelli basati su MIL hanno il potenziale per triaggiare i pazienti per i test molecolari, riducendo i costi e i tempi di risposta \cite{song2023}.

Il paradigma MIL consente anche la \emph{predizione della sopravvivenza} e la \emph{classificazione dei sottotipi molecolari} direttamente da immagini E\&E, compiti che sarebbero impossibili con la sola supervisione a livello di patch. Aggregando le caratteristiche dei patch attraverso l'intero microambiente tumorale, i modelli MIL possono catturare caratteristiche globali dell'architettura tissutale---come il rapporto tra tumore e stroma, la densità dei linfociti infiltranti il tumore e l'organizzazione spaziale della necrosi---che sono predittive degli esiti clinici. Queste capacità posizionano il MIL come framework fondamentale per la prossima generazione di strumenti prognostici basati sull'IA in patologia.

% ============================================================
\section{Modelli Transformer e Auto-Attenzione in Patologia}
\label{sec:ch4_transformer}
% ============================================================

\subsection{L'Architettura Transformer}

L'architettura transformer, introdotta da Vaswani e colleghi nel 2017 per l'elaborazione del linguaggio naturale, è diventata da allora il paradigma dominante nell'apprendimento automatico attraverso le modalità. L'innovazione centrale del transformer è il \emph{meccanismo di auto-attenzione}, che consente a ogni elemento di una sequenza di input di prestare attenzione---ed essere influenzato da---ogni altro elemento, indipendentemente dalla loro distanza nella sequenza. Questo è in netto contrasto con le CNN, i cui kernel convoluzionali hanno un campo recettivo fisso e locale: un kernel $3 \times 3$ può osservare direttamente solo un vicinato $3 \times 3$, e le dipendenze a lungo raggio devono essere costruite gradualmente attraverso molti strati impilati.

Nel meccanismo di auto-attenzione, ogni elemento di input viene proiettato in tre vettori: una \emph{query} ($Q$), una \emph{chiave} ($K$) e un \emph{valore} ($V$). Il peso di attenzione tra due elementi viene calcolato come il prodotto scalare scalato dei loro vettori query e chiave, normalizzato da una funzione softmax. L'output per ogni elemento è la somma ponderata di tutti i vettori valore, dove i pesi sono i punteggi di attenzione. Formalmente:
\[
  \text{Attention}(Q, K, V) = \text{softmax}\!\left(\frac{QK^\top}{\sqrt{d_k}}\right)V,
\]
dove $d_k$ è la dimensionalità dei vettori chiave. Il fattore di scala $1/\sqrt{d_k}$ impedisce ai prodotti scalari di crescere in magnitudine, il che spingerebbe il softmax in regioni di gradiente molto piccolo.

L'\emph{attenzione multi-testa} estende questo meccanismo eseguendo $h$ operazioni di attenzione parallele (``teste'') con diverse proiezioni apprese, concatenando poi e proiettando linearmente gli output. Ogni testa può prestare attenzione a diversi aspetti dell'input---per esempio, una testa potrebbe catturare le relazioni di morfologia nucleare mentre un'altra cattura l'organizzazione stromale. La \emph{codifica posizionale} viene aggiunta agli embedding di input per iniettare informazioni sulla posizione spaziale di ogni elemento, poiché il meccanismo di auto-attenzione è intrinsecamente invariante alla permutazione e sarebbe altrimenti cieco all'ordine spaziale. Nel transformer originale vengono utilizzate codifiche posizionali sinusoidali; nei Vision Transformer, gli embedding posizionali apprendibili sono più comuni.

\subsection{Vision Transformer}

Il \emph{Vision Transformer} (ViT), introdotto da Dosovitskiy e colleghi, adatta l'architettura transformer ai dati di immagini dividendo un'immagine in una griglia di patch non sovrapposti (tipicamente $16 \times 16$ pixel), incorporando linearmente ogni patch in un vettore ed elaborando la sequenza risultante di embedding di patch con un encoder transformer standard. Un token \texttt{[CLS]} apprendibile viene preposto alla sequenza; la sua rappresentazione finale viene utilizzata per la classificazione. Il ViT ha dimostrato che, con dati di addestramento sufficienti, un'architettura transformer pura può eguagliare o superare le prestazioni delle CNN sui benchmark di classificazione delle immagini, senza alcun bias induttivo convoluzionale.

Il vantaggio chiave del ViT rispetto alla CNN per la patologia è il suo \emph{campo recettivo globale}: fin dal primo strato, ogni patch può prestare attenzione a ogni altro patch nell'immagine. Questo è particolarmente prezioso in istopatologia, dove le caratteristiche diagnostiche spesso coinvolgono relazioni spaziali a lungo raggio---per esempio, la distribuzione spaziale dei linfociti infiltranti il tumore rispetto al fronte invasivo, o la relazione tra architettura ghiandolare e desmoplasia stromale. Studi comparativi hanno confermato che i modelli basati su ViT superano i modelli basati su CNN su diversi benchmark di patologia, in particolare quando pre-addestrati su grandi dataset specifici per la patologia \cite{deininger2022}.

Una limitazione pratica del ViT è la sua complessità computazionale quadratica rispetto alla lunghezza della sequenza: l'operazione di auto-attenzione richiede il calcolo dei punteggi di attenzione a coppie tra tutte le coppie di patch, il che diventa proibitivamente costoso per sequenze molto lunghe. Per l'analisi a livello di WSI, dove il numero di patch può raggiungere decine di migliaia, questa limitazione è significativa. Le architetture transformer gerarchiche---come Swin Transformer, che calcola l'attenzione all'interno di finestre locali e sposta le finestre tra gli strati---affrontano questa limitazione riducendo la lunghezza effettiva della sequenza ad ogni strato di attenzione, consentendo comunque il flusso di informazioni a lungo raggio attraverso la struttura gerarchica \cite{xu2024}.

\subsection{Modelli Fondazionali Specifici per la Patologia}

Un \emph{modello fondazionale} è una grande rete neurale pre-addestrata su un dataset ampio e diversificato utilizzando l'apprendimento auto-supervisionato, che può poi essere sottoposta a fine-tuning o utilizzata come estrattore di caratteristiche per un'ampia gamma di compiti a valle. Il concetto, originariamente sviluppato nell'elaborazione del linguaggio naturale (ad esempio, BERT, GPT), è stato applicato con grande successo all'imaging patologico.

\textbf{UNI} è un modello fondazionale visivo pre-addestrato su oltre 100.000 WSI di diversi tipi di tessuto e protocolli di colorazione utilizzando il framework di apprendimento auto-supervisionato DINOv2. UNI produce embedding di caratteristiche a livello di patch che si generalizzano su un'ampia gamma di compiti a valle---inclusa la classificazione dei sottotipi di cancro, la predizione della sopravvivenza e la predizione delle mutazioni---senza fine-tuning specifico per il compito, raggiungendo prestazioni all'avanguardia su più benchmark \cite{wang2024}.

\textbf{CONCH} (CONtrastive learning from Captions for Histopathology) è un modello fondazionale visione-linguaggio addestrato su dati immagine-testo accoppiati provenienti da referti patologici e risorse educative. Allineando le rappresentazioni visive e testuali in uno spazio di embedding condiviso, CONCH consente la classificazione zero-shot utilizzando prompt in linguaggio naturale (ad esempio, ``carcinoma duttale invasivo, grado 3'') e supporta compiti di recupero cross-modale \cite{wang2024}.

\textbf{Prov-GigaPath} è un modello fondazionale per vetrini interi sviluppato da Microsoft Research e Providence Health, pre-addestrato su oltre 1,3 miliardi di tile di immagini patologiche provenienti da 171.189 WSI \cite{xu2024}. Prov-GigaPath impiega un'architettura gerarchica che elabora i tile a livello di patch utilizzando un encoder ViT e poi aggrega le rappresentazioni dei patch a livello di vetrino utilizzando un transformer basato su LongNet, consentendogli di modellare le dipendenze a lungo raggio attraverso l'intero vetrino. Prov-GigaPath ha raggiunto prestazioni all'avanguardia in 26 su 39 compiti di sottotipizzazione del cancro e pathomics, dimostrando il valore della scala nel pre-addestramento dei modelli fondazionali per la patologia.

Questi modelli fondazionali rappresentano un cambiamento di paradigma nella patologia computazionale: invece di addestrare modelli specifici per il compito da zero, i professionisti possono sfruttare potenti rappresentazioni pre-addestrate e adattarle a nuovi compiti con dati etichettati minimi---una capacità di particolare valore nei tipi di tumore rari e negli ambienti con risorse limitate \cite{song2023}.

% ============================================================
\section{Modelli Linguistici di Grandi Dimensioni per la Generazione di Referti Patologici}
\label{sec:ch4_llm}
% ============================================================

\subsection{Architettura e Capacità dei Modelli Linguistici di Grandi Dimensioni}

I modelli linguistici di grandi dimensioni (LLM) sono reti neurali basate su transformer addestrate su vasti corpus di testo---tipicamente centinaia di miliardi di token tratti da libri, letteratura scientifica, cartelle cliniche e internet---utilizzando un obiettivo auto-supervisionato come la predizione del token successivo (come nei modelli in stile GPT) o la modellazione del linguaggio mascherato (come nei modelli in stile BERT). La scala di questi modelli---che va da centinaia di milioni a centinaia di miliardi di parametri---li dota di capacità straordinarie: possono generare testo fluente e contestualmente appropriato, rispondere a domande, riassumere documenti, tradurre tra lingue ed eseguire ragionamenti in più fasi, tutto senza addestramento specifico per il compito.

Nei modelli in stile GPT (generative pre-trained transformer), la rete viene addestrata a predire il token successivo in una sequenza dati tutti i token precedenti. Durante l'inferenza, il testo viene generato in modo autoregressivo: il modello produce un token alla volta, aggiungendo ogni token generato al contesto di input prima di predire il successivo. La qualità del testo generato è controllata da un parametro di \emph{temperatura} che modula la nitidezza della distribuzione di probabilità dell'output: temperature basse producono output più deterministici e conservativi, mentre temperature elevate introducono maggiore diversità e creatività. Il \emph{fine-tuning con istruzioni} e l'\emph{apprendimento per rinforzo dal feedback umano} (RLHF) vengono utilizzati per allineare il comportamento degli LLM alle preferenze umane, producendo modelli che seguono le istruzioni in modo affidabile ed evitano output dannosi.

L'architettura transformer che è alla base degli LLM è identica in linea di principio a quella descritta nella sezione precedente, con la differenza fondamentale che gli LLM operano su sequenze di token di testo piuttosto che su patch di immagini. Ogni token viene mappato in un vettore di embedding ad alta dimensione, e l'encoder o decoder transformer elabora la sequenza di embedding utilizzando l'auto-attenzione multi-testa e strati feed-forward. L'enorme scala dell'addestramento degli LLM---sia in termini di parametri del modello che di dati di addestramento---consente \emph{capacità emergenti} non presenti nei modelli più piccoli, tra cui l'apprendimento nel contesto (la capacità di eseguire nuovi compiti da pochi esempi forniti nel prompt) e il ragionamento a catena di pensiero (la capacità di risolvere problemi complessi generando passaggi di ragionamento intermedi).

\subsection{Applicazioni nella Generazione di Referti Patologici}

I referti patologici sono documenti strutturati che comunicano i risultati diagnostici, le informazioni prognostiche e le raccomandazioni cliniche ai medici curanti. La generazione di referti di alta qualità, completi e formattati in modo coerente è una componente che richiede molto tempo del carico di lavoro del patologo, e gli errori o le omissioni nei referti possono avere conseguenze dirette sulla sicurezza del paziente. Gli LLM offrono il potenziale per automatizzare o semi-automatizzare diversi aspetti della generazione di referti, inclusa la trascrizione dei risultati dettati, la compilazione di modelli di referto \emph{sinottico} (strutturato, basato su checklist) e la sintesi di riassunti narrativi da dati strutturati.

La \emph{refertazione sinottica}---l'uso di modelli standardizzati e itemizzati che catturano tutti gli elementi di dati richiesti per un dato tipo di campione---è stata sostenuta da organismi professionali come il Royal College of Pathologists e il College of American Pathologists come mezzo per migliorare la completezza e la coerenza dei referti. Gli LLM possono essere sottoposti a fine-tuning per estrarre elementi di dati strutturati da dettature in testo libero e compilare automaticamente i modelli sinotici, riducendo lo sforzo manuale richiesto e diminuendo il tasso di elementi di dati mancanti. Gli studi hanno dimostrato che gli LLM possono raggiungere un'elevata accuratezza nei compiti di estrazione di informazioni dai referti patologici, inclusi tipo di tumore, grado, stadio, stato dei margini e invasione linfovascolare \cite{bankhead2022}.

L'\emph{elaborazione del linguaggio naturale} (NLP) degli archivi di referti patologici esistenti consente studi retrospettivi su larga scala che sarebbero impraticabili con l'estrazione manuale dei dati. Gli LLM possono essere utilizzati per identificare coorti di pazienti con diagnosi specifiche, estrarre elementi di dati quantitativi e segnalare referti che richiedono una revisione---per esempio, referti che menzionano un reperto inatteso o che mancano di un elemento di dati richiesto. Questa capacità è sempre più importante per la garanzia della qualità, l'audit e la ricerca nei dipartimenti di patologia. Inoltre, gli LLM possono assistere nell'armonizzazione della terminologia tra le istituzioni, mappando i termini diagnostici locali su ontologie standardizzate come SNOMED-CT, facilitando così la condivisione dei dati e la ricerca multicentrica.

\subsection{LLM Multimodali e Modelli Visione-Linguaggio}

Un'estensione naturale degli LLM alla patologia è lo sviluppo di \emph{modelli visione-linguaggio} (VLM) che possono elaborare sia immagini che testo. Modelli come GPT-4V, LLaVA e varianti specifiche per la patologia possono accettare un'immagine istopatologica come input insieme a un prompt testuale e generare una risposta testuale---per esempio, descrivendo le caratteristiche morfologiche visibili nell'immagine, suggerendo una diagnosi differenziale o rispondendo a una specifica domanda clinica. Questi modelli vengono addestrati allineando le rappresentazioni visive e testuali, tipicamente utilizzando obiettivi di apprendimento contrastivo simili a quelli utilizzati in CONCH \cite{wang2024}. Il modello CONCH, per esempio, può essere sollecitato con una descrizione in linguaggio naturale di un tipo di tessuto e recupererà i patch visivamente più simili da un grande database, consentendo la ricerca basata su immagini e la classificazione zero-shot.

\subsection{Limitazioni e Necessità di Validazione}

Nonostante le loro impressionanti capacità, gli LLM presentano diverse importanti limitazioni che devono essere attentamente considerate prima dell'impiego clinico. La più significativa è l'\emph{allucinazione}: la tendenza degli LLM a generare affermazioni plausibili ma fattualmente errate. In un contesto clinico, una diagnosi allucinata, un grado tumorale errato o un riferimento fabbricato a una linea guida clinica potrebbero avere gravi conseguenze per la sicurezza del paziente. L'allucinazione si verifica perché gli LLM sono addestrati a produrre testo fluente e contestualmente coerente, non a verificare l'accuratezza fattuale; non hanno alcun meccanismo per distinguere tra informazioni che hanno appreso in modo affidabile e informazioni che stanno confabulando.

Le limitazioni aggiuntive includono il \emph{limite di conoscenza}---gli LLM sono addestrati su dataset statici e potrebbero non riflettere i criteri diagnostici o le linee guida cliniche più recenti---e lo \emph{spostamento della distribuzione}, per cui un modello addestrato su referti di un'istituzione potrebbe avere prestazioni scadenti su referti di un'altra a causa di differenze nella terminologia, nella formattazione e nella pratica clinica. Il \emph{bias} è anche una preoccupazione: se il corpus di addestramento sovra-rappresenta certi gruppi demografici, tipi di tumore o contesti clinici, le prestazioni del modello potrebbero essere diseguali tra le popolazioni di pazienti. Per queste ragioni, una rigorosa validazione prospettica nell'ambiente clinico target è un prerequisito essenziale per l'impiego degli LLM in patologia, e la supervisione umana deve essere mantenuta in tutte le fasi del flusso di lavoro di refertazione \cite{song2023}.

% ============================================================
\section{Modelli Multimodali: Combinazione di Immagini e Dati Clinici}
\label{sec:ch4_multimodal}
% ============================================================

\subsection{Razionale per l'Integrazione Multimodale}

Una limitazione fondamentale dei modelli di IA basati solo su immagini è che accedono solo a una frazione delle informazioni disponibili per il processo decisionale clinico. Nella pratica oncologica di routine, il referto del patologo viene interpretato nel contesto della storia clinica, dei risultati radiologici, degli esami di laboratorio e---sempre più---dei dati di profilazione molecolare. I modelli di IA multimodali che integrano informazioni da più modalità di dati hanno il potenziale per produrre predizioni più accurate e clinicamente azionabili rispetto ai modelli unimodali, sfruttando informazioni complementari non disponibili in nessuna singola modalità \cite{song2023}.

L'integrazione delle caratteristiche delle immagini istopatologiche con i dati genomici è di particolare interesse a causa della relazione ben consolidata tra morfologia tumorale e biologia molecolare. La morfologia tumorale riflette le alterazioni genetiche ed epigenetiche sottostanti che guidano l'oncogenesi, e viceversa, i profili molecolari influenzano l'aspetto istologico dei tumori. I modelli multimodali che modellano congiuntamente queste relazioni possono, in linea di principio, predire le caratteristiche molecolari dalla morfologia (consentendo la profilazione molecolare virtuale), predire gli esiti clinici con maggiore accuratezza rispetto a entrambe le modalità singolarmente, e identificare nuovi correlati morfologici delle alterazioni molecolari che potrebbero non essere evidenti all'occhio umano.

\subsection{Predizione dei Sottotipi Molecolari da Immagini E\&E}

Diversi studi hanno dimostrato che i modelli di apprendimento profondo addestrati su immagini E\&E possono predire le caratteristiche molecolari dei tumori con un'accuratezza clinicamente utile. Coudray e colleghi hanno dimostrato che una CNN addestrata su immagini E\&E di adenocarcinoma polmonare da TCGA poteva predire lo stato di mutazione di sei geni---inclusi \emph{STK11}, \emph{EGFR} e \emph{KRAS}---con valori AUC compresi tra 0,73 e 0,85 \cite{coudray2018}. Questo risultato suggerisce che le conseguenze morfologiche di queste mutazioni sono rilevabili dall'apprendimento profondo, anche se non sono immediatamente evidenti all'osservatore umano. Risultati simili sono stati riportati per la predizione dell'instabilità dei microsatelliti (MSI) da immagini E\&E di cancro colorettale, con implicazioni per la valutazione dell'idoneità all'immunoterapia.

La predizione dei sottotipi molecolari---come i sottotipi intrinseci PAM50 del cancro al seno (Luminale A, Luminale B, HER2-arricchito, Basale) o i sottotipi molecolari di consenso del cancro colorettale---da immagini E\&E è un compito più complesso che richiede al modello di integrare più caratteristiche morfologiche attraverso l'intero tumore. I modelli basati su MIL con meccanismi di attenzione sono ben adatti a questo compito, poiché possono aggregare le caratteristiche a livello di patch da tutta la WSI concentrando l'attenzione sulle regioni più informative. I modelli multimodali che combinano le caratteristiche E\&E con variabili cliniche (età, stadio, storia del trattamento) superano costantemente i modelli basati solo su immagini nei compiti di predizione dei sottotipi \cite{lu2020}.

\subsection{Predizione della Sopravvivenza Combinando Istologia e Genomica}

La predizione della sopravvivenza---la stima del tempo alla morte o alla recidiva della malattia di un paziente dai dati clinici e patologici di base---è una delle applicazioni clinicamente più importanti dell'IA multimodale in oncologia. I modelli unimodali basati sull'istologia o sulla genomica da soli hanno dimostrato valore prognostico, ma i modelli multimodali che combinano entrambi i tipi di dati raggiungono costantemente prestazioni superiori. Gli studi che integrano le caratteristiche WSI (estratte utilizzando un framework MIL) con i dati di sequenziamento dell'RNA in bulk hanno raggiunto valori dell'indice di concordanza (C-index) significativamente migliori per la predizione della sopravvivenza globale attraverso più tipi di cancro TCGA rispetto a entrambi i modelli unimodali \cite{song2023}.

La logica biologica di questo miglioramento è che l'istologia e la genomica catturano aspetti complementari della biologia tumorale: l'istologia riflette l'organizzazione spaziale del microambiente tumorale---inclusa la densità e la distribuzione dei linfociti infiltranti il tumore, l'entità della necrosi e il grado di desmoplasia stromale---mentre la genomica cattura i driver molecolari della crescita tumorale e della resistenza al trattamento. Modellando congiuntamente questi segnali complementari, i modelli multimodali possono identificare sottogruppi di pazienti con prognosi distinte che non sono separabili da nessuna delle due modalità singolarmente. L'integrazione dei dati di imaging radiologico---come le caratteristiche CT o MRI estratte da un modello di apprendimento profondo separato---come terza modalità ha dimostrato di fornire informazioni prognostiche aggiuntive oltre all'istologia e alla genomica in diversi tipi di cancro.

\subsection{Strategie di Fusione delle Caratteristiche}

L'integrazione delle caratteristiche da più modalità può essere ottenuta attraverso tre strategie principali, che differiscono nella fase in cui le rappresentazioni specifiche della modalità vengono combinate.

La \emph{fusione precoce} (chiamata anche \emph{fusione a livello di dati}) concatena le caratteristiche grezze o minimamente elaborate da tutte le modalità in un singolo vettore di input prima di qualsiasi elaborazione specifica della modalità. Questo approccio è semplice da implementare ma richiede che tutte le modalità siano rappresentate in uno spazio di caratteristiche comune, il che potrebbe non essere semplice quando le modalità differiscono sostanzialmente in dimensionalità e struttura (ad esempio, un'immagine gigapixel rispetto a un vettore di espressione genica a 20.000 dimensioni).

La \emph{fusione tardiva} (chiamata anche \emph{fusione a livello di decisione}) addestra modelli separati per ogni modalità e combina le loro predizioni---tipicamente tramite media, votazione o una combinazione appresa---nella fase di output. La fusione tardiva è robusta alle modalità mancanti (un modello può ancora produrre una predizione se una modalità non è disponibile) e consente alle architetture specifiche della modalità di essere ottimizzate indipendentemente. Tuttavia, non consente al modello di apprendere le interazioni cross-modali.

La \emph{fusione intermedia} (chiamata anche \emph{fusione a livello di caratteristiche}) combina le rappresentazioni specifiche della modalità in uno strato intermedio della rete, consentendo al modello di apprendere le interazioni cross-modali beneficiando comunque dell'estrazione di caratteristiche specifica della modalità. I meccanismi basati sull'attenzione---come l'attenzione cross-modale, in cui i vettori query di una modalità prestano attenzione alle coppie chiave-valore di un'altra---sono particolarmente efficaci per la fusione intermedia e sono stati adottati in diversi modelli di patologia multimodale all'avanguardia \cite{song2023,xu2024}. La scelta della strategia di fusione dovrebbe essere guidata dalla dimensione del dataset disponibile, dal grado di complementarità delle modalità e dal requisito clinico di robustezza ai dati mancanti.

% ============================================================
\section{Supporto Tecnico nei Flussi di Lavoro Assistiti dall'IA}
\label{sec:ch4_technical}
% ============================================================

\subsection{Il Ruolo del Biologo nelle Pipeline di IA}

L'impiego di strumenti di IA in un laboratorio di patologia clinica non diminuisce l'importanza del biologo (BMS); piuttosto, trasforma ed estende il ruolo del BMS in modi che richiedono nuove competenze tecniche. Man mano che i sistemi di IA passano da prototipi di ricerca a strumenti clinici validati, il BMS è sempre più chiamato ad agire come interfaccia tra il sistema di IA e il flusso di lavoro clinico---garantendo che i dati siano preparati correttamente, che gli output dell'IA siano interpretati in modo appropriato e che il patologo riceva le informazioni necessarie per prendere una decisione diagnostica sicura e informata \cite{bankhead2022}.

Il BMS deve sviluppare una comprensione operativa dei sistemi di IA impiegati nel proprio laboratorio, inclusi i compiti per cui sono progettati, le popolazioni su cui sono stati validati, le loro caratteristiche di prestazione note (sensibilità, specificità, AUC) e le loro modalità di fallimento note. Questa comprensione è essenziale per l'interpretazione appropriata degli output dell'IA e per l'identificazione dei casi in cui l'output dell'IA dovrebbe essere trattato con cautela. Non richiede che il BMS abbia competenze in apprendimento automatico o ingegneria del software, ma richiede l'impegno con la documentazione tecnica e i rapporti di validazione forniti dallo sviluppatore del sistema di IA.

\subsection{Preparazione dei Dati per le Pipeline di IA}

La qualità degli output dell'IA dipende fondamentalmente dalla qualità dei dati di input. Per i sistemi di IA basati su WSI, ciò significa che i vetrini devono essere scansionati al corretto ingrandimento, con la messa a fuoco appropriata e senza artefatti come pieghe del tessuto, bolle d'aria o eterogeneità di colorazione. Il BMS è responsabile dell'implementazione e del monitoraggio delle procedure di controllo della qualità nelle fasi pre-analitiche e analitiche---incluso il processamento del tessuto, il sezionamento, la colorazione e la scansione---per garantire che i vetrini soddisfino gli standard di qualità richiesti dal sistema di IA. Molti sistemi di IA includono moduli di controllo della qualità automatizzati che segnalano i vetrini con area tissutale insufficiente, regioni fuori fuoco o artefatti di colorazione; il BMS deve comprendere i criteri utilizzati da questi moduli e adottare misure correttive quando i vetrini vengono rifiutati.

La preparazione dei dati comprende anche la gestione dei metadati---identificatori del paziente, tipo di campione, indicazione clinica, protocollo di colorazione---che devono essere correttamente associati a ogni WSI nel sistema informativo di laboratorio (LIS). Gli errori nei metadati possono portare a una classificazione errata dei campioni, all'applicazione errata dei modelli di IA (ad esempio, applicare un modello di classificazione del cancro alla prostata a una biopsia del seno) o al mancato recupero dei risultati dell'IA al momento della refertazione. Il BMS deve quindi mantenere pratiche rigorose di governance dei dati e garantire che il LIS sia correttamente configurato per supportare i flussi di lavoro assistiti dall'IA.

\subsection{Esecuzione dell'Inferenza e Interpretazione degli Output dell'IA}

Una volta che un vetrino è stato scansionato e sottoposto a controllo di qualità, il BMS può essere responsabile dell'avvio della pipeline di inferenza dell'IA---inviando la WSI al sistema di IA, monitorando il progresso dell'analisi e recuperando i risultati. In alcune configurazioni di laboratorio, questo processo è completamente automatizzato e integrato nel LIS; in altri, richiede un intervento manuale. Il BMS deve comprendere gli output prodotti dal sistema di IA---inclusi i punteggi di probabilità, le mappe di calore di attenzione e gli elementi di dati strutturati---ed essere in grado di comunicare questi output chiaramente al patologo refertante.

La \emph{segnalazione dei casi discordanti} è una responsabilità particolarmente importante. Un caso discordante è quello in cui la predizione dell'IA è incoerente con la valutazione clinica o morfologica preliminare---per esempio, un vetrino di un paziente con un forte sospetto clinico di malignità che l'IA classifica come benigno, o viceversa. Il BMS dovrebbe essere formato per riconoscere tali discordanze e per segnalarle al patologo per la revisione, piuttosto che accettare acriticamente l'output dell'IA. Ciò richiede una comprensione di base delle caratteristiche di prestazione del sistema di IA---la sua sensibilità, specificità e modalità di fallimento note---nonché il contesto clinico di ogni caso.

\subsection{Tracce di Audit e Governance}

I framework di governance clinica richiedono che tutte le decisioni diagnostiche assistite dall'IA siano documentate in modo da supportare l'audit, la revisione e la responsabilità. Il BMS è responsabile del mantenimento di registrazioni accurate su quale sistema di IA è stato utilizzato per ogni caso, la versione del modello, la data e l'ora dell'analisi, l'output dell'IA e la decisione finale del patologo. Questi registri devono essere conservati in modo sicuro, in conformità con la legislazione sulla protezione dei dati, e devono essere accessibili per l'audit retrospettivo. In caso di errore diagnostico, la traccia di audit consente di indagare se il sistema di IA ha contribuito all'errore e se l'errore era prevedibile date le caratteristiche di prestazione note del sistema.

Il BMS svolge anche un ruolo nel monitoraggio continuo delle prestazioni del sistema di IA nell'ambiente di impiego. I modelli di IA possono degradare nel tempo a causa dello \emph{spostamento della distribuzione}---cambiamenti nella popolazione di pazienti, nei protocolli di colorazione o nelle apparecchiature di scansione che causano la divergenza della distribuzione dei dati di input dalla distribuzione di addestramento. Il monitoraggio regolare delle prestazioni, utilizzando un campione rappresentativo di casi con esiti noti, è essenziale per rilevare tale degradazione e per attivare il riaddestramento o la ricalibrazione del modello quando necessario \cite{song2023}. Il BMS dovrebbe avere familiarità con le metriche di prestazione chiave utilizzate per monitorare i sistemi di IA e dovrebbe partecipare a riunioni di revisione regolari in cui vengono discussi i dati di monitoraggio.

% ============================================================
\section{Attività Pratiche}
\label{sec:ch4_practical}
% ============================================================

\subsection{Attività A: Classificazione Tissutale Basata su CNN in QuPath}

\textbf{Obiettivo:} Acquisire esperienza pratica con una rete neurale convoluzionale per la classificazione automatizzata del tessuto in un'immagine di vetrino intero, utilizzando la piattaforma di patologia digitale open-source QuPath \cite{bankhead2022}.

\textbf{Contesto:} QuPath fornisce un'interfaccia intuitiva per l'addestramento e l'applicazione di classificatori di pixel e classificatori di oggetti basati sull'apprendimento automatico, inclusi i classificatori basati su CNN tramite la sua estensione di apprendimento profondo. Questa attività utilizza una WSI colorata con E\&E di tessuto di cancro colorettale disponibile pubblicamente per addestrare un classificatore CNN che distingue l'epitelio tumorale, la mucosa normale, lo stroma e lo sfondo.

\textbf{Procedura:}
\begin{enumerate}
  \item Scaricare e installare QuPath (versione 0.5 o successiva) da \url{https://qupath.github.io}. Scaricare la WSI dimostrativa dal sito di documentazione di QuPath.
  \item Aprire la WSI in QuPath e utilizzare lo strumento \textit{Pixel Classifier} (\textit{Classify} $\rightarrow$ \textit{Pixel classification} $\rightarrow$ \textit{Train pixel classifier}) per creare un nuovo classificatore.
  \item Selezionare un'architettura di classificatore basata su CNN (ad esempio, una ResNet leggera) e definire quattro classi tissutali: \textit{Tumore}, \textit{Mucosa normale}, \textit{Stroma} e \textit{Sfondo}.
  \item Annotare regioni rappresentative di ogni classe disegnando poligoni sulla WSI utilizzando gli strumenti di annotazione. Puntare ad almeno 10 annotazioni per classe, distribuite in diverse regioni del vetrino per catturare la variabilità morfologica.
  \item Addestrare il classificatore e valutarne le prestazioni su una regione del vetrino tenuta da parte. Esaminare la sovrapposizione di classificazione e identificare le regioni in cui il classificatore funziona bene e le regioni in cui commette errori.
  \item Esportare i risultati della classificazione come file GeoJSON e calcolare la percentuale di area tissutale occupata da ogni classe. Discutere come queste misurazioni potrebbero essere utilizzate come biomarcatori quantitativi in un contesto di ricerca o clinico.
  \item Riflettere sulle seguenti domande: In che modo la scelta delle annotazioni di addestramento influisce sulle prestazioni del classificatore? Quali tipi di artefatti o caratteristiche morfologiche causano errori di classificazione? Come si validerebbe questo classificatore per l'uso clinico?
\end{enumerate}

\textbf{Risultati attesi:} Gli studenti dovrebbero essere in grado di addestrare un classificatore CNN di base in QuPath, interpretare l'output della classificazione e valutare criticamente le prestazioni e le limitazioni del classificatore.

\subsection{Attività B: Generazione di Referti Patologici Assistita da LLM}

\textbf{Obiettivo:} Esplorare le capacità e le limitazioni di un modello linguistico di grandi dimensioni per la generazione strutturata di referti patologici e l'elaborazione del linguaggio naturale delle note cliniche.

\textbf{Contesto:} Questa attività utilizza un'interfaccia LLM accessibile pubblicamente (ad esempio, un sistema basato su GPT-4 o un equivalente open-source come Llama 3) per dimostrare come gli LLM possano assistere nella generazione di referti sinotici e nell'estrazione di informazioni da referti patologici in testo libero. Gli studenti valuteranno l'accuratezza, la completezza e le potenziali modalità di fallimento degli output generati dagli LLM.

\textbf{Procedura:}
\begin{enumerate}
  \item Accedere all'interfaccia LLM fornita dalla propria istituzione o utilizzare un'API disponibile pubblicamente. Preparare tre estratti di referti patologici in testo libero anonimizzati e de-identificati (forniti dal coordinatore del corso) che rappresentano: (a) una resezione di adenocarcinoma colorettale, (b) una biopsia con ago sottile del seno e (c) una biopsia prostatica con ago.
  \item Per ogni referto, sollecitare l'LLM a estrarre i seguenti elementi di dati strutturati e compilare un modello sinotico: tipo di tumore, grado istologico, dimensione del tumore, stato dei margini, invasione linfovascolare, invasione perineurale, stato dei linfonodi (ove applicabile) e stadio patologico (pTNM).
  \item Confrontare gli elementi di dati estratti dall'LLM con i valori di riferimento forniti dal coordinatore del corso. Calcolare l'accuratezza dell'estrazione per ogni elemento di dati e identificare quali elementi sono estratti in modo più e meno affidabile.
  \item Sollecitare l'LLM a generare un referto sinotico completo per un nuovo caso, fornendo solo la descrizione macroscopica e una breve storia clinica. Valutare il referto generato per accuratezza, completezza e presenza di informazioni allucinatorie.
  \item Discutere le seguenti domande in gruppo: In quali circostanze ci si fiderebbe di un elemento di referto generato dall'LLM senza verifica? Quali salvaguardie dovrebbero essere in atto prima che la refertazione assistita da LLM possa essere utilizzata in un contesto clinico? Come dovrebbero essere documentate e gestite le discrepanze tra gli output dell'LLM e i risultati del patologo?
\end{enumerate}

\textbf{Risultati attesi:} Gli studenti dovrebbero essere in grado di descrivere le capacità e le limitazioni degli LLM per la generazione di referti patologici, identificare le modalità di fallimento comuni inclusa l'allucinazione, e articolare i requisiti di governance per l'impiego clinico.

% ============================================================
\section{Sintesi}
\label{sec:ch4_summary}
% ============================================================

Questo capitolo ha fornito un'introduzione completa alle architetture di apprendimento profondo che stanno trasformando la patologia computazionale. Le reti neurali convoluzionali, con le loro capacità di apprendimento gerarchico delle caratteristiche e di transfer learning, hanno raggiunto prestazioni di livello clinico su benchmark di patologia di riferimento inclusi CAMELYON (rilevamento delle metastasi linfonodali) e PANDA (classificazione di Gleason del cancro alla prostata). Il framework di apprendimento multi-istanza affronta la sfida fondamentale dell'analisi WSI gigapixel decomponendo i vetrini in patch e aggregando le caratteristiche a livello di patch utilizzando meccanismi di attenzione, consentendo l'apprendimento debolmente supervisionato dalle sole etichette a livello di vetrino. Le architetture basate su transformer, con il loro meccanismo di auto-attenzione globale, hanno soppiantato le CNN come stato dell'arte in molti compiti di patologia, e i modelli fondazionali specifici per la patologia---inclusi UNI, CONCH e Prov-GigaPath---forniscono potenti rappresentazioni pre-addestrate che si generalizzano su un'ampia gamma di compiti a valle.

I modelli linguistici di grandi dimensioni offrono il potenziale per automatizzare la generazione di referti strutturati, la refertazione sinottica e l'NLP degli archivi clinici, ma il loro impiego richiede un'attenta attenzione ai rischi di allucinazione, limite di conoscenza e spostamento della distribuzione. I modelli multimodali che integrano le caratteristiche delle immagini istopatologiche con dati genomici, clinici e radiologici superano costantemente i modelli unimodali nei compiti di predizione della sopravvivenza e classificazione dei sottotipi molecolari, sfruttando informazioni complementari tra le modalità attraverso strategie di fusione delle caratteristiche precoce, tardiva o intermedia.

Il biologo svolge un ruolo indispensabile nei flussi di lavoro di patologia assistiti dall'IA, che comprende la preparazione dei dati, il controllo della qualità, la gestione dell'inferenza, l'interpretazione degli output, la segnalazione delle discordanze e il mantenimento della traccia di audit. Man mano che gli strumenti di IA diventano integrati nella pratica clinica di routine, il BMS deve sviluppare nuove competenze tecniche mantenendo le capacità di valutazione critica necessarie per identificare e segnalare i casi in cui gli output dell'IA potrebbero essere inaffidabili. Le attività pratiche in questo capitolo forniscono un'introduzione a queste competenze attraverso l'esperienza pratica con la classificazione tissutale basata su CNN in QuPath e la generazione di referti assistita da LLM.

% ============================================================
\section{Domande di Revisione}
\label{sec:ch4_questions}
% ============================================================

\begin{enumerate}

  \item \label{ch4:q1} Quale componente di una rete neurale convoluzionale è principalmente responsabile dell'introduzione della non-linearità nel modello?
  \begin{enumerate}[label=\Alph*.]
    \item Strato di pooling
    \item Kernel convoluzionale
    \item Funzione di attivazione (ReLU)
    \item Strato completamente connesso
  \end{enumerate}

  \item \label{ch4:q2} Nel contesto dell'architettura CNN, qual è lo scopo principale del max-pooling?
  \begin{enumerate}[label=\Alph*.]
    \item Aumentare la risoluzione spaziale delle mappe delle caratteristiche
    \item Ridurre le dimensioni spaziali e conferire invarianza alla traslazione
    \item Normalizzare le attivazioni su un mini-batch
    \item Prevenire la co-adattazione dei neuroni durante l'addestramento
  \end{enumerate}

  \item \label{ch4:q3} La sfida CAMELYON è stata progettata per valutare gli algoritmi di apprendimento profondo per quale compito diagnostico?
  \begin{enumerate}[label=\Alph*.]
    \item Classificazione di Gleason del cancro alla prostata
    \item Rilevamento delle metastasi del cancro al seno nei linfonodi sentinella
    \item Classificazione dei sottotipi di cancro polmonare
    \item Predizione dell'instabilità dei microsatelliti da immagini E\&E di cancro colorettale
  \end{enumerate}

  \item \label{ch4:q4} Qual è la sfida fondamentale che motiva l'uso dell'apprendimento multi-istanza (MIL) per l'analisi delle immagini di vetrini interi?
  \begin{enumerate}[label=\Alph*.]
    \item Le WSI sono memorizzate in formati proprietari incompatibili con i framework standard di apprendimento profondo
    \item Le WSI gigapixel non possono essere contenute nella memoria GPU e richiedono l'elaborazione basata su patch con etichette a livello di vetrino
    \item Le WSI contengono troppo pochi pixel per addestrare efficacemente una rete neurale convoluzionale
    \item Il MIL è necessario perché i patologi annotano le singole cellule piuttosto che interi vetrini
  \end{enumerate}

  \item \label{ch4:q5} Nel framework CLAM (Clustering-constrained Attention Multiple instance learning), qual è il ruolo della rete di attenzione?
  \begin{enumerate}[label=\Alph*.]
    \item Segmentare le regioni tissutali dal vetro di sfondo
    \item Assegnare un peso scalare a ogni patch che riflette la sua rilevanza per la predizione a livello di vetrino
    \item Eseguire l'aumento dei dati durante l'addestramento
    \item Convertire i vettori di caratteristiche dei patch in immagini RGB per la visualizzazione
  \end{enumerate}

  \item \label{ch4:q6} Quale delle seguenti descrive meglio la differenza architetturale chiave tra i Vision Transformer (ViT) e le reti neurali convoluzionali (CNN)?
  \begin{enumerate}[label=\Alph*.]
    \item Il ViT utilizza connessioni ricorrenti, mentre le CNN utilizzano connessioni feed-forward
    \item Il ViT ha un campo recettivo globale tramite auto-attenzione, mentre le CNN hanno un campo recettivo locale tramite kernel convoluzionali
    \item Il ViT richiede più dati di addestramento delle CNN in tutte le circostanze
    \item Il ViT non può elaborare immagini a colori, mentre le CNN possono
  \end{enumerate}

  \item \label{ch4:q7} Prov-GigaPath è meglio descritto come quale tipo di modello?
  \begin{enumerate}[label=\Alph*.]
    \item Una rete neurale ricorrente per la generazione di referti patologici
    \item Un modello fondazionale per vetrini interi pre-addestrato su oltre 1,3 miliardi di tile di immagini patologiche
    \item Un modello classico di apprendimento automatico basato su caratteristiche morfologiche costruite manualmente
    \item Una rete generativa avversariale per la sintesi di immagini E\&E sintetiche
  \end{enumerate}

  \item \label{ch4:q8} Quale delle seguenti è la preoccupazione di sicurezza più significativa associata all'uso dei modelli linguistici di grandi dimensioni per la generazione di referti patologici?
  \begin{enumerate}[label=\Alph*.]
    \item Gli LLM sono troppo lenti per essere utilizzati in un contesto clinico
    \item Gli LLM possono generare informazioni plausibili ma fattualmente errate (allucinazione)
    \item Gli LLM non possono elaborare testo più lungo di una singola frase
    \item Gli LLM richiedono l'accesso ai dati genomici del paziente per generare un referto
  \end{enumerate}

  \item \label{ch4:q9} Nei modelli di IA multimodali per la predizione della sopravvivenza, quale strategia di fusione addestra modelli separati per ogni modalità e combina le loro predizioni nella fase di output?
  \begin{enumerate}[label=\Alph*.]
    \item Fusione precoce
    \item Fusione intermedia
    \item Fusione tardiva
    \item Fusione con attenzione cross-modale
  \end{enumerate}

  \item \label{ch4:q10} Quale delle seguenti descrive meglio la responsabilità del biologo quando un sistema di IA produce una predizione incoerente con la valutazione clinica preliminare?
  \begin{enumerate}[label=\Alph*.]
    \item Accettare la predizione dell'IA, poiché si basa su più dati rispetto alla valutazione clinica
    \item Scartare la predizione dell'IA e procedere con la sola valutazione clinica
    \item Segnalare il caso discordante e segnalarlo al patologo per la revisione
    \item Ripetere l'analisi dell'IA fino a quando la predizione è coerente con la valutazione clinica
  \end{enumerate}

\end{enumerate}

% ============================================================
%  Capitolo 5 – Sfide, Etica e Prospettive Future
%              nella Patologia Assistita dall'IA
%  University Pathology Course Handbook
%  Included file — do NOT add \documentclass or \begin{document}
% ============================================================

\chapter{Sfide, Etica e Prospettive Future nella Patologia Assistita dall'IA}
\label{ch:5}

% ============================================================
\section{Introduzione}
% ============================================================

I capitoli precedenti di questo manuale hanno introdotto le basi tecniche della patologia digitale, l'architettura dei modelli di apprendimento automatico e i flussi di lavoro pratici attraverso cui l'intelligenza artificiale (IA) viene integrata nella pratica istopatologica. Sarebbe allettante, a questo punto, concludere che la traiettoria dell'IA in patologia sia quella di un progresso lineare e ininterrotto — una storia in cui dataset sempre più grandi e reti neurali sempre più profonde convergono verso un futuro di diagnosi automatizzata e priva di errori. Tale conclusione sarebbe al tempo stesso prematura e intellettualmente irresponsabile. La storia della tecnologia medica è ricca di innovazioni che promettevano trasformazioni ma hanno prodotto complessità, e l'IA non fa eccezione. La riflessione critica sulle limitazioni, le dimensioni etiche e le conseguenze sociali della patologia assistita dall'IA non è una preoccupazione marginale; è una competenza fondamentale per qualsiasi professionista che lavorerà a fianco di questi sistemi.

Questo capitolo affronta le sfide ancora irrisolte nel settore, i quadri etici che devono governare l'implementazione degli strumenti di IA e le questioni filosofiche che emergono quando alle macchine viene chiesto di svolgere compiti precedentemente considerati di esclusiva competenza di esperti umani qualificati. Guarda anche al futuro, esaminando come i ruoli dei patologi, dei biologi e dei tecnici di patologia siano destinati a evolversi in un panorama sempre più plasmato dall'automazione e dal supporto decisionale basato sui dati. Il capitolo è strutturato per passare dalle limitazioni tecniche (Sezioni~\ref{sec:ch5_limitations} e~\ref{sec:ch5_bias}), attraverso l'interpretabilità e l'epistemologia (Sezioni~\ref{sec:ch5_xai} e~\ref{sec:ch5_philosophy}), all'etica e alla governance (Sezione~\ref{sec:ch5_ethics}), e infine ai futuri pratici dell'automazione e della pratica professionale (Sezioni~\ref{sec:ch5_robotics} e~\ref{sec:ch5_technician}).

Un tema centrale in tutto questo capitolo è che l'adozione dell'IA in patologia non è semplicemente un evento tecnico ma sociotecnico — ridisegna le relazioni tra clinici e pazienti, tra istituzioni e regolatori, e tra competenza umana e inferenza algoritmica. Comprendere queste dimensioni è essenziale non solo per coloro che svilupperanno strumenti di IA, ma, forse ancora più importante, per coloro che li utilizzeranno, valideranno e ne saranno responsabili nelle impostazioni cliniche quotidiane. Lo studente post-laurea che completa questo capitolo dovrebbe essere in grado di confrontarsi criticamente con le affermazioni relative all'IA nella letteratura, di identificare potenziali fonti di danno nei sistemi implementati e di contribuire in modo significativo alle discussioni istituzionali sull'adozione responsabile delle tecnologie di patologia digitale \cite{hanna2024}.

% ============================================================
\section{Limitazioni degli Attuali Sistemi di IA in Patologia}
\label{sec:ch5_limitations}
% ============================================================

\subsection{Fallimento della Generalizzazione e Spostamento del Dominio}

Una delle limitazioni più costantemente documentate dei modelli di IA in patologia computazionale è la loro tendenza a ottenere buone prestazioni su dati che assomigliano al loro set di addestramento, ma a degradarsi sostanzialmente quando applicati a dati provenienti da fonti diverse. Questo fenomeno, noto come \emph{spostamento del dominio} o \emph{spostamento della distribuzione}, sorge perché i modelli di apprendimento profondo apprendono associazioni statistiche specifiche per le condizioni in cui i dati di addestramento sono stati acquisiti. Un modello addestrato su immagini di vetrino intero (WSI) digitalizzate con uno scanner Leica Aperio presso un'istituzione può codificare caratteristiche cromatiche specifiche dello scanner, artefatti di compressione e proprietà di messa a fuoco come parte della sua rappresentazione appresa. Quando lo stesso modello viene applicato a vetrini scansionati con un Hamamatsu NanoZoomer presso un'istituzione diversa, queste caratteristiche di basso livello differiscono e le prestazioni del modello possono calare drasticamente — a volte a livelli non migliori del caso \cite{bankhead2022}.

Lo spostamento del dominio non è limitato alle differenze tra scanner. Le variazioni nei protocolli di fissazione del tessuto, nello spessore delle sezioni, nelle procedure di colorazione con ematossilina ed eosina (H\&E), nei mezzi di copertura e persino nell'età dei reagenti possono introdurre differenze sistematiche tra i dataset. Le istituzioni che hanno sviluppato i propri protocolli di colorazione nel corso di decenni possono produrre vetrini con palette cromatiche caratteristiche che sono invisibili all'occhio umano come fonti di bias, ma che vengono prontamente sfruttate dalle reti neurali come caratteristiche discriminative. Ciò significa che un modello può sembrare classificare accuratamente i sottotipi tumorali mentre in realtà si basa su idiosincrasie di colorazione piuttosto che su segnali morfologici genuini. La conseguenza pratica è che i modelli devono essere rigorosamente validati su coorti esterne prima di qualsiasi implementazione clinica, e che la validazione su un singolo sito esterno è insufficiente per stabilire la generalizzabilità attraverso la diversità della pratica di laboratorio nel mondo reale \cite{song2023}.

Affrontare lo spostamento del dominio è diventato un'area di ricerca attiva. Gli algoritmi di normalizzazione della colorazione — che tentano di mappare la distribuzione cromatica di un'immagine sorgente su uno standard di riferimento — possono parzialmente mitigare la variazione di scanner e colorazione, ma introducono i propri artefatti e non risolvono tutte le fonti di spostamento del dominio. Le tecniche di adattamento del dominio, in cui i modelli vengono messi a punto su piccole quantità di dati del dominio target, offrono un'altra via, ma richiedono dati etichettati dal sito target, che possono essere scarsi o costosi da ottenere. L'apprendimento federato, in cui i modelli vengono addestrati in modo collaborativo tra più istituzioni senza condividere i dati grezzi, rappresenta un approccio promettente per costruire modelli più generalizzabili, ma introduce le proprie sfide tecniche e di governance. La lezione fondamentale per il professionista è che le prestazioni riportate di un modello su un dataset benchmark rappresentano un limite superiore, non inferiore, per le sue prestazioni in un nuovo ambiente clinico.

\subsection{Bias del Dataset ed Effetti Batch}

Grandi dataset pubblici come The Cancer Genome Atlas (TCGA) sono stati preziosi nel consentire la ricerca in patologia computazionale, fornendo migliaia di vetrini digitalizzati collegati a dati genomici, clinici e di sopravvivenza. Tuttavia, questi dataset portano bias sistematici che sono sempre più ben caratterizzati. Howard e colleghi hanno dimostrato che i vetrini TCGA contengono firme specifiche del sito — pattern statistici che codificano l'istituzione contribuente piuttosto che le proprietà biologiche del tumore — e che i modelli addestrati su TCGA possono inavvertitamente imparare a classificare per sito piuttosto che per malattia \cite{howard2021}. Questo effetto batch è una forma di confondimento: il modello apprende un'associazione spuria tra l'origine istituzionale e l'esito di interesse, producendo metriche di prestazione gonfiate che non riflettono un apprendimento biologico genuino.

Gli effetti batch nei dataset di patologia derivano dalle stesse fonti dello spostamento del dominio — tipo di scanner, protocollo di colorazione, processazione del tessuto — ma sono particolarmente insidiosi nei dataset multi-sito perché sono correlati con le variabili biologiche di interesse. Ad esempio, se un sito contribuente è specializzato in un particolare sottotipo tumorale e utilizza anche un protocollo di colorazione distintivo, un modello può imparare ad associare quel pattern di colorazione con il sottotipo, raggiungendo un'elevata accuratezza sui dati trattenuti dello stesso sito mentre fallisce sui dati di altri siti. Un design sperimentale attento — inclusa la cross-validazione stratificata per sito e la validazione esterna prospettica — è essenziale per rilevare e mitigare questi effetti, ma tali design non sono universalmente adottati nella letteratura pubblicata, portando a una sovrastima della prontezza degli strumenti di IA per l'uso clinico.

Le conseguenze degli effetti batch si estendono oltre le metriche di prestazione gonfiate. Quando un modello addestrato su un dataset distorto viene implementato in un contesto clinico, i suoi errori possono essere sistematici piuttosto che casuali — può fallire costantemente su vetrini di particolari scanner, particolari protocolli di colorazione o particolari popolazioni di pazienti. Gli errori sistematici sono più pericolosi degli errori casuali perché è meno probabile che vengano rilevati dal monitoraggio della qualità di routine, che in genere si concentra sulle prestazioni medie piuttosto che sulla distribuzione degli errori tra i sottogruppi. Rilevare e caratterizzare gli errori sistematici richiede un'analisi deliberata dei sottogruppi, che dovrebbe essere una componente standard di qualsiasi studio di validazione dell'IA \cite{howard2021}.

\subsection{L'Illusione dell'Oggettività}

Un malinteso persistente nelle discussioni sulla patologia assistita dall'IA è che i metodi digitali e computazionali siano intrinsecamente più oggettivi della valutazione umana. Questa assunzione confonde la quantificazione con l'oggettività e l'automazione con la neutralità. In realtà, ogni fase della pipeline dell'IA — dalla selezione dei casi di addestramento alla scelta del protocollo di annotazione, dalla progettazione della funzione di perdita alla soglia utilizzata per binarizzare una previsione continua — comporta decisioni umane che incorporano valori, assunzioni e potenziali bias nel sistema. L'output di un modello di IA non è una visione da nessun luogo; è il prodotto di una serie di scelte fatte da sviluppatori, annotatori e implementatori, ognuno dei quali porta le proprie prospettive e vincoli \cite{hanna2024}.

L'illusione dell'oggettività è particolarmente consequenziale perché può sopprimere il controllo critico. Quando un patologo non è d'accordo con la diagnosi di un collega, il disaccordo è visibile e può essere discusso. Quando un patologo non è d'accordo con l'output di un sistema di IA, può esserci una pressione istituzionale o psicologica a deferire all'algoritmo, in particolare se il sistema è stato commercializzato come più accurato o coerente della valutazione umana. Questo fenomeno — a volte chiamato \emph{bias da automazione} — è stato documentato nell'aviazione, nella radiologia e in altri domini ad alto rischio, e non c'è motivo di supporre che la patologia ne sia immune. Formare i professionisti a mantenere un impegno critico con gli output dell'IA, piuttosto che un'accettazione passiva, è quindi una componente essenziale di un'implementazione responsabile dell'IA.

Inoltre, le metriche utilizzate per valutare le prestazioni dell'IA sono esse stesse scelte cariche di valori. Accuratezza, sensibilità, specificità e area sotto la curva caratteristica operativa del ricevitore (AUC-ROC) catturano ciascuna aspetti diversi delle prestazioni del modello e possono portare a conclusioni diverse sull'idoneità di un modello per l'uso clinico. Un modello ottimizzato per l'accuratezza complessiva può avere prestazioni scadenti su casi rari ma clinicamente importanti; un modello ottimizzato per la sensibilità può generare un tasso inaccettabile di falsi positivi. La scelta della metrica di valutazione dovrebbe essere guidata dal contesto clinico e dalle conseguenze dei diversi tipi di errore, non dal desiderio di massimizzare un singolo numero a fini di pubblicazione \cite{bankhead2022}.

\subsection{Obsolescenza del Modello e Deriva del Concetto}

I modelli di IA vengono addestrati su dati che riflettono lo stato della pratica patologica in un particolare momento nel tempo. Man mano che la pratica clinica si evolve — attraverso cambiamenti nei criteri diagnostici, l'introduzione di nuovi biomarcatori, aggiornamenti agli schemi di classificazione dell'OMS o cambiamenti nella popolazione di pazienti — le relazioni statistiche codificate in un modello possono diventare obsolete. Questo fenomeno, noto come \emph{deriva del concetto}, significa che un modello che era stato validato come accurato al momento dell'implementazione può gradualmente diventare meno affidabile senza alcuna modifica al modello stesso. A differenza di un libro di testo, che è visibilmente datato dall'anno di pubblicazione, un modello di IA implementato può continuare a generare previsioni sicure molto tempo dopo che il contesto clinico è cambiato in modi che invalidano le sue assunzioni.

La gestione della deriva del concetto richiede un monitoraggio continuo delle prestazioni del modello in fase di implementazione — una pratica nota come \emph{sorveglianza post-commercializzazione} nella terminologia normativa. Questo è tecnicamente impegnativo, richiedendo la raccolta di etichette di verità di base per un campione rappresentativo di casi elaborati dal modello, ed è organizzativamente impegnativo, richiedendo linee chiare di responsabilità per il monitoraggio e per l'attivazione del riaddestramento o del ritiro del modello. Poche istituzioni dispongono attualmente di quadri robusti per questo tipo di governance longitudinale del modello, e i requisiti normativi per la sorveglianza post-commercializzazione dei dispositivi medici basati su IA/ML sono ancora in evoluzione. Lo studente post-laurea dovrebbe essere consapevole che l'implementazione di un modello di IA non è un evento una tantum ma l'inizio di una responsabilità di gestione continua.

\subsection{Il Divario tra Prova di Concetto e Implementazione Clinica}

La letteratura di patologia computazionale contiene migliaia di articoli che dimostrano che i modelli di IA possono eseguire compiti specifici — rilevamento di tumori, classificazione, previsione di biomarcatori — con un'accuratezza paragonabile o superiore a quella di patologi esperti su dataset benchmark curati. Eppure il numero di strumenti di IA che sono stati integrati con successo nei flussi di lavoro clinici di routine rimane piccolo rispetto al volume della ricerca pubblicata. Questo divario tra prova di concetto e implementazione clinica riflette una serie di sfide che sono in gran parte invisibili nelle pubblicazioni accademiche: la necessità di approvazione normativa, il requisito di integrazione con i sistemi informativi di laboratorio (LIS) e le piattaforme di patologia digitale, le esigenze della validazione clinica in contesti prospettici, i costi di implementazione e manutenzione e la gestione del cambiamento organizzativo necessaria per incorporare nuovi strumenti nei flussi di lavoro consolidati \cite{song2023}.

La sola approvazione normativa può richiedere anni e richiede prove di utilità clinica — non solo accuratezza tecnica — nella popolazione di uso previsto. L'integrazione con l'infrastruttura IT esistente è frequentemente sottovalutata come sfida; molti ospedali operano piattaforme LIS legacy che non sono state progettate per interfacciarsi con strumenti di IA, e il lavoro ingegneristico necessario per costruire interfacce robuste e sicure può essere sostanziale. La validazione clinica in contesti prospettici richiede la cooperazione di patologi, clinici e pazienti, e deve essere progettata per rilevare non solo le prestazioni medie ma anche le modalità di fallimento — i casi in cui il modello è sicuramente sbagliato. Lo studente post-laurea che aspira a contribuire alla traduzione clinica degli strumenti di IA dovrebbe essere preparato alla realtà che il percorso da un algoritmo promettente a uno strumento clinico implementato è lungo, costoso e incerto.

\subsection{Requisiti Computazionali e Costi Infrastrutturali}

I modelli di IA all'avanguardia per la patologia computazionale — in particolare quelli basati su grandi modelli fondazionali o vision transformer addestrati su milioni di WSI — richiedono risorse computazionali sostanziali sia per l'addestramento che per l'inferenza. Addestrare un modello di grandi dimensioni da zero può richiedere centinaia di ore GPU su cluster di calcolo ad alte prestazioni, ponendo tale lavoro al di là della portata della maggior parte dei laboratori individuali o delle istituzioni più piccole. Anche l'inferenza — applicare un modello addestrato a nuovi vetrini — può essere computazionalmente impegnativa durante l'elaborazione di WSI gigapixel, richiedendo un'infrastruttura GPU locale o servizi di calcolo basati su cloud, entrambi con costi finanziari.

Questi requisiti infrastrutturali hanno implicazioni per l'equità nell'accesso agli strumenti di IA. I centri medici accademici ben dotati di risorse nei paesi ad alto reddito sono nella posizione migliore per sviluppare, validare e implementare sistemi di IA avanzati, mentre gli ospedali più piccoli, i laboratori comunitari e le istituzioni nei paesi a basso e medio reddito possono mancare dell'infrastruttura computazionale, delle competenze tecniche o delle risorse finanziarie per partecipare alla rivoluzione dell'IA in patologia. Se gli strumenti di IA vengono sviluppati principalmente su dati di istituzioni ben dotate di risorse e implementati principalmente in quegli stessi contesti, rischiano di esacerbare le disuguaglianze esistenti nella qualità diagnostica piuttosto che ridurle. Affrontare questa sfida richiede uno sforzo deliberato per includere istituzioni diverse nei consorzi di sviluppo dell'IA, per sviluppare modelli computazionalmente efficienti adatti a contesti con risorse limitate e per esplorare modelli open-source e di infrastruttura condivisa che riducano il costo di accesso \cite{chen2023}.

% ============================================================
\section{Bias nell'IA: Fonti, Tipologie e Conseguenze}
\label{sec:ch5_bias}
% ============================================================

\subsection{Comprendere il Bias Algoritmico}

Il bias algoritmico si riferisce a errori sistematici negli output di un sistema di IA che producono risultati ingiusti o non equi per particolari gruppi di individui. Nel contesto della patologia assistita dall'IA, il bias può manifestarsi come accuratezza diagnostica differenziale tra sottogruppi di pazienti definiti per età, sesso, etnia, status socioeconomico o origine geografica. È importante capire che il bias algoritmico non è principalmente una proprietà dell'algoritmo stesso, ma dell'intero sistema sociotecnico in cui l'algoritmo è incorporato — inclusi i dati utilizzati per addestrarlo, le etichette assegnate a quei dati, i processi di misurazione che hanno generato i dati e il contesto in cui il modello viene implementato \cite{hanna2024}. Identificare e mitigare il bias richiede quindi attenzione a ciascuno di questi componenti, non solo alle proprietà matematiche del modello.

Le conseguenze del bias algoritmico in patologia non sono astratte. Un modello che sistematicamente ottiene prestazioni inferiori per i pazienti di gruppi demografici sottorappresentati produrrà più errori diagnostici per quei pazienti — tumori mancati, classificazione errata, raccomandazioni di trattamento inappropriate — contribuendo così alle disuguaglianze sanitarie già ben documentate nella pratica patologica convenzionale. Poiché i sistemi di IA possono elaborare grandi volumi di casi rapidamente, un modello distorto implementato su larga scala può causare danni su una scala che sarebbe impossibile per un singolo professionista distorto. Questo effetto di amplificazione rende l'identificazione e la mitigazione del bias nei sistemi di IA una questione di urgente preoccupazione etica e di salute pubblica \cite{seyyedkalantari2021}.

\subsection{Bias dei Dati: Sottorappresentazione dei Gruppi Demografici}

Il bias dei dati sorge quando il dataset di addestramento non rappresenta adeguatamente la diversità della popolazione a cui il modello verrà applicato. In patologia computazionale, questo si manifesta più comunemente come sottorappresentazione di certi gruppi demografici nelle coorti di addestramento. Grandi dataset pubblici come TCGA sono fortemente orientati verso pazienti di centri medici accademici nordamericani, che sono in modo sproporzionato di ascendenza europea e hanno accesso a cure sanitarie di alta qualità. La morfologia tumorale, la composizione dello stroma, i pattern di infiltrazione immunitaria e altre caratteristiche istologiche che i modelli di IA imparano a riconoscere possono variare sistematicamente tra popolazioni con diversi background genetici, esposizioni ambientali e profili di comorbilità. Un modello addestrato prevalentemente su vetrini di un gruppo demografico può quindi avere prestazioni meno accurate quando applicato a pazienti di gruppi sottorappresentati \cite{chen2023}.

Seyyed-Kalantari e colleghi hanno fornito una dimostrazione sorprendente di questo fenomeno nel contesto dell'interpretazione delle radiografie del torace, mostrando che i modelli di IA addestrati su grandi dataset pubblici hanno sistematicamente sottodiagnosticato le condizioni nei pazienti di gruppi emarginati — incluse donne, pazienti neri e pazienti con status socioeconomico inferiore — anche quando l'accuratezza complessiva del modello era elevata \cite{seyyedkalantari2021}. Sebbene questo studio si sia concentrato sulla radiologia piuttosto che sulla patologia, i meccanismi sottostanti — dati di addestramento che non riflettono la diversità della popolazione di pazienti e metriche di valutazione che aggregano le prestazioni tra i gruppi senza esaminare le disparità dei sottogruppi — sono direttamente applicabili alla patologia computazionale. La lezione è che un'elevata accuratezza complessiva è un criterio necessario ma non sufficiente per un sistema di IA equo; l'analisi delle prestazioni dei sottogruppi è essenziale.

\subsection{Bias delle Etichette: Variabilità Inter-Osservatore nelle Annotazioni}

I modelli di IA per la patologia vengono tipicamente addestrati su dataset in cui le etichette di verità di base — tumore presente/assente, grado 1/2/3, positivo/negativo per un biomarcatore — sono state assegnate da patologi umani. Queste etichette non sono infallibili. La variabilità inter-osservatore nella diagnosi patologica è ben documentata in un'ampia gamma di condizioni, dalla classificazione del cancro alla prostata (dove il sistema di Gleason ha subito molteplici revisioni per migliorare la riproducibilità) alla classificazione dei linfomi e alla valutazione dei linfociti infiltranti il tumore. Quando le etichette di addestramento vengono assegnate da un singolo annotatore, o da un piccolo gruppo di annotatori di una particolare istituzione o tradizione formativa, il modello impara a replicare lo stile diagnostico di quegli annotatori — inclusi i loro bias sistematici e le loro idiosincrasie \cite{bankhead2022}.

Il bias delle etichette è particolarmente problematico quando gli annotatori le cui etichette vengono utilizzate per l'addestramento non sono rappresentativi della più ampia comunità di patologi che utilizzeranno il modello. Se un modello viene addestrato su etichette di subspecialisti esperti di un centro di riferimento terziario, potrebbe non generalizzarsi bene alla pratica diagnostica dei patologi generali negli ospedali comunitari, che possono applicare criteri diversi o avere soglie diverse per i casi equivoci. Al contrario, se le etichette di addestramento vengono raccolte da un gran numero di annotatori con diversi livelli di competenza, il modello può apprendere una regressione verso la media che non riesce a catturare i giudizi sfumati degli specialisti esperti. Non esiste una soluzione semplice al bias delle etichette, ma le migliori pratiche includono l'uso di più annotatori indipendenti, l'aggiudicazione esplicita dei disaccordi e la segnalazione delle statistiche di accordo inter-annotatore insieme alle metriche di prestazione del modello.

\subsection{Bias di Misurazione: Variazione di Scanner e Colorazione}

Il bias di misurazione sorge quando il processo di generazione dei dati introduce differenze sistematiche correlate con l'esito di interesse. In patologia digitale, le fonti più importanti di bias di misurazione sono il tipo di scanner e la variazione della colorazione. Come discusso nella Sezione~\ref{sec:ch5_limitations}, diversi scanner producono immagini con diverse caratteristiche cromatiche, risoluzione e proprietà di compressione. Se la distribuzione dei tipi di scanner in un dataset di addestramento è correlata con la distribuzione delle etichette diagnostiche — ad esempio, se i vetrini di uno scanner hanno maggiori probabilità di provenire da un particolare sottotipo tumorale — il modello può imparare a utilizzare le caratteristiche dello scanner come proxy per la diagnosi \cite{howard2021}.

La variazione della colorazione introduce problemi analoghi. La colorazione H\&E non è una procedura standardizzata; il colore preciso dei depositi di ematossilina ed eosina dipende dalla concentrazione e dall'età dei reagenti, dalla durata delle fasi di colorazione, dal pH delle soluzioni e dalla competenza del tecnico. La colorazione immunoistochimica è ancora più variabile, con differenze nel clone dell'anticorpo, nel protocollo di recupero dell'antigene, nel sistema di rilevamento e nel cromogeno che contribuiscono tutti alla variazione nell'intensità e nel pattern di colorazione. I modelli di IA addestrati su immagini immunoistochimiche di un laboratorio possono quindi non generalizzarsi alle immagini di un altro laboratorio, anche quando la biologia sottostante è identica. La standardizzazione dei protocolli di colorazione e l'uso di strumenti di calibrazione cromatica digitale possono ridurre ma non eliminare questa fonte di bias.

\subsection{Bias di Implementazione e Conseguenze per l'Equità Sanitaria}

Il bias di implementazione si riferisce alla discrepanza tra la popolazione su cui un modello è stato sviluppato e validato e la popolazione a cui viene applicato nella pratica. Anche un modello privo di bias rispetto ai suoi dati di addestramento può produrre risultati non equi se viene implementato in un contesto in cui la popolazione di pazienti, il flusso di lavoro clinico o le risorse istituzionali differiscono da quelle del contesto di sviluppo. Ad esempio, un modello validato in un centro accademico ad alto volume con rapido turnover dei vetrini e revisione da parte di patologi esperti può avere prestazioni diverse in un ospedale comunitario dove i vetrini vengono elaborati meno frequentemente, il controllo della qualità è meno rigoroso e il patologo che esamina gli output dell'IA ha meno competenza subspecialistica \cite{zhang2023}.

L'effetto cumulativo del bias dei dati, del bias delle etichette, del bias di misurazione e del bias di implementazione è un rischio che i sistemi di IA svantaggino sistematicamente i pazienti già più vulnerabili agli errori diagnostici — quelli di gruppi demografici sottorappresentati, quelli serviti da istituzioni con risorse insufficienti e quelli con condizioni rare o scarsamente rappresentate nei dataset di addestramento. Affrontare questo rischio richiede un impegno per l'equità come criterio di progettazione esplicito nello sviluppo dell'IA, non semplicemente come un ripensamento. Ciò significa reclutare attivamente coorti di addestramento diverse, condurre analisi delle prestazioni dei sottogruppi come componente standard della valutazione del modello, coinvolgere le comunità che probabilmente saranno interessate dall'implementazione dell'IA e stabilire requisiti normativi per la validazione incentrata sull'equità \cite{chen2023}.

% ============================================================
\section{Interpretabilità e IA Spiegabile}
\label{sec:ch5_xai}
% ============================================================

\subsection{Il Problema della Scatola Nera}

I moderni modelli di apprendimento profondo — in particolare le reti neurali convoluzionali (CNN) e i vision transformer — raggiungono le loro impressionanti prestazioni attraverso l'apprendimento di rappresentazioni altamente complesse e non lineari distribuite su milioni o miliardi di parametri. Questa complessità è precisamente ciò che consente loro di catturare sottili pattern morfologici che possono sfuggire agli osservatori umani, ma li rende anche opachi: in generale non è possibile ispezionare il funzionamento interno di un modello addestrato e capire, in termini interpretabili dall'uomo, perché ha prodotto un particolare output per un particolare input. Questa opacità è comunemente indicata come il \emph{problema della scatola nera}, e ha implicazioni significative per l'implementazione clinica dei sistemi di IA in patologia \cite{amann2020}.

Il problema della scatola nera non è semplicemente un inconveniente tecnico. In medicina clinica, la capacità di spiegare una diagnosi — di articolare il ragionamento che ha portato dalle osservazioni alle conclusioni — è fondamentale per la pratica della medicina e per la fiducia che i pazienti e i clinici ripongono nei processi diagnostici. Un patologo che riporta una diagnosi di carcinoma duttale invasivo può indicare caratteristiche morfologiche specifiche — pleomorfismo nucleare, figure mitotiche, architettura ghiandolare — che giustificano la diagnosi e possono essere esaminate da un collega che cerca un secondo parere. Un modello di IA che produce la stessa diagnosi senza alcuna indicazione di quali caratteristiche dell'immagine abbiano guidato la previsione non può essere esaminato allo stesso modo. Ciò limita la capacità dei patologi di rilevare errori, di imparare dagli output del modello e di assumersi la responsabilità delle diagnosi informate dall'assistenza dell'IA.

\subsection{Metodi di Attribuzione Basati sul Gradiente e Mappe di Attenzione}

Il campo dell'\emph{IA spiegabile} (XAI) ha sviluppato una serie di metodi per generare spiegazioni post-hoc degli output dei modelli di apprendimento profondo. Tra i più ampiamente utilizzati in patologia computazionale vi sono i metodi di attribuzione basati sul gradiente, di cui la Mappatura dell'Attivazione di Classe Ponderata dal Gradiente (Grad-CAM) è il più prominente. Grad-CAM calcola il gradiente dell'output del modello rispetto alle attivazioni di uno strato convoluzionale specificato, e utilizza questi gradienti per produrre una mappa di calore che evidenzia le regioni dell'immagine di input che hanno influenzato più fortemente la previsione. Quando applicato a patch di WSI, le mappe di calore Grad-CAM possono indicare quali regioni del tessuto il modello ha considerato quando ha preso una decisione di classificazione, fornendo una forma di spiegazione spaziale che può essere sovrapposta all'immagine originale per l'ispezione da parte di un patologo \cite{amann2020}.

I modelli basati sull'attenzione — incluse le architetture di apprendimento multi-istanza (MIL) con pooling di attenzione e i vision transformer — offrono una forma correlata di interpretabilità attraverso le mappe di attenzione. In questi modelli, i pesi di attenzione assegnati a diverse patch di immagini o token durante il passaggio in avanti possono essere visualizzati per mostrare quali parti del vetrino hanno contribuito maggiormente alla previsione finale. A differenza di Grad-CAM, che è una spiegazione post-hoc applicata a un modello che non è stato progettato con l'interpretabilità in mente, le mappe di attenzione sono una proprietà intrinseca dell'architettura del modello e possono quindi fornire una rappresentazione più fedele del processo decisionale del modello. Tuttavia, la ricerca ha dimostrato che i pesi di attenzione non sempre corrispondono a caratteristiche interpretabili dall'uomo, e che pesi di attenzione elevati possono essere assegnati a regioni diagnosticamente irrilevanti \cite{bankhead2022}.

\subsection{Valori SHAP e Attribuzione delle Caratteristiche}

SHapley Additive exPlanations (SHAP) è un framework game-teorico per attribuire l'output di un modello alle sue caratteristiche di input, basato sul concetto di valori di Shapley dalla teoria dei giochi cooperativi. Nel contesto della patologia computazionale, i valori SHAP possono essere utilizzati per quantificare il contributo delle singole caratteristiche dell'immagine — o, nei modelli che incorporano metadati clinici, delle singole variabili cliniche — alla previsione di un modello per un caso specifico. I valori SHAP hanno la proprietà desiderabile di soddisfare diversi assiomi matematici di equità (efficienza, simmetria, dummy e additività), rendendoli una base principiata per l'attribuzione delle caratteristiche. Possono essere calcolati per qualsiasi modello, indipendentemente dall'architettura, rendendoli applicabili all'intera gamma di strumenti di IA utilizzati in patologia \cite{amann2020}.

Nonostante il loro fascino teorico, i valori SHAP e altri metodi XAI hanno importanti limitazioni che devono essere comprese dai professionisti. I metodi di spiegazione post-hoc generano approssimazioni del comportamento del modello piuttosto che descrizioni esatte dei calcoli interni del modello; possono essere sensibili alla scelta della distribuzione di base o di riferimento, e possono produrre spiegazioni plausibili ma errate. Inoltre, la risoluzione spaziale delle spiegazioni basate sul gradiente e sull'attenzione è limitata dall'architettura del modello, e le spiegazioni generate a livello di patch potrebbero non tradursi direttamente in interpretazioni a livello di vetrino. Il campo dell'XAI in patologia è attivo e in rapida evoluzione, e i professionisti dovrebbero approcciare i metodi di spiegazione attuali come strumenti utili ma imperfetti piuttosto che come resoconti definitivi del ragionamento del modello.

\subsection{Spiegabilità, Fiducia Clinica e Approvazione Normativa}

L'importanza della spiegabilità nella patologia assistita dall'IA si estende oltre il dominio tecnico nei regni della fiducia clinica, della responsabilità professionale e della conformità normativa. I clinici sono più propensi a fidarsi e utilizzare appropriatamente gli strumenti di IA quando possono comprendere la base delle raccomandazioni dello strumento, e meno propensi a impegnarsi in un bias da automazione acritico. I pazienti hanno un legittimo interesse a capire come l'IA ha contribuito alla loro diagnosi, in particolare quando la diagnosi assistita dall'IA viene utilizzata per informare decisioni di trattamento consequenziali. Gli organi normativi in più giurisdizioni hanno iniziato a richiedere che i dispositivi medici basati su IA/ML forniscano una qualche forma di spiegazione per i loro output, riflettendo un crescente consenso che l'opacità è incompatibile con un'IA clinica sicura e responsabile \cite{zhang2023}.

Le linee guida di reporting SPIRIT-AI e CONSORT-AI, sviluppate da Cruz Rivera e colleghi, forniscono quadri standardizzati per riportare il design e i risultati degli studi clinici che coinvolgono interventi di IA \cite{cruzirivera2020}. SPIRIT-AI (Standard Protocol Items: Recommendations for Interventional Trials — Artificial Intelligence) estende le linee guida SPIRIT per i protocolli degli studi clinici per affrontare considerazioni specifiche dell'IA, inclusa la specifica dell'intervento di IA, la gestione degli aggiornamenti del modello durante lo studio e la segnalazione delle prestazioni tra i sottogruppi. CONSORT-AI (Consolidated Standards of Reporting Trials — Artificial Intelligence) estende le linee guida CONSORT per la segnalazione degli studi per richiedere la divulgazione di informazioni specifiche dell'IA, inclusa la versione del modello utilizzato, la disponibilità del modello per la validazione indipendente e la natura di qualsiasi supervisione umana degli output dell'IA. L'aderenza a queste linee guida è sempre più attesa dalle riviste ad alto impatto e dagli organi normativi, e la familiarità con esse è una competenza importante per i professionisti coinvolti nella ricerca o nella valutazione dell'IA.

% ============================================================
\section{Riflessioni Filosofiche ed Epistemologiche}
\label{sec:ch5_philosophy}
% ============================================================

\subsection{Cosa Significa per una Macchina ``Diagnosticare''?}

La parola \emph{diagnosi} deriva dal greco \emph{diagign\={o}skein} — distinguere, discernere — e nel suo uso clinico porta connotazioni che si estendono ben oltre l'identificazione di un pattern. Una diagnosi è un'affermazione sulla natura della condizione di un paziente, fatta da un agente responsabile che può essere ritenuto responsabile dell'affermazione, comunicata al paziente e al suo team clinico, e utilizzata come base per decisioni consequenziali sul trattamento e sulla prognosi. Quando diciamo che un sistema di IA ha ``diagnosticato'' un tumore, stiamo usando la parola in un senso che elimina la maggior parte di queste connotazioni: il sistema ha identificato un pattern in un'immagine che corrisponde alla firma statistica di una particolare classe nei suoi dati di addestramento. Se questo costituisca una diagnosi in qualsiasi senso significativo è una questione che merita un attento esame filosofico \cite{hanna2024}.

La distinzione tra riconoscimento dei pattern e comprensione è centrale per questa domanda. I modelli di apprendimento profondo sono, nella loro essenza, sofisticati sistemi di corrispondenza dei pattern: imparano ad associare le rappresentazioni di input con le etichette di output ottimizzando una funzione obiettivo matematica su un grande dataset. Non possiedono credenze, intenzioni o comprensione in alcun senso filosoficamente robusto. Quando una CNN classifica correttamente un vetrino come mostrante carcinoma invasivo, lo fa perché le proprietà statistiche dell'immagine rientrano nella regione dello spazio delle caratteristiche associata a quell'etichetta nei dati di addestramento — non perché il modello capisca cos'è il cancro, perché insorge, come colpisce il paziente o cosa significa la diagnosi per la vita del paziente. Questa distinzione è importante perché la comprensione, nel senso umano, è ciò che consente a un patologo di gestire casi nuovi, di riconoscere quando un caso esula dalla portata della propria esperienza e di comunicare l'incertezza e la sfumatura di una diagnosi ai colleghi clinici \cite{song2023}.

\subsection{Conoscenza Tacita e i Limiti della Formalizzazione}

Il filosofo Michael Polanyi ha introdotto il concetto di \emph{conoscenza tacita} per descrivere la dimensione dell'expertise umana che non può essere completamente articolata in regole o procedure esplicite — la conoscenza che è incorporata nella pratica qualificata e che viene trasmessa attraverso l'apprendistato, l'osservazione e l'esperienza piuttosto che attraverso l'istruzione formale. La diagnosi patologica è ricca di conoscenza tacita: il patologo esperto sviluppa un senso intuitivo di come appare un vetrino ``normale'', una capacità di rilevare sottili deviazioni dalla normalità che resistono a una descrizione precisa, e una capacità di integrare le osservazioni morfologiche con il contesto clinico, la storia del paziente e l'esperienza precedente in modi che non sono facilmente formalizzabili. Questa dimensione tacita dell'expertise patologica è precisamente ciò che è più difficile da catturare in un dataset di addestramento e più resistente alla formalizzazione algoritmica.

Le implicazioni di questa osservazione per l'IA in patologia sono significative. I modelli di IA possono essere addestrati per replicare i componenti espliciti e articolabili della diagnosi patologica — l'identificazione di caratteristiche morfologiche specifiche, l'applicazione di criteri di classificazione definiti — ma non possono, al momento, replicare la dimensione tacita della pratica esperta. Ciò significa che gli strumenti di IA sono probabilmente più affidabili in compiti ben definiti e standardizzati dove i criteri diagnostici sono espliciti e i dati di addestramento sono rappresentativi, e meno affidabili in casi nuovi, ambigui o contestualmente complessi — precisamente i casi in cui il giudizio umano esperto è più prezioso. Riconoscere questa asimmetria è essenziale per progettare appropriati quadri di collaborazione uomo-IA che sfruttino i punti di forza sia dell'intelligenza umana che di quella artificiale \cite{amann2020}.

\subsection{Potenziamento Versus Sostituzione: Prospettive Filosofiche}

La questione se l'IA sostituirà o potenzierà il patologo è stata ampiamente dibattuta nella letteratura e nella stampa popolare, spesso generando più calore che luce. Da una prospettiva filosofica, la domanda non è semplicemente empirica — non è semplicemente una questione di se l'IA possa raggiungere un'accuratezza sufficiente — ma normativa: riguarda che tipo di pratica diagnostica vogliamo, quali valori desideriamo incorporare nei nostri sistemi sanitari e quale ruolo crediamo che l'expertise umana e la responsabilità umana debbano svolgere nel processo decisionale medico. Queste sono domande a cui non si può rispondere solo con le metriche di prestazione del benchmark \cite{hanna2024}.

Diversi quadri filosofici sono rilevanti per questo dibattito. Il concetto di \emph{cognizione distribuita}, sviluppato da Edwin Hutchins, suggerisce che i compiti cognitivi non vengono eseguiti da menti individuali in isolamento ma sono distribuiti attraverso reti di persone, strumenti e artefatti. Da questo punto di vista, il sistema patologo-IA è un sistema cognitivo in cui le capacità umane e quelle della macchina sono complementari e mutuamente costitutive, piuttosto che competitive. Il concetto di IA \emph{human-in-the-loop} formalizza questa intuizione in termini ingegneristici, specificando che la supervisione e l'intervento umano dovrebbero essere mantenuti nei punti decisionali critici in qualsiasi flusso di lavoro assistito dall'IA. Il concetto di \emph{controllo umano significativo}, sviluppato nel contesto dei sistemi d'arma autonomi ma sempre più applicato all'IA medica, va oltre, richiedendo non solo che un essere umano sia presente nel ciclo ma che l'essere umano abbia una comprensione sufficiente degli output del sistema di IA per esercitare un controllo genuino e informato sulla decisione finale \cite{zhang2023}.

\subsection{L'Epistemologia della Diagnosi Assistita dall'IA}

Da una prospettiva epistemologica, la diagnosi assistita dall'IA solleva domande sulla natura e la giustificazione della conoscenza diagnostica. Quando un patologo accetta la raccomandazione di un sistema di IA, su quale base è giustificata tale accettazione? Se il patologo non può ispezionare il ragionamento del modello — perché il modello è una scatola nera — allora l'accettazione si basa sulla fiducia nel track record del sistema piuttosto che sulla comprensione del suo ragionamento. Questo non è necessariamente problematico; i clinici si affidano abitualmente a test di laboratorio e modalità di imaging i cui meccanismi sottostanti non comprendono pienamente. Ma significa che lo status epistemologico della diagnosi assistita dall'IA è diverso da quello della diagnosi patologica tradizionale, e che le condizioni in cui tale fiducia è giustificata — validazione rigorosa, monitoraggio continuo, segnalazione trasparente delle prestazioni — devono essere chiaramente specificate e applicate \cite{cruzirivera2020}.

Le sfide epistemologiche della diagnosi assistita dall'IA sono aggravate dal fenomeno dell'\emph{ingiustizia epistemica}, descritto dalla filosofa Miranda Fricker, che si verifica quando una persona viene danneggiata nella sua capacità di conoscitore. Nel contesto della patologia assistita dall'IA, l'ingiustizia epistemica potrebbe sorgere se i pazienti di gruppi emarginati ricevono diagnosi meno accurate perché i sistemi di IA addestrati su dati non rappresentativi sono meno affidabili per i loro casi, o se i patologi di istituzioni meno prestigiose vengono sistematicamente esclusi dallo sviluppo di strumenti di IA che plasmeranno la pratica diagnostica a livello globale. Prestare attenzione a queste dimensioni della giustizia epistemica fa parte della più ampia responsabilità etica di coloro che sviluppano e implementano sistemi di IA in sanità.

% ============================================================
\section{Etica nella Diagnosi Assistita dall'IA}
\label{sec:ch5_ethics}
% ============================================================

\subsection{Consenso Informato e Diritti del Paziente}

Il principio del consenso informato — che i pazienti hanno il diritto di comprendere e acconsentire alle procedure utilizzate nella loro cura — si applica alla diagnosi assistita dall'IA come a qualsiasi altro intervento medico. Tuttavia, l'applicazione di questo principio all'IA solleva nuove domande. Quali informazioni devono essere divulgate a un paziente affinché il suo consenso alla diagnosi assistita dall'IA sia genuinamente informato? Come minimo, i pazienti dovrebbero essere informati che gli strumenti di IA vengono utilizzati nell'analisi dei loro campioni, qual è lo scopo di tali strumenti e quali sono le limitazioni e i rischi noti degli strumenti. Dovrebbero anche essere informati del loro diritto di richiedere una revisione solo umana se hanno preoccupazioni riguardo al coinvolgimento dell'IA nella loro cura \cite{zhang2023}.

In pratica, l'implementazione del consenso informato per la diagnosi assistita dall'IA è complicata da diversi fattori. Gli strumenti di IA possono essere utilizzati in più fasi del flusso di lavoro diagnostico — nel controllo della qualità, nello screening primario, nella quantificazione dei biomarcatori — e potrebbe non essere pratico ottenere un consenso specifico per ogni applicazione. Il linguaggio utilizzato per descrivere gli strumenti di IA ai pazienti deve essere accessibile e non tecnico, il che è impegnativo data la complessità dei sistemi sottostanti. Inoltre, il concetto di consenso è complicato dal fatto che i modelli di IA vengono addestrati su dati di pazienti passati che potrebbero non aver acconsentito all'uso dei loro dati per lo sviluppo dell'IA. L'uso retrospettivo dei dati dei pazienti per l'addestramento dell'IA è disciplinato dalla legislazione sulla protezione dei dati e dai quadri etici della ricerca, ma i confini dell'uso consentito non sono sempre chiari e le pratiche variano tra le giurisdizioni \cite{hanna2024}.

\subsection{Privacy dei Dati e Proprietà Intellettuale}

Le immagini di vetrino intero del tessuto del paziente sono dati sanitari personali e sono soggette alle protezioni previste dalla legislazione sulla protezione dei dati, incluso il Regolamento Generale sulla Protezione dei Dati (GDPR) nell'Unione Europea e la legislazione equivalente in altre giurisdizioni. L'uso delle WSI dei pazienti per l'addestramento dell'IA richiede basi giuridiche appropriate — tipicamente il consenso esplicito o una determinazione di interesse legittimo — e deve rispettare i requisiti di minimizzazione dei dati, limitazione delle finalità e sicurezza. La de-identificazione delle WSI è tecnicamente impegnativa perché le immagini possono contenere informazioni — come i pattern di morfologia del tessuto che sono unici per un individuo — che potrebbero in linea di principio essere utilizzate per re-identificare il paziente, anche in assenza di identificatori espliciti \cite{zhang2023}.

Lo status di proprietà intellettuale (PI) dei modelli di IA addestrati sui dati dei pazienti è un'area del diritto complessa e in evoluzione. Quando un ospedale o un istituto di ricerca addestra un modello di IA utilizzando i dati dei pazienti, chi possiede il modello risultante? L'istituzione che ha fornito i dati? Gli sviluppatori che hanno progettato e addestrato il modello? I pazienti i cui dati sono stati utilizzati? Gli attuali quadri di PI non sono stati progettati tenendo conto dell'IA, e c'è una significativa incertezza giuridica sulla proprietà degli output generati dall'IA e dei modelli che li producono. Questa incertezza ha conseguenze pratiche: influisce sulla capacità delle istituzioni di commercializzare strumenti di IA sviluppati utilizzando i dati dei pazienti, sulla capacità dei pazienti di beneficiare del valore commerciale dei loro dati e sulla capacità dei regolatori di richiedere l'accesso ai pesi del modello per la valutazione della sicurezza \cite{hanna2024}.

\subsection{Responsabilità e Quadri Normativi}

Quando un sistema di IA contribuisce a un errore diagnostico — un tumore mancato, un grado errato, un risultato falso positivo che porta a un trattamento non necessario — chi è responsabile? Il patologo che ha esaminato l'output dell'IA e lo ha accettato? L'istituzione che ha implementato il sistema di IA? L'azienda che ha sviluppato e venduto il sistema? L'organo normativo che lo ha approvato? Queste domande di responsabilità non sono meramente teoriche; hanno implicazioni dirette per la responsabilità legale dei professionisti e delle istituzioni, per gli incentivi degli sviluppatori di IA e per la progettazione dei quadri normativi \cite{elendu2023}.

Il panorama normativo per i dispositivi medici basati su IA/ML si sta evolvendo rapidamente. Negli Stati Uniti, la Food and Drug Administration (FDA) ha sviluppato un quadro per la regolamentazione del Software come Dispositivo Medico (SaMD) basato su IA/ML, che distingue tra algoritmi \emph{bloccati} (i cui parametri non cambiano dopo l'implementazione) e algoritmi \emph{adattativi} (che continuano ad apprendere da nuovi dati dopo l'implementazione). Gli algoritmi adattativi presentano particolari sfide normative perché la versione dell'algoritmo che è stata validata può differire dalla versione in uso in un dato momento. La FDA ha proposto un quadro di \emph{piano di controllo delle modifiche predeterminato}, in cui gli sviluppatori specificano in anticipo i tipi di modifiche che possono essere apportate a un modello e le procedure di validazione che verranno applicate prima che tali modifiche vengano implementate \cite{zhang2023}.

Nell'Unione Europea, il Regolamento UE sull'IA — adottato nel 2024 — stabilisce un quadro normativo basato sul rischio per i sistemi di IA, classificando le applicazioni di IA in sanità come \emph{ad alto rischio} e sottoponendole a requisiti di trasparenza, supervisione umana, robustezza e accuratezza. I sistemi di IA utilizzati per decisioni diagnostiche o terapeutiche devono essere registrati in un database pubblico, devono essere sottoposti a valutazione di conformità prima dell'implementazione e devono essere accompagnati da documentazione tecnica che consenta il controllo normativo. Il Regolamento UE sull'IA stabilisce anche requisiti per il monitoraggio post-commercializzazione e per la segnalazione di incidenti gravi che coinvolgono sistemi di IA. La conformità a questi requisiti imporrà obblighi significativi agli sviluppatori di IA e alle istituzioni che li implementano, e richiederà lo sviluppo di nuove strutture di governance e capacità tecniche \cite{zhang2023}.

\subsection{Responsabilità Professionale e Pratica Etica}

Al di là dei requisiti formali della legge e della regolamentazione, l'implementazione dell'IA in patologia solleva questioni di responsabilità professionale e pratica etica. I patologi e i biologi hanno obblighi professionali — sanciti nei codici di condotta emanati da organismi come il Royal College of Pathologists e l'Institute of Biomedical Science — di mantenere la qualità e l'accuratezza del loro lavoro diagnostico, di agire nel migliore interesse dei pazienti e di impegnarsi nello sviluppo professionale continuo. Questi obblighi non diminuiscono quando gli strumenti di IA vengono introdotti nel flusso di lavoro; se mai, si intensificano, perché il professionista deve ora assumersi la responsabilità non solo dei propri giudizi ma anche dell'uso appropriato e della valutazione critica degli output dell'IA \cite{hanna2024}.

La responsabilità professionale nel contesto della patologia assistita dall'IA include l'obbligo di comprendere le capacità e le limitazioni degli strumenti di IA utilizzati, di mantenere le competenze necessarie per rilevare gli errori dell'IA, di segnalare le preoccupazioni sulle prestazioni dell'IA attraverso i canali appropriati e di difendere i pazienti i cui interessi possono essere compromessi da sistemi di IA distorti, inaccurati o implementati in modo inappropriato. Include anche l'obbligo di confrontarsi con le più ampie dimensioni etiche dell'IA in sanità — di partecipare alle discussioni istituzionali e professionali sull'adozione responsabile dell'IA, di contribuire allo sviluppo di linee guida etiche e quadri normativi e di garantire che le voci dei pazienti e delle comunità sottorappresentate siano ascoltate in queste discussioni \cite{elendu2023}.

% ============================================================
\section{Robotica, Automazione e Sistemi Autonomi}
\label{sec:ch5_robotics}
% ============================================================

\subsection{Dissezione Robotica e Processazione Automatizzata del Tessuto}

L'automazione dei flussi di lavoro di patologia si estende oltre l'analisi digitale dei vetrini per comprendere la manipolazione fisica dei campioni di tessuto. I sistemi di dissezione robotica — che utilizzano bracci robotici, visione artificiale e protocolli di taglio guidati dall'IA per eseguire l'esame macroscopico e la sezionatura dei campioni chirurgici — sono in fase di sviluppo attivo e sono stati dimostrati in contesti di ricerca. Questi sistemi mirano a standardizzare il processo di dissezione, ridurre la variabilità introdotta dalle differenze nella tecnica del tecnico e migliorare la riproducibilità del campionamento del tessuto. Hanno anche il potenziale per ridurre le richieste fisiche sui tecnici di patologia, che eseguono compiti di taglio ripetitivi che comportano rischi di lesioni muscoloscheletriche e esposizione accidentale a oggetti taglienti \cite{elendu2023}.

Le piattaforme di colorazione automatizzata sono in uso routinario nei laboratori di patologia da molti anni, eseguendo la colorazione H\&E, l'immunoistochimica e l'ibridazione in situ con un grado di standardizzazione e throughput che sarebbe impossibile raggiungere manualmente. Sviluppi più recenti includono sistemi completamente integrati di processazione e colorazione del tessuto che possono portare una cassetta di tessuto dalla fissazione attraverso l'embedding, la sezionatura, la colorazione e la copertura con un intervento umano minimo. I sistemi robotici di gestione dei vetrini — che trasportano i vetrini tra le stazioni di processazione, caricano i vetrini sugli scanner e gestiscono lo stoccaggio e il recupero dei vetrini — sono sempre più comuni nei laboratori di patologia digitale ad alto volume. Insieme, queste tecnologie stanno creando una visione del laboratorio di patologia come un ambiente altamente automatizzato e gestito roboticamente in cui l'intervento umano è riservato al controllo della qualità, alla gestione delle eccezioni e al processo decisionale diagnostico \cite{song2023}.

\subsection{Diagnosi Autonoma tramite IA: Concetto e Controversia}

Il concetto di \emph{diagnosi autonoma tramite IA} — in cui un sistema di IA produce un referto diagnostico senza alcuna revisione umana — rappresenta il punto finale logico della traiettoria di automazione in patologia. I sostenitori sostengono che la diagnosi autonoma tramite IA potrebbe aumentare drasticamente il throughput diagnostico, ridurre i tempi di risposta ed estendere l'accesso a diagnosi di alta qualità in contesti in cui l'expertise del patologo è scarsa. Nei paesi a basso e medio reddito, dove il rapporto tra patologi e popolazione è molto inferiore rispetto ai paesi ad alto reddito, la diagnosi autonoma tramite IA potrebbe in linea di principio fornire un livello di servizio diagnostico che altrimenti sarebbe irraggiungibile \cite{elendu2023}.

Tuttavia, il concetto di diagnosi autonoma tramite IA affronta sostanziali barriere normative, etiche e pratiche. I quadri normativi nella maggior parte delle giurisdizioni richiedono che i referti diagnostici siano firmati da un medico qualificato che si assume la responsabilità professionale della diagnosi. Un sistema di IA autonomo non può assumersi la responsabilità professionale in questo senso; non può essere ritenuto responsabile, non può essere disciplinato e non può essere citato in giudizio per negligenza. L'approvazione normativa dei sistemi diagnostici autonomi di IA richiederebbe quindi o una revisione fondamentale del concetto di responsabilità professionale in medicina o lo sviluppo di nuovi quadri giuridici per la responsabilità dell'IA che non esistono ancora. Inoltre, le prestazioni tecniche degli attuali sistemi di IA — sebbene impressionanti in compiti ben definiti — non sono sufficienti a giustificare la rimozione della supervisione umana nell'intera gamma di scenari diagnostici incontrati nella pratica patologica di routine \cite{zhang2023}.

\subsection{Barriere Normative ed Etiche alla Piena Automazione}

Le barriere etiche alla piena automazione in patologia sono altrettanto significative di quelle normative. La rimozione della supervisione umana dai processi diagnostici solleva profonde domande sui valori incorporati nella pratica medica. La medicina non è semplicemente un'impresa tecnica; è un'impresa morale, in cui i professionisti hanno obblighi verso i pazienti che vanno oltre l'identificazione accurata della malattia. Questi obblighi includono il dovere di comunicare le diagnosi con compassione e sensibilità, di considerare i valori e le preferenze del paziente nel contesto dell'incertezza diagnostica e di esercitare il giudizio professionale nei casi che esulano dall'ambito dei protocolli stabiliti. Queste dimensioni della pratica medica non possono essere automatizzate, e la loro importanza non dovrebbe essere diminuita dall'entusiasmo per l'efficienza tecnica \cite{hanna2024}.

Da una prospettiva pratica, l'implementazione di sistemi di IA autonomi in patologia richiederebbe un livello di affidabilità tecnica — in termini di accuratezza, robustezza allo spostamento del dominio e resistenza agli input avversariali — che i sistemi attuali non raggiungono. Richiederebbe anche meccanismi robusti per rilevare e rispondere ai guasti del sistema, per gestire i casi che esulano dalla competenza del modello e per mantenere la qualità dei dati di addestramento da cui il sistema dipende. Lo sviluppo di questi meccanismi è un'area di ricerca attiva, ma è lontano dall'essere completo. Nel frattempo, il modello appropriato per l'IA in patologia è quello della \emph{collaborazione uomo-IA}, in cui gli strumenti di IA potenziano e supportano l'expertise umana piuttosto che sostituirla, e in cui la supervisione umana viene mantenuta come salvaguardia contro gli errori dell'IA \cite{song2023}.

\subsection{Il Futuro dell'Automazione di Laboratorio}

Guardando più avanti, l'integrazione di robotica, IA e automazione di laboratorio è destinata a trasformare l'ambiente fisico del laboratorio di patologia in modi difficili da prevedere in dettaglio ma i cui contorni generali stanno diventando visibili. I laboratori del futuro potrebbero essere caratterizzati da un funzionamento continuo, 24 ore su 24, abilitato da sistemi robotici che non richiedono riposo; da un monitoraggio della qualità in tempo reale che rileva i guasti di processazione prima che influenzino la qualità diagnostica; da un'integrazione senza soluzione di continuità di dati istologici, molecolari e di imaging in piattaforme diagnostiche unificate; e da sistemi di IA in grado di triaggiare i casi per urgenza, instradare i vetrini agli specialisti appropriati e segnalare i casi per la discussione multidisciplinare. Questi sviluppi richiederanno ai tecnici di patologia e ai biologi di sviluppare nuove competenze e di assumere nuovi ruoli nella gestione e supervisione dei sistemi automatizzati \cite{elendu2023}.

% ============================================================
\section{Il Futuro Ruolo del Tecnico di Patologia}
\label{sec:ch5_technician}
% ============================================================

\subsection{L'IA come Trasformazione, non Eliminazione}

Un'ansia comune tra i tecnici di patologia e i biologi è che l'IA e l'automazione renderanno i loro ruoli obsoleti. Questa ansia è comprensibile ma, sulla base delle prove disponibili, esagerata. La storia dell'automazione di laboratorio fornisce precedenti istruttivi: l'introduzione degli analizzatori ematologici automatizzati non ha eliminato il ruolo del tecnico di ematologia ma lo ha trasformato, spostando l'attenzione dal conteggio manuale delle cellule alla gestione degli strumenti, al controllo della qualità e alla revisione dei risultati anomali segnalati. Allo stesso modo, l'introduzione dell'IA nella patologia digitale è destinata a trasformare il ruolo del tecnico piuttosto che eliminarlo, spostando l'attenzione dai compiti routinari e ripetitivi a funzioni di livello superiore che richiedono giudizio umano, comprensione contestuale e responsabilità professionale \cite{song2023}.

I compiti più probabilmente automatizzati nel breve termine sono quelli ben definiti, ripetitivi e suscettibili di formalizzazione algoritmica: screening primario dei vetrini per la presenza di tessuto, controllo della qualità della colorazione e della scansione, quantificazione dei biomarcatori standard e triage dei casi per urgenza. I compiti meno probabilmente automatizzati — e quindi più probabilmente destinati a rimanere di competenza dei professionisti umani — sono quelli che richiedono giudizio contestuale, comunicazione con i colleghi clinici, gestione di risultati inattesi e supervisione dei sistemi automatizzati. Il tecnico del futuro dovrà essere competente sia nelle competenze tradizionali della tecnica istologica che nelle nuove competenze della patologia digitale e della gestione dell'IA \cite{bankhead2022}.

\subsection{Nuove Competenze per l'Era della Patologia Digitale}

L'integrazione dell'IA nei flussi di lavoro di patologia crea una domanda di nuove competenze tecniche tra i tecnici di patologia e i biologi. L'\emph{alfabetizzazione ai dati} — la capacità di comprendere, interpretare e valutare criticamente i dati e i sistemi che li generano — è sempre più essenziale. I tecnici che lavorano con le piattaforme di patologia digitale devono comprendere i principi dell'acquisizione delle immagini e del controllo della qualità, i fattori che influenzano la qualità delle immagini e le implicazioni della qualità delle immagini per le prestazioni dell'IA. Devono essere in grado di identificare gli artefatti — pieghe del tessuto, bolle d'aria, regioni fuori fuoco, irregolarità di colorazione — che possono compromettere l'analisi dell'IA e di intraprendere azioni correttive quando tali artefatti vengono rilevati \cite{bankhead2022}.

La \emph{validazione dell'IA} è un'altra competenza critica. I tecnici responsabili del funzionamento dei sistemi diagnostici assistiti dall'IA devono comprendere i principi della validazione del modello, le metriche utilizzate per valutare le prestazioni dell'IA e le procedure per rilevare e segnalare i guasti del modello. Devono essere in grado di interpretare criticamente gli output dell'IA — di riconoscere quando una raccomandazione dell'IA è plausibile e quando dovrebbe essere messa in discussione — e di escalare le preoccupazioni ai patologi o agli amministratori di sistema quando appropriato. Le competenze di base di \emph{scripting} — la capacità di scrivere semplici script in Python o R per automatizzare le attività di gestione dei dati, generare report di controllo della qualità o eseguire analisi statistiche di base — sono sempre più apprezzate nei laboratori di patologia digitale e rappresentano un significativo differenziatore di carriera per i tecnici che le acquisiscono \cite{song2023}.

\subsection{Sviluppo Professionale Continuo e Percorsi di Carriera}

Il rapido ritmo del cambiamento nella patologia digitale e nell'IA crea sia una sfida che un'opportunità per lo sviluppo professionale continuo (CPD). La sfida è che le conoscenze e le competenze necessarie per lavorare efficacemente in un ambiente di patologia digitale si stanno evolvendo più velocemente di quanto i tradizionali quadri CPD possano accogliere; un tecnico che si è qualificato cinque anni fa potrebbe scoprire che aspetti significativi della sua formazione sono già obsoleti. L'opportunità è che lo sviluppo di nuove competenze nella patologia digitale e nell'IA apre nuovi percorsi di carriera che non erano disponibili per le generazioni precedenti di tecnici di patologia \cite{hanna2024}.

I percorsi di carriera nella patologia digitale includono ruoli nell'implementazione e gestione della patologia digitale, nella validazione dell'IA e nell'assicurazione della qualità, nell'informatica di laboratorio e nel supporto alla ricerca. Alcuni tecnici possono progredire verso ruoli di specialisti di patologia digitale o coordinatori dell'IA, assumendo la responsabilità della gestione delle piattaforme di patologia digitale e dei sistemi di IA in un laboratorio o in una rete ospedaliera. Altri possono passare a ruoli industriali con aziende che sviluppano hardware, software o strumenti di IA per la patologia digitale. Gli organismi professionali, incluso l'Institute of Biomedical Science (IBMS) nel Regno Unito, hanno iniziato a sviluppare quadri CPD e qualifiche specialistiche nella patologia digitale che forniscono un riconoscimento formale di queste nuove competenze. Il coinvolgimento con questi quadri è fortemente incoraggiato per i tecnici che desiderano sviluppare la propria carriera in questo settore.

\subsection{L'Importanza della Supervisione Umana}

Qualunque sia la futura traiettoria dell'IA e dell'automazione in patologia, l'importanza della supervisione umana non può essere sopravvalutata. I sistemi di IA possono fallire in modi inaspettati, sottili e difficili da rilevare senza un monitoraggio attivo. Le conseguenze dei guasti dell'IA non rilevati in un contesto diagnostico possono essere gravi — diagnosi mancate, decisioni di trattamento errate, danni ai pazienti. La supervisione umana — il coinvolgimento attivo e informato dei professionisti qualificati con gli output dell'IA — è la principale salvaguardia contro queste conseguenze. Questa supervisione non è passiva; richiede professionisti che comprendano le capacità e le limitazioni dei sistemi di IA con cui lavorano, che mantengano le competenze necessarie per rilevare gli errori dell'IA e che siano autorizzati a mettere in discussione gli output dell'IA quando il loro giudizio professionale indica che qualcosa non va \cite{zhang2023}.

Il concetto di controllo umano significativo, introdotto nella Sezione~\ref{sec:ch5_philosophy}, è direttamente rilevante qui. Affinché la supervisione umana sia significativa, deve essere più di un esercizio di timbro di gomma in cui un professionista approva formalmente gli output dell'IA senza confrontarsi genuinamente con essi. Richiede che i professionisti abbiano tempo, informazioni e competenze sufficienti per valutare criticamente gli output dell'IA, e che le culture istituzionali e i flussi di lavoro supportino piuttosto che minino questo impegno critico. Creare le condizioni per una supervisione umana significativa è una responsabilità condivisa degli sviluppatori di IA, delle istituzioni che li implementano, degli organismi professionali e dei regolatori — ed è una responsabilità che il tecnico di patologia del futuro sarà chiamato a esercitare ogni giorno \cite{elendu2023}.

% ============================================================
%  AI Ethics Framework Figure
% ============================================================

\begin{figure}[htbp]
    \centering
    \usetikzlibrary{arrows.meta, positioning, calc}

\begin{tikzpicture}[
    every node/.style={font=\small\sffamily},
    main/.style={
        rectangle, rounded corners=5pt,
        fill={rgb,255:red,46;green,134;blue,171},
        text=white,
        minimum width=3.2cm, minimum height=0.85cm,
        align=center, font=\small\sffamily\bfseries
    },
    ann/.style={
        rectangle, rounded corners=3pt,
        fill={rgb,255:red,255;green,183;blue,77},
        text={rgb,255:red,70;green,45;blue,0},
        minimum width=2.9cm, minimum height=0.62cm,
        align=center, font=\footnotesize\sffamily,
        draw={rgb,255:red,220;green,145;blue,25},
        line width=0.5pt
    },
    arr/.style={
        -{Stealth[length=5pt,width=4pt]},
        line width=0.9pt,
        color={rgb,255:red,46;green,134;blue,171}
    },
    dann/.style={
        dashed,
        -{Stealth[length=4pt,width=3pt]},
        line width=0.5pt,
        color={rgb,255:red,180;green,120;blue,30}
    },
    stg/.style={
        font=\footnotesize\sffamily\itshape,
        text={rgb,255:red,46;green,134;blue,171}
    }
]

%% ── Stage labels ────────────────────────────────────────────────────
\node[stg] at (0,   0.72) {Stage 1};
\node[stg] at (4.2, 0.72) {Stage 2};
\node[stg] at (8.4, 0.72) {Stage 3};

%% ── Main pipeline ───────────────────────────────────────────────────
\node[main] (dc) at (0,   0) {Data\\Collection};
\node[main] (md) at (4.2, 0) {Model\\Development};
\node[main] (cd) at (8.4, 0) {Clinical\\Deployment};

\draw[arr] (dc) -- (md);
\draw[arr] (md) -- (cd);

%% ── Annotation boxes: Data Collection ──────────────────────────────
\node[ann] (dc1) at (0.0, -1.55) {Bias in training data};
\node[ann] (dc2) at (0.0, -2.28) {Privacy \& consent};
\node[ann] (dc3) at (0.0, -3.01) {Representativeness};

\draw[dann] (dc.south)  -- (dc1.north);
\draw[dann] (dc1.south) -- (dc2.north);
\draw[dann] (dc2.south) -- (dc3.north);

%% ── Annotation boxes: Model Development ────────────────────────────
\node[ann] (md1) at (4.2, -1.55) {Algorithmic fairness};
\node[ann] (md2) at (4.2, -2.28) {Interpretability};
\node[ann] (md3) at (4.2, -3.01) {Validation};

\draw[dann] (md.south)  -- (md1.north);
\draw[dann] (md1.south) -- (md2.north);
\draw[dann] (md2.south) -- (md3.north);

%% ── Annotation boxes: Clinical Deployment ───────────────────────────
\node[ann] (cd1) at (8.4, -1.55) {Accountability};
\node[ann] (cd2) at (8.4, -2.28) {Human oversight};
\node[ann] (cd3) at (8.4, -3.01) {Regulatory compliance};

\draw[dann] (cd.south)  -- (cd1.north);
\draw[dann] (cd1.south) -- (cd2.north);
\draw[dann] (cd2.south) -- (cd3.north);

%% ── Bottom rule ─────────────────────────────────────────────────────
\draw[gray!30, line width=0.4pt]
    (-1.8, -3.42) -- (10.2, -3.42);

\end{tikzpicture}

    \caption{Il quadro etico dell'IA in patologia, che illustra la pipeline dalla raccolta dei dati attraverso lo sviluppo del modello fino all'implementazione clinica. In ogni fase emergono considerazioni etiche fondamentali: la raccolta dei dati è soggetta al \emph{bias dei dati} (sottorappresentazione dei gruppi demografici, variazione degli scanner ed effetti batch); lo sviluppo del modello deve affrontare l'\emph{equità} (prestazioni eque tra i sottogruppi) e la \emph{spiegabilità} (output interpretabili tramite Grad-CAM, mappe di attenzione e valori SHAP); e l'implementazione clinica richiede la \emph{responsabilità} (linee chiare di responsabilità per gli errori dell'IA), la \emph{conformità normativa} (Regolamento UE sull'IA, quadro FDA SaMD) e la continua \emph{sorveglianza post-commercializzazione}. La supervisione umana è rappresentata come un requisito trasversale che abbraccia tutte le fasi della pipeline.}
    \label{fig:ethics_framework}
\end{figure}

% ============================================================
\section{Attività Pratiche}
\label{sec:ch5_activities}
% ============================================================

\subsection{Attività 1: Dibattito di Gruppo --- Fiducia e Responsabilità nella Patologia Assistita dall'IA}

\textbf{Obiettivo:} Sviluppare la capacità degli studenti di ragionamento etico critico sull'IA in patologia ed esplorare le tensioni tra i potenziali benefici dell'IA e i rischi di bias, errori e lacune di responsabilità.

\textbf{Formato:} Dibattito strutturato in gruppi da sei a otto studenti, con due squadre che sostengono posizioni opposte su una mozione. Le mozioni suggerite includono: \emph{``Questa assemblea ritiene che i sistemi di IA debbano essere autorizzati a emettere referti diagnostici senza revisione umana obbligatoria in contesti con risorse limitate''} o \emph{``Questa assemblea ritiene che i rischi del bias algoritmico nella patologia assistita dall'IA superino i potenziali benefici per le popolazioni svantaggiate.''}

\textbf{Preparazione:} Ogni squadra dovrebbe prepararsi per 30 minuti, attingendo al materiale di questo capitolo e ad almeno due fonti di letteratura primaria. Le squadre dovrebbero considerare le prospettive di pazienti, patologi, tecnici, regolatori e sviluppatori di IA. Ogni squadra dovrebbe designare un oratore principale, un oratore di replica e un verbalizzatore.

\textbf{Procedura:} Il dibattito dovrebbe seguire un formato strutturato: dichiarazioni di apertura (5 minuti per squadra), replica (3 minuti per squadra), domande aperte (10 minuti) e dichiarazioni di chiusura (2 minuti per squadra). Un facilitatore dovrebbe garantire che tutti i partecipanti abbiano l'opportunità di contribuire e dovrebbe sollecitare le squadre a considerare prospettive che potrebbero aver trascurato.

\textbf{Debriefing:} Dopo il dibattito, il facilitatore dovrebbe condurre una discussione di debriefing di 15 minuti esplorando le seguenti domande: (1) Quali argomenti sono stati più persuasivi e perché? (2) C'erano prospettive che nessuna delle due squadre ha adeguatamente rappresentato? (3) Come ha cambiato il dibattito il vostro pensiero sulle dimensioni etiche dell'IA in patologia? (4) Quali meccanismi istituzionali o normativi sarebbero necessari per affrontare le preoccupazioni sollevate? Gli studenti dovrebbero presentare una riflessione individuale di 500 parole sul dibattito entro una settimana.

\textbf{Risultati di apprendimento:} Al completamento di questa attività, gli studenti saranno in grado di articolare i principali argomenti etici a favore e contro la diagnosi autonoma tramite IA; identificare le prospettive degli stakeholder rilevanti per le decisioni di implementazione dell'IA; e applicare quadri etici a scenari nuovi nella patologia assistita dall'IA.

\subsection{Attività 2: Workshop --- Carriere e Tendenze Future nella Patologia Digitale}

\textbf{Obiettivo:} Aiutare gli studenti a comprendere il panorama in evoluzione delle opportunità di carriera nella patologia digitale e nell'IA, e a sviluppare un piano personale di sviluppo professionale che incorpori competenze digitali e di IA.

\textbf{Formato:} Workshop di mezza giornata che combina una discussione a panel con relatori invitati, esercizi in piccoli gruppi e riflessione individuale.

\textbf{Discussione a panel (60 minuti):} Invitare da tre a quattro relatori che rappresentano diversi percorsi di carriera nella patologia digitale: un patologo consulente con un interesse subspecialistico nella patologia digitale; un biologo o tecnico di patologia che lavora in un ruolo di implementazione della patologia digitale; un data scientist o ingegnere dell'IA che lavora in un'azienda di IA per la patologia; e, se possibile, un difensore dei pazienti o un rappresentante di un'organizzazione di pazienti con interesse nell'IA in sanità. Ogni relatore dovrebbe tenere una presentazione di 10 minuti sul proprio percorso di carriera e sulla propria prospettiva sul futuro della patologia digitale, seguita da 20 minuti di domande degli studenti.

\textbf{Esercizio in piccoli gruppi (45 minuti):} Gli studenti lavorano in gruppi da quattro a cinque per mappare le competenze richieste per tre ruoli di carriera nella patologia digitale di loro scelta, utilizzando un modello strutturato che copre: (a) competenze tecniche (analisi delle immagini, scripting, validazione dell'IA); (b) conoscenze scientifiche (istologia, patologia molecolare, statistica); (c) competenze professionali (comunicazione, gestione della qualità, consapevolezza normativa); e (d) percorsi di sviluppo professionale continuo (qualifiche IBMS, corsi online, iscrizione a società professionali). I gruppi presentano i loro risultati all'intero gruppo in un riassunto di cinque minuti.

\textbf{Riflessione individuale (30 minuti):} Ogni studente completa un piano personale di sviluppo professionale identificando: (a) due competenze di patologia digitale o IA che desidera sviluppare nei prossimi 12 mesi; (b) le azioni specifiche che intraprenderà per sviluppare queste competenze (corsi, letture, esperienza pratica); e (c) come dimostrerà queste competenze a un futuro datore di lavoro o organismo professionale. Gli studenti sono incoraggiati a riesaminare e aggiornare questo piano a intervalli di sei mesi durante la loro formazione post-laurea.

\textbf{Risultati di apprendimento:} Al completamento di questa attività, gli studenti saranno in grado di descrivere almeno tre percorsi di carriera nella patologia digitale e nell'IA; identificare le competenze tecniche, scientifiche e professionali richieste per i ruoli nella patologia digitale; e costruire un piano personale di sviluppo professionale che incorpori competenze digitali e di IA.

% ============================================================
\section{Sintesi}
\label{sec:ch5_summary}
% ============================================================

Questo capitolo ha esaminato le sfide, le dimensioni etiche e le prospettive future della patologia assistita dall'IA da molteplici angolazioni. Le limitazioni tecniche degli attuali sistemi di IA — inclusi il fallimento della generalizzazione, il bias del dataset, l'illusione dell'oggettività, l'obsolescenza del modello, il divario tra prova di concetto e implementazione clinica e i costi dell'infrastruttura computazionale — ci ricordano che gli strumenti di IA sono potenti ma imperfetti, e che la loro implementazione richiede una validazione rigorosa, un monitoraggio continuo e un impegno critico da parte dei professionisti \cite{bankhead2022,howard2021}.

L'analisi del bias nei sistemi di IA — che comprende il bias dei dati, il bias delle etichette, il bias di misurazione e il bias di implementazione — evidenzia il rischio che gli strumenti di IA esacerbino le disuguaglianze sanitarie esistenti se vengono sviluppati e implementati senza un'attenzione esplicita all'equità \cite{chen2023,seyyedkalantari2021}. Il campo dell'IA spiegabile offre soluzioni parziali ma imperfette al problema della scatola nera, fornendo strumenti come Grad-CAM, mappe di attenzione e valori SHAP che possono rendere gli output dell'IA più interpretabili, mentre le linee guida SPIRIT-AI e CONSORT-AI forniscono quadri per la segnalazione trasparente della ricerca sull'IA \cite{amann2020,cruzirivera2020}.

La riflessione filosofica sulla natura della diagnosi automatica, il ruolo della conoscenza tacita in patologia e lo status epistemologico della diagnosi assistita dall'IA rivela che l'integrazione dell'IA in patologia non è semplicemente un evento tecnico ma una trasformazione del panorama epistemico e morale della medicina diagnostica \cite{hanna2024,song2023}. Le dimensioni etiche della diagnosi assistita dall'IA — che comprendono il consenso informato, la privacy dei dati, la proprietà intellettuale, la responsabilità e la conformità normativa — richiedono lo sviluppo di nuovi quadri di governance e nuove competenze professionali \cite{zhang2023,elendu2023}.

La discussione sulla robotica, l'automazione e i sistemi autonomi illustra la traiettoria verso flussi di lavoro di patologia sempre più automatizzati, evidenziando al contempo le barriere normative ed etiche alla piena automazione e l'importanza duratura della supervisione umana. Infine, l'analisi del futuro ruolo del tecnico di patologia sottolinea che l'IA e l'automazione trasformeranno piuttosto che elimineranno il ruolo del tecnico, creando nuove esigenze di alfabetizzazione ai dati, competenze di validazione dell'IA e competenze di patologia digitale, e aprendo nuovi percorsi di carriera per coloro che abbracciano questi cambiamenti \cite{bankhead2022,song2023}.

Il messaggio generale di questo capitolo è che l'integrazione responsabile dell'IA in patologia richiede non solo eccellenza tecnica ma anche consapevolezza etica, pensiero critico e un impegno per i valori di equità, trasparenza e responsabilità. Questi non sono optional; sono competenze fondamentali per il professionista della patologia del ventunesimo secolo.

% ============================================================
\section{Domande di Revisione}
\label{sec:ch5_questions}
% ============================================================

\begin{enumerate}

    \item \label{ch5:q1} Un modello di apprendimento profondo addestrato su immagini di vetrino intero di una singola istituzione raggiunge il 95\% di accuratezza su un set di test trattenuto della stessa istituzione, ma solo il 72\% di accuratezza quando applicato a vetrini di un'istituzione diversa che utilizza uno scanner diverso. Questo fenomeno è meglio descritto come:
    \begin{enumerate}[label=\Alph*.]
        \item Bias delle etichette derivante dalla variabilità inter-osservatore nelle annotazioni
        \item Spostamento del dominio dovuto a differenze nelle condizioni di acquisizione delle immagini
        \item Deriva del concetto causata da cambiamenti nei criteri diagnostici nel tempo
        \item Bias di implementazione risultante da differenze nei dati demografici dei pazienti
    \end{enumerate}

    \item \label{ch5:q2} Howard e colleghi hanno dimostrato che i modelli addestrati sul dataset The Cancer Genome Atlas (TCGA) possono inavvertitamente imparare a classificare i vetrini per istituzione contribuente piuttosto che per biologia tumorale. La causa principale di questo fenomeno è:
    \begin{enumerate}[label=\Alph*.]
        \item Capacità insufficiente del modello per apprendere caratteristiche specifiche del tumore
        \item Firme specifiche del sito derivanti da differenze nella processazione del tessuto e nella scansione tra le istituzioni contribuenti
        \item Etichettatura deliberatamente errata dei vetrini da parte delle istituzioni contribuenti
        \item L'uso del transfer learning da dataset di immagini non patologiche
    \end{enumerate}

    \item \label{ch5:q3} Quale delle seguenti descrive meglio il concetto di ``bias da automazione'' nel contesto della patologia assistita dall'IA?
    \begin{enumerate}[label=\Alph*.]
        \item La tendenza dei modelli di IA a ottenere prestazioni migliori nei compiti automatizzati rispetto ai compiti che richiedono il giudizio umano
        \item Le prestazioni sistematicamente inferiori dei modelli di IA per i pazienti di gruppi demografici sottorappresentati
        \item La tendenza dei professionisti umani ad accettare acriticamente le raccomandazioni dell'IA senza una valutazione indipendente
        \item Il bias introdotto nei modelli di IA dall'uso di strumenti di annotazione automatizzati
    \end{enumerate}

    \item \label{ch5:q4} La Mappatura dell'Attivazione di Classe Ponderata dal Gradiente (Grad-CAM) è un metodo utilizzato nell'IA spiegabile. Quale delle seguenti descrive più accuratamente come Grad-CAM genera il suo output?
    \begin{enumerate}[label=\Alph*.]
        \item Riaddestra il modello su un sottoinsieme dei dati di addestramento per identificare gli esempi di addestramento più influenti
        \item Calcola il gradiente dell'output del modello rispetto alle attivazioni di uno strato convoluzionale per produrre una mappa di calore spaziale
        \item Utilizza i valori di Shapley game-teorici per attribuire l'output del modello alle singole caratteristiche di input
        \item Genera immagini controfattuali perturbando l'input e osservando i cambiamenti nell'output del modello
    \end{enumerate}

    \item \label{ch5:q5} Seyyed-Kalantari e colleghi hanno dimostrato che i modelli di IA addestrati su grandi dataset pubblici hanno sistematicamente sottodiagnosticato le condizioni nei pazienti di gruppi emarginati. Questa scoperta è più direttamente rilevante per quale tipo di bias dell'IA?
    \begin{enumerate}[label=\Alph*.]
        \item Bias di misurazione derivante dalla variazione degli scanner
        \item Bias delle etichette derivante dalla variabilità inter-osservatore
        \item Bias dei dati derivante dalla sottorappresentazione dei gruppi demografici nei dataset di addestramento
        \item Bias di implementazione derivante da differenze nel flusso di lavoro clinico
    \end{enumerate}

    \item \label{ch5:q6} Le linee guida di reporting CONSORT-AI sono state sviluppate per affrontare quale delle seguenti lacune nella segnalazione della ricerca sull'IA?
    \begin{enumerate}[label=\Alph*.]
        \item L'assenza di metodi standardizzati per l'addestramento di modelli di apprendimento profondo su dataset di patologia
        \item La mancanza di requisiti per la divulgazione di informazioni specifiche dell'IA nei report degli studi clinici che coinvolgono interventi di IA
        \item Il fallimento delle linee guida esistenti nell'affrontare i requisiti computazionali dell'addestramento dei modelli di IA
        \item L'assenza di quadri normativi per l'approvazione dei dispositivi medici basati sull'IA
    \end{enumerate}

    \item \label{ch5:q7} Nel contesto della patologia assistita dall'IA, il concetto di ``controllo umano significativo'' richiede quale delle seguenti?
    \begin{enumerate}[label=\Alph*.]
        \item Che un professionista umano sia fisicamente presente in laboratorio quando i sistemi di IA sono in funzione
        \item Che i sistemi di IA siano progettati per replicare i processi decisionali umani nel modo più fedele possibile
        \item Che i professionisti umani abbiano una comprensione sufficiente degli output dell'IA per esercitare una supervisione genuina e informata delle decisioni diagnostiche finali
        \item Che i sistemi di IA vengano riaddestrati da esperti umani a intervalli regolari per prevenire la deriva del concetto
    \end{enumerate}

    \item \label{ch5:q8} Il Regolamento UE sull'IA classifica i sistemi di IA utilizzati a fini diagnostici in sanità in quale categoria di rischio, e qual è la principale conseguenza di questa classificazione?
    \begin{enumerate}[label=\Alph*.]
        \item Rischio basso; i sistemi diagnostici di IA sono esenti dai requisiti normativi ma devono mostrare un avviso di trasparenza
        \item Rischio moderato; i sistemi diagnostici di IA devono essere registrati presso le autorità nazionali ma non richiedono una valutazione di conformità
        \item Rischio alto; i sistemi diagnostici di IA sono soggetti a requisiti di trasparenza, supervisione umana, robustezza e valutazione di conformità prima dell'implementazione
        \item Rischio inaccettabile; i sistemi diagnostici di IA sono vietati dall'uso in contesti clinici senza esenzione speciale
    \end{enumerate}

    \item \label{ch5:q9} Quale delle seguenti descrive meglio il concetto di ``deriva del concetto'' nel contesto dei modelli di IA implementati in patologia?
    \begin{enumerate}[label=\Alph*.]
        \item Il graduale miglioramento delle prestazioni del modello che si verifica quando il modello viene esposto a più dati dopo l'implementazione
        \item Il degrado delle prestazioni del modello che si verifica quando le relazioni statistiche nell'ambiente di implementazione cambiano rispetto a quelle nell'ambiente di addestramento
        \item La tendenza dei modelli di IA a generare previsioni sempre più sicure nel tempo, indipendentemente dalla qualità dell'input
        \item Il cambiamento nell'architettura del modello che si verifica quando un modello viene aggiornato per incorporare nuovi dati di addestramento
    \end{enumerate}

    \item \label{ch5:q10} Un laboratorio di patologia sta valutando l'implementazione di un sistema di IA per lo screening primario dei vetrini di citologia cervicale. Quale delle seguenti combinazioni di azioni rifletterebbe meglio i principi di implementazione responsabile dell'IA discussi in questo capitolo?
    \begin{enumerate}[label=\Alph*.]
        \item Implementare il sistema immediatamente sulla base delle prestazioni benchmark pubblicate e monitorare i reclami del personale clinico
        \item Validare il sistema su un campione rappresentativo di vetrini della popolazione di pazienti locale, stabilire un protocollo di revisione umana per tutti gli output dell'IA, implementare la sorveglianza post-commercializzazione e ottenere l'approvazione normativa appropriata
        \item Fare affidamento sui dati di validazione del produttore e sull'approvazione normativa, senza condurre una validazione locale, poiché ciò duplicherebbe le prove esistenti
        \item Implementare il sistema in modalità completamente autonoma per massimizzare il throughput, con revisione umana riservata ai casi segnalati come incerti dall'IA
    \end{enumerate}

\end{enumerate}


%% ---- MATERIALE FINALE ----
\backmatter

%% Bibliografia
\cleardoublepage
\addcontentsline{toc}{chapter}{Bibliografia}
\bibliography{main}

%% Soluzioni
\cleardoublepage
\setcounter{chapter}{5}
%% ============================================================
%%  Soluzioni alle Domande di Revisione
%%  Dispensa del Corso di Patologia Post-Laurea
%%  Questo file è incluso con \input{} da main.tex; NON aggiungere il preambolo qui.
%% ============================================================

\chapter*{Soluzioni alle Domande di Revisione}
\addcontentsline{toc}{chapter}{Soluzioni alle Domande di Revisione}
\label{ch:solutions}

Le seguenti soluzioni forniscono la risposta corretta e una breve spiegazione per ciascuna delle
domande a scelta multipla presentate alla fine di ogni capitolo. Le spiegazioni sono concepite
per rinforzare i concetti chiave e chiarire i motivi per cui le risposte errate sono scorrette.

% ─────────────────────────────────────────────────────────────
\section*{Capitolo 1: Introduzione alla Patologia Digitale}
\addcontentsline{toc}{section}{Capitolo 1}

\begin{enumerate}

    \item \textbf{Risposta: C --- Immagini di vetrini interi (WSI) ad alta risoluzione.}\\
    La patologia digitale si basa sulla digitalizzazione dei vetrini istologici in immagini di
    vetrini interi (WSI) ad alta risoluzione, che possono essere archiviate, trasmesse e
    analizzate computazionalmente. Le opzioni A, B e D descrivono tecnologie o approcci non
    specifici della patologia digitale o non pertinenti alla sua definizione fondamentale.

    \item \textbf{Risposta: B --- Hamamatsu NanoZoomer, Leica Aperio, Philips IntelliSite.}\\
    Questi sono i principali produttori commerciali di scanner per vetrini interi (WSI) utilizzati
    nella patologia digitale. Le altre opzioni elencano strumenti di microscopia convenzionale,
    sistemi di imaging non correlati o combinazioni non corrette.

    \item \textbf{Risposta: C --- Circa 0,25 \textmu m per pixel.}\\
    Alla scansione a 40$\times$, la dimensione del pixel è tipicamente di circa 0,25~\textmu m,
    sufficiente a risolvere i dettagli subcellulari fini come la cromatina nucleare e le membrane
    cellulari. La scansione a 20$\times$ produce pixel di circa 0,5~\textmu m.

    \item \textbf{Risposta: B --- TIFF piramidale con tessere (ad es., SVS, NDPI, MRXS).}\\
    I formati WSI sono implementazioni proprietarie o aperte del TIFF piramidale con tessere, che
    consente la visualizzazione efficiente a diverse risoluzioni senza caricare l'intera immagine
    in memoria. I formati JPEG e PNG non supportano la struttura piramidale necessaria per le WSI.

    \item \textbf{Risposta: C --- OpenSlide.}\\
    OpenSlide è la libreria open source standard per la lettura di file WSI in formati proprietari
    (SVS, NDPI, MRXS, ecc.) in Python e altri linguaggi di programmazione. Le altre opzioni non
    sono librerie per la lettura di WSI.

    \item \textbf{Risposta: B --- Consente la visualizzazione efficiente a diverse risoluzioni
    senza caricare l'intera immagine in memoria.}\\
    La struttura piramidale delle WSI memorizza l'immagine a più livelli di risoluzione (livelli
    di zoom), consentendo ai visualizzatori di caricare solo le tessere necessarie per il livello
    di zoom corrente. Questo è essenziale per la gestione di file WSI che possono superare i
    100.000 $\times$ 100.000 pixel.

    \item \textbf{Risposta: C --- QuPath.}\\
    QuPath è la piattaforma open source più ampiamente utilizzata per l'analisi delle immagini di
    patologia digitale. Supporta la visualizzazione di WSI, l'annotazione, la segmentazione
    cellulare e tissutale, e l'integrazione con script Groovy e Python per l'analisi
    personalizzata.

    \item \textbf{Risposta: B --- Ematossilina ed eosina (E\&E).}\\
    La colorazione E\&E è il cardine della diagnosi istopatologica. L'ematossilina colora i nuclei
    di blu-viola e l'eosina colora il citoplasma e la matrice extracellulare di rosa. Tutte le
    altre colorazioni elencate sono colorazioni speciali o immunoistochimiche utilizzate per scopi
    diagnostici specifici.

    \item \textbf{Risposta: C --- Ridurre la variabilità cromatica inter-laboratorio per migliorare
    la generalizzabilità dei modelli di IA.}\\
    La normalizzazione della colorazione è una fase di pre-elaborazione che trasforma i colori
    delle WSI verso uno standard di riferimento comune, riducendo il \emph{domain shift} cromatico
    che può degradare le prestazioni dei modelli di IA quando applicati a dati di laboratori
    diversi da quelli di addestramento.

    \item \textbf{Risposta: B --- Apprendimento profondo, in particolare le reti neurali
    convoluzionali (CNN).}\\
    Le reti neurali convoluzionali (CNN) e le loro varianti (ResNet, EfficientNet, Vision
    Transformer) sono i metodi di apprendimento automatico dominanti nella patologia computazionale
    per compiti come la classificazione dei tumori, la segmentazione cellulare e la predizione
    degli esiti. I metodi tradizionali di apprendimento automatico (SVM, foreste casuali) sono
    stati in gran parte superati dall'apprendimento profondo per i compiti di analisi delle
    immagini.

\end{enumerate}

% ─────────────────────────────────────────────────────────────
\section*{Capitolo 2: Flusso di Lavoro Digitale e Controllo Qualità}
\addcontentsline{toc}{section}{Capitolo 2}

\begin{enumerate}

    \item \textbf{Risposta: C --- Formalina neutra tamponata al 10\%.}\\
    La formalina neutra tamponata al 10\% (FNT) è il fissativo standard per la processazione
    tissutale istologica di routine. Reticola le proteine preservando l'architettura cellulare e
    prevenendo l'autolisi. La glutaraldeide è utilizzata per la microscopia elettronica; l'alcool
    etilico è usato per alcuni campioni citologici; la soluzione di Bouin è un fissativo
    alternativo usato raramente nella pratica di routine moderna.

    \item \textbf{Risposta: B --- Ridurre la variazione nell'aspetto della colorazione causata da
    laboratori, protocolli e scanner diversi.}\\
    La normalizzazione della colorazione è una fase di pre-elaborazione computazionale che
    trasforma la distribuzione cromatica di un'immagine sorgente per farla corrispondere a quella
    di un'immagine di riferimento, riducendo la variabilità cromatica inter-laboratorio. Non
    aumenta la risoluzione, non elimina la necessità di IIC e non converte le immagini in scala
    di grigi.

    \item \textbf{Risposta: C --- Spessore tissutale non uniforme.}\\
    Lo spessore tissutale non uniforme --- causato da pieghe tissutali, bolle d'aria sotto il
    coprioggetto o un coprioggetto non uniforme --- impedisce all'algoritmo di messa a fuoco
    automatica dello scanner di mantenere la messa a fuoco sull'intera superficie del vetrino,
    producendo regioni sfocate nell'immagine. La sovra-colorazione e la contaminazione batterica
    non causano sfocatura; l'uso di un obiettivo a basso ingrandimento riduce la risoluzione ma
    non produce sfocatura nel senso tecnico.

    \item \textbf{Risposta: B --- Separare i contributi dei singoli coloranti (ad es., ematossilina
    e DAB) da un'immagine RGB.}\\
    La deconvoluzione cromatica di Ruifrok e Johnston utilizza la legge di Beer-Lambert per
    separare matematicamente i contributi dei singoli coloranti al colore di ogni pixel in
    un'immagine RGB. Questo consente l'analisi quantitativa separata dei canali dell'ematossilina
    (nucleare) e dell'eosina o DAB (citoplasmatico/membranoso).

    \item \textbf{Risposta: B --- Uno strumento open source per la valutazione automatizzata della
    qualità delle immagini di vetrini interi.}\\
    HistoQC è una pipeline open source e configurabile per il controllo qualità delle WSI,
    sviluppata da Janowczyk et al. Calcola metriche di qualità dell'immagine (messa a fuoco,
    intensità della colorazione, copertura tissutale, artefatti) e genera rapporti di qualità per
    vetrino. Non è un sistema commerciale, un protocollo di colorazione o un formato di file.

    \item \textbf{Risposta: C --- Tutte le fasi dalla ricezione del campione alla scansione del
    vetrino.}\\
    La fase pre-analitica comprende tutte le fasi che precedono la scansione: ricezione del
    campione, descrizione macroscopica, fissazione, processazione tissutale, inclusione,
    sezionamento, colorazione e coprioggetto. Queste fasi esercitano la maggiore influenza sulla
    qualità finale dell'immagine.

    \item \textbf{Risposta: B --- Il tempo di fissazione influenza la morfologia tissutale, la
    preservazione degli antigeni e le analisi IIC e molecolari successive.}\\
    Il tempo di fissazione è critico perché determina il grado di reticolazione proteica. La
    sotto-fissazione porta a scarso dettaglio nucleare e IIC inaffidabile; la sovra-fissazione
    causa una reticolazione eccessiva che può mascherare gli epitopi antigenici. Il tempo di
    fissazione non influenza la dimensione del file WSI, il numero di livelli della piramide o
    la velocità di scansione.

    \item \textbf{Risposta: C --- Riscansionare il vetrino dopo aver verificato la calibrazione
    della messa a fuoco dello scanner.}\\
    Il primo passo appropriato quando si scopre che una WSI presenta significative regioni sfocate
    è riscansionare il vetrino, dopo aver verificato la calibrazione della messa a fuoco dello
    scanner e ispezionato il vetrino per eventuali problemi fisici (pieghe, bolle d'aria). Accettare
    l'immagine sfocata per la diagnosi è inappropriato; eliminare l'immagine senza tentare la
    rimediazione è prematuro; applicare un filtro di nitidezza computazionale non recupera le
    informazioni perse a causa della sfocatura.

    \item \textbf{Risposta: B --- Garantire la tracciabilità dei campioni e dei vetrini attraverso
    il flusso di lavoro.}\\
    La documentazione della catena di custodia garantisce che ogni campione, cassetta, vetrino e
    immagine digitale possa essere tracciato attraverso il flusso di lavoro, dalla ricezione alla
    diagnosi. Questo è un requisito legale ed etico che protegge i pazienti dagli errori di
    identificazione del campione.

    \item \textbf{Risposta: B --- ISO 15189.}\\
    ISO 15189 è lo standard internazionale specifico per i laboratori medici, che specifica i
    requisiti per la qualità e la competenza nei test di laboratorio medico. ISO 9001 è uno
    standard generico per i sistemi di gestione della qualità; ISO 27001 riguarda la sicurezza
    delle informazioni; ISO 13485 riguarda i dispositivi medici.

\end{enumerate}

% ─────────────────────────────────────────────────────────────
\section*{Capitolo 3: Fondamenti dell'Intelligenza Artificiale nella Patologia}
\addcontentsline{toc}{section}{Capitolo 3}

\begin{enumerate}

    \item \textbf{Risposta: C --- Una rete neurale con molti strati nascosti che apprende
    rappresentazioni gerarchiche dei dati.}\\
    L'apprendimento profondo si riferisce alle reti neurali con molti strati nascosti (da qui
    ``profondo''), che apprendono rappresentazioni gerarchiche dei dati di input. Ogni strato
    apprende caratteristiche di complessità crescente, dalle caratteristiche di basso livello
    (bordi, texture) alle caratteristiche di alto livello (strutture cellulari, pattern tissutali).

    \item \textbf{Risposta: B --- Reti neurali convoluzionali (CNN).}\\
    Le CNN sono l'architettura di apprendimento profondo dominante per l'analisi delle immagini
    nella patologia computazionale. Utilizzano filtri convoluzionali per estrarre caratteristiche
    spaziali locali dalle immagini, rendendole particolarmente adatte all'analisi di immagini
    istologiche.

    \item \textbf{Risposta: C --- Apprendimento debolmente supervisionato con etichette a livello
    di vetrino.}\\
    L'apprendimento debolmente supervisionato (o apprendimento con supervisione debole) utilizza
    etichette a livello di vetrino (ad es., diagnosi del caso) anziché annotazioni pixel per pixel,
    riducendo il carico di annotazione. Questo è particolarmente rilevante nella patologia, dove
    le annotazioni dettagliate richiedono molto tempo da parte di patologi esperti.

    \item \textbf{Risposta: B --- Apprendimento per istanze multiple (MIL).}\\
    Il MIL è un paradigma di apprendimento debolmente supervisionato in cui un vetrino è
    rappresentato come un ``sacchetto'' di tessere (istanze), e l'etichetta del sacchetto (diagnosi
    del vetrino) è nota ma le etichette delle singole tessere non lo sono. Il modello apprende a
    identificare le tessere diagnosticamente rilevanti (istanze positive) all'interno del sacchetto.

    \item \textbf{Risposta: C --- Trasferimento dell'apprendimento da modelli pre-addestrati su
    grandi dataset di immagini naturali.}\\
    Il trasferimento dell'apprendimento (transfer learning) è la tecnica più comune per affrontare
    la scarsità di dati annotati nella patologia computazionale. I modelli pre-addestrati su
    ImageNet o altri grandi dataset vengono adattati (fine-tuned) su dataset di patologia più
    piccoli, sfruttando le rappresentazioni di caratteristiche apprese durante il pre-addestramento.

    \item \textbf{Risposta: B --- La capacità di un modello di generalizzare a nuovi dati non visti
    durante l'addestramento.}\\
    La generalizzazione si riferisce alla capacità di un modello di apprendimento automatico di
    applicare le conoscenze acquisite durante l'addestramento a nuovi dati non visti. Un modello
    con scarsa generalizzazione può avere prestazioni eccellenti sui dati di addestramento ma
    scarse sui dati di test, un fenomeno noto come overfitting.

    \item \textbf{Risposta: C --- Aumentazione dei dati.}\\
    L'aumentazione dei dati (data augmentation) è la tecnica di generare artificialmente nuovi
    esempi di addestramento applicando trasformazioni casuali (rotazione, riflessione, variazione
    del colore, ecc.) alle immagini esistenti. Questo aumenta la dimensione effettiva del dataset
    di addestramento e migliora la robustezza del modello.

    \item \textbf{Risposta: B --- Mappe di attivazione di classe (CAM) e metodi di attribuzione
    del gradiente (Grad-CAM).}\\
    Le CAM e i metodi basati sul gradiente come Grad-CAM sono le tecniche di interpretabilità più
    comunemente utilizzate nella patologia computazionale. Producono mappe di calore che
    evidenziano le regioni dell'immagine che hanno contribuito maggiormente alla decisione del
    modello, consentendo ai patologi di verificare se il modello si concentra sulle caratteristiche
    morfologiche corrette.

    \item \textbf{Risposta: C --- Bias nei dati di addestramento che riflette disparità nella
    raccolta dei dati o nelle pratiche di annotazione.}\\
    Il bias nei modelli di IA per la patologia deriva spesso da bias nei dati di addestramento,
    che possono riflettere disparità nella raccolta dei dati (ad es., sottorappresentazione di
    certi gruppi demografici), nelle pratiche di annotazione (ad es., variabilità inter-osservatore)
    o nelle pratiche di laboratorio (ad es., variabilità della colorazione).

    \item \textbf{Risposta: B --- Validazione esterna su dataset indipendenti di laboratori
    diversi.}\\
    La validazione esterna su dataset indipendenti di laboratori diversi è il gold standard per
    valutare la generalizzabilità di un modello di IA per la patologia. La validazione interna
    (su dati dello stesso laboratorio) può sovrastimare le prestazioni del modello a causa del
    \emph{domain shift} tra laboratori.

\end{enumerate}

% ─────────────────────────────────────────────────────────────
\section*{Capitolo 4: Apprendimento Profondo e Analisi Multimodale}
\addcontentsline{toc}{section}{Capitolo 4}

\begin{enumerate}

    \item \textbf{Risposta: C --- Vision Transformer (ViT).}\\
    I Vision Transformer (ViT) applicano il meccanismo di attenzione dei Transformer, originalmente
    sviluppato per il linguaggio naturale, alle immagini suddividendole in patch. Hanno dimostrato
    prestazioni competitive o superiori alle CNN su molti compiti di analisi delle immagini,
    inclusa la patologia computazionale.

    \item \textbf{Risposta: B --- Segmentazione semantica.}\\
    La segmentazione semantica assegna un'etichetta di classe a ogni pixel dell'immagine,
    consentendo la delineazione precisa di strutture tissutali (ad es., ghiandole, stroma, tumore)
    nelle WSI. Questo è distinto dalla classificazione (etichetta a livello di immagine) e dal
    rilevamento di oggetti (bounding box).

    \item \textbf{Risposta: C --- Integrazione di dati di imaging con dati genomici, clinici o
    di altra natura per migliorare la predizione degli esiti.}\\
    L'analisi multimodale nella patologia computazionale si riferisce all'integrazione di dati di
    imaging (WSI) con altri tipi di dati --- genomici, trascrittomici, clinici, radiologici ---
    per migliorare la predizione degli esiti rispetto a quanto possibile con un singolo tipo di
    dato.

    \item \textbf{Risposta: B --- Reti generative avversariali (GAN).}\\
    Le GAN sono il metodo di apprendimento profondo più comunemente utilizzato per la
    normalizzazione della colorazione e l'aumentazione dei dati nella patologia computazionale.
    Le CycleGAN, in particolare, consentono la traduzione di immagini non appaiata tra domini
    (ad es., da un laboratorio a un altro) senza richiedere coppie di immagini corrispondenti.

    \item \textbf{Risposta: C --- Predizione della sopravvivenza e degli esiti clinici da WSI.}\\
    La predizione della sopravvivenza e degli esiti clinici da WSI è uno dei compiti più
    clinicamente rilevanti nella patologia computazionale. I modelli di apprendimento profondo
    possono estrarre caratteristiche morfologiche prognostiche dalle WSI che non sono facilmente
    quantificabili dall'osservazione umana.

    \item \textbf{Risposta: B --- Modelli fondazionali pre-addestrati su grandi dataset di
    patologia (ad es., UNI, CONCH, Prov-GigaPath).}\\
    I modelli fondazionali per la patologia, pre-addestrati su milioni di tessere WSI mediante
    apprendimento auto-supervisionato, hanno dimostrato prestazioni superiori ai modelli
    pre-addestrati su ImageNet per molti compiti di patologia computazionale. Esempi includono
    UNI, CONCH e Prov-GigaPath.

    \item \textbf{Risposta: C --- Apprendimento auto-supervisionato su grandi dataset non
    etichettati.}\\
    L'apprendimento auto-supervisionato (self-supervised learning) consente ai modelli di
    apprendere rappresentazioni utili da grandi dataset non etichettati, senza richiedere
    annotazioni manuali costose. Questo è particolarmente vantaggioso nella patologia, dove le
    annotazioni richiedono molto tempo da parte di patologi esperti.

    \item \textbf{Risposta: B --- Attenzione multi-testa che modella le relazioni a lungo raggio
    tra le tessere WSI.}\\
    Il meccanismo di attenzione multi-testa dei Transformer consente di modellare le relazioni
    a lungo raggio tra le tessere WSI, catturando il contesto spaziale globale che le CNN con
    campi recettivi locali non possono facilmente catturare. Questo è particolarmente utile per
    compiti che richiedono la comprensione dell'architettura tissutale a scala multipla.

    \item \textbf{Risposta: C --- Fusione di caratteristiche estratte da WSI con dati genomici
    o clinici mediante reti neurali multimodali.}\\
    La fusione multimodale nella patologia computazionale combina caratteristiche estratte da
    diverse modalità di dati (WSI, genomica, clinica) mediante architetture di rete neurale
    progettate per integrare rappresentazioni eterogenee. Questo approccio ha dimostrato di
    migliorare la predizione degli esiti rispetto ai modelli unimodali.

    \item \textbf{Risposta: B --- Modelli linguistici di grandi dimensioni (LLM) e modelli
    visione-linguaggio (VLM) per la generazione di referti e il ragionamento diagnostico.}\\
    I modelli linguistici di grandi dimensioni (LLM) e i modelli visione-linguaggio (VLM) stanno
    emergendo come strumenti promettenti per la generazione automatica di referti patologici, il
    ragionamento diagnostico e l'interazione in linguaggio naturale con i sistemi di patologia
    digitale. Esempi includono GPT-4V, LLaVA e modelli specifici per la patologia come PathChat.

\end{enumerate}

% ─────────────────────────────────────────────────────────────
\section*{Capitolo 5: Etica, Bias e il Futuro della Patologia Digitale}
\addcontentsline{toc}{section}{Capitolo 5}

\begin{enumerate}

    \item \textbf{Risposta: C --- Bias algoritmico che porta a prestazioni diseguali tra gruppi
    demografici.}\\
    Il bias algoritmico nella patologia digitale si riferisce alla tendenza dei modelli di IA a
    avere prestazioni diseguali tra gruppi demografici (ad es., per razza, sesso, età), spesso
    a causa di una rappresentazione non equilibrata di questi gruppi nei dati di addestramento.
    Questo può portare a disparità nelle cure sanitarie se i modelli vengono implementati
    clinicamente senza un'adeguata valutazione dell'equità.

    \item \textbf{Risposta: B --- Interpretabilità e spiegabilità dei modelli di IA.}\\
    L'interpretabilità e la spiegabilità dei modelli di IA sono requisiti fondamentali per
    l'implementazione clinica sicura. I patologi devono essere in grado di comprendere e
    verificare le basi delle decisioni del modello per poterle integrare nella loro pratica
    diagnostica in modo responsabile.

    \item \textbf{Risposta: C --- Consenso informato, privacy dei dati e proprietà dei dati.}\\
    Le questioni etiche relative all'uso dei dati dei pazienti nella patologia digitale includono
    il consenso informato (i pazienti devono essere informati su come vengono utilizzati i loro
    dati), la privacy dei dati (protezione delle informazioni identificabili) e la proprietà dei
    dati (chi possiede i dati dei pazienti e chi può trarne beneficio economico).

    \item \textbf{Risposta: B --- Validazione prospettica in studi clinici controllati.}\\
    La validazione prospettica in studi clinici controllati è il gold standard per dimostrare
    l'utilità clinica di un sistema di IA per la patologia. La validazione retrospettiva su
    dataset storici, sebbene necessaria, non è sufficiente per dimostrare che il sistema migliora
    gli esiti clinici nella pratica reale.

    \item \textbf{Risposta: C --- Supervisione umana e responsabilità professionale del patologo.}\\
    Nonostante i progressi dell'IA, la supervisione umana e la responsabilità professionale del
    patologo rimangono fondamentali. I sistemi di IA sono strumenti di supporto decisionale, non
    sostituti del giudizio clinico del patologo. La responsabilità legale e professionale della
    diagnosi rimane con il patologo.

    \item \textbf{Risposta: B --- Regolamento sui Dispositivi Medici (MDR) e Regolamento sui
    Dispositivi Medico-Diagnostici In Vitro (IVDR).}\\
    Nell'Unione Europea, i sistemi di IA per la patologia digitale utilizzati per la diagnosi
    primaria sono regolamentati come dispositivi medici ai sensi del MDR (UE 2017/745) o come
    dispositivi medico-diagnostici in vitro ai sensi dell'IVDR (UE 2017/746), a seconda della
    loro funzione specifica.

    \item \textbf{Risposta: C --- Federated learning per l'addestramento di modelli su dati
    distribuiti senza condivisione dei dati grezzi.}\\
    Il federated learning consente l'addestramento di modelli di IA su dati distribuiti in più
    istituzioni senza condividere i dati grezzi dei pazienti, affrontando le preoccupazioni sulla
    privacy e sulla sovranità dei dati. I modelli vengono addestrati localmente in ogni istituzione
    e solo i parametri del modello (non i dati) vengono condivisi con un server centrale.

    \item \textbf{Risposta: B --- Integrazione dell'IA come strumento di supporto decisionale
    nella pratica diagnostica di routine.}\\
    Il ruolo più probabile dell'IA nella patologia nel prossimo futuro è come strumento di
    supporto decisionale che assiste il patologo nella diagnosi di routine, piuttosto che come
    sostituto del patologo. L'IA può aumentare l'efficienza, la riproducibilità e la sensibilità
    diagnostica, ma la responsabilità finale della diagnosi rimane con il patologo.

    \item \textbf{Risposta: C --- Disparità nell'accesso alle tecnologie di patologia digitale
    tra paesi ad alto e basso reddito.}\\
    Le disparità nell'accesso alle tecnologie di patologia digitale tra paesi ad alto e basso
    reddito sono una preoccupazione etica significativa. Sebbene la patologia digitale abbia il
    potenziale di migliorare l'accesso alla diagnosi esperta attraverso la telemedicina, i costi
    dell'infrastruttura tecnologica possono esacerbare le disparità esistenti se non vengono
    affrontate deliberatamente.

    \item \textbf{Risposta: B --- Collaborazione interdisciplinare tra patologi, informatici,
    bioinformatici e clinici.}\\
    Lo sviluppo e l'implementazione responsabile dell'IA nella patologia richiedono una
    collaborazione interdisciplinare tra patologi (che forniscono competenza di dominio e
    annotazioni), informatici e bioinformatici (che sviluppano e validano i modelli) e clinici
    (che definiscono le esigenze cliniche e valutano l'utilità). Nessuna singola disciplina può
    affrontare da sola le sfide tecniche, cliniche ed etiche della patologia digitale basata
    sull'IA.

\end{enumerate}


\end{document}
